% Transformations chimiques forcées : l'électrolyse — Chimie 1ère D — Cours signé OnBuch+
% Programme officiel MINESEC, Première C, D, E, TI. Compiler avec : tectonic N09-transformations-chimiques-forcees-electrolyse.tex
\newcommand{\DOCMATIERE}{Chimie}
\newcommand{\DOCNIVEAU}{1ère D}
\newcommand{\DOCMODULE}{Module 2 : Oxydoréduction}
\newcommand{\DOCLECON}{09}
\newcommand{\DOCTITRE}{Transformations chimiques forcées : l'électrolyse}
% OnBuch+ — préambule « cours signés » ENRICHI (compilé avec tectonic / XeLaTeX)
% Système visuel « Pop » d'OnBuch+ : fond crème (jamais blanc), bordures encre
% épaisses, ombres « dures » décalées, titres Archivo Black en capitales.
% Version enrichie pour les cours de chimie : chimie (mhchem, chemfig),
% graphes (pgfplots), boîtes pédagogiques supplémentaires (définition,
% propriété, méthode, attention, le savais-tu, exercice résolu, exercice,
% TP, bilan), page de garde refaite et signature OnBuch+ ancrée en bas.
%
% Macros attendues dans main.tex AVANT \input{preamble} :
%   \DOCMATIERE  \DOCNIVEAU  \DOCMODULE  \DOCLECON (numéro)  \DOCTITRE
\documentclass[11pt]{article}

\usepackage{fontspec}
\usepackage[french]{babel}
\usepackage[a4paper,margin=2cm,top=3.4cm,bottom=2.8cm,headheight=1.4cm,footskip=1.3cm]{geometry}
\usepackage{amsmath,amssymb,amsfonts}
\usepackage[table]{xcolor} % [table] : fournit aussi \rowcolors (alternance de lignes)
\usepackage{pagecolor}
\usepackage{tikz}
\usetikzlibrary{arrows.meta,positioning,calc,shapes.geometric,shapes.misc,shapes.arrows,decorations.pathmorphing,decorations.markings,patterns,shadows,mindmap,trees,fit,backgrounds,matrix,intersections}
\usepackage{pgfplots}
\pgfplotsset{compat=1.18}
\usepackage[version=4]{mhchem}
\usepackage{chemfig}
\usepackage{enumitem}
\usepackage{fancyhdr}
% [most] charge toutes les bibliothèques tcolorbox (listings, minted, raster,
% documentation...) et gonfle inutilement la mémoire TeX sur un document long
% (~300 boîtes) : on ne charge que ce qui est réellement utilisé.
\usepackage[skins,breakable]{tcolorbox}
\usepackage{graphicx}
\usepackage{colortbl}
\usepackage{booktabs}
\usepackage{tabularx}
\usepackage{array}
\usepackage{multicol}
\usepackage{siunitx}
% parse-numbers=false : accepte aussi \SI{2,5 \times 10^{-3}}{...}
\sisetup{locale=FR,output-decimal-marker={,},group-minimum-digits=4,parse-numbers=false}
\DeclareSIUnit\ohm{\ensuremath{\symup{\Omega}}} % U+2126 absent de la police
\usepackage{qrcode}
% \degree : commande fréquemment utilisée par les modèles pour un angle ou une
% température en dehors de \SI{}{} (non fournie par les paquets chargés).
\providecommand{\degree}{^{\circ}}
\usepackage{lastpage}
\usepackage{hyperref}
\hypersetup{hidelinks}

% ============================================================
% Palette « Pop » OnBuch+ — mêmes hex que la classe P (plus_ui.dart)
% ============================================================
\definecolor{poporange}{HTML}{F59321}
\definecolor{popdark}{HTML}{E07A0C}
\definecolor{popcream}{HTML}{FFF6E8}
\definecolor{popink}{HTML}{1C1714}
\definecolor{poppurple}{HTML}{7A5AE0}
\definecolor{popgreen}{HTML}{2E9E6A}
\definecolor{poppink}{HTML}{E0457B}
\definecolor{popblue}{HTML}{2F7DE1}
\definecolor{popgold}{HTML}{C98A12}
\definecolor{popmuted}{HTML}{8A7F76}
\definecolor{popline}{HTML}{EADFD2}
\definecolor{popgray}{HTML}{8A7F76}
\colorlet{popgrey}{popgray}
% Alias tolérants (le modèle nomme parfois les couleurs différemment)
\colorlet{popwhite}{white}
\colorlet{popblack}{popink}
\colorlet{popred}{poppink}
\colorlet{popyellow}{popgold}
% Teintes claires pour fonds de tableaux / zones
\colorlet{popblueL}{popblue!10!white}
\colorlet{poporangeL}{poporange!14!white}
\colorlet{popgreenL}{popgreen!12!white}
\colorlet{poppurpleL}{poppurple!12!white}
\colorlet{poppinkL}{poppink!12!white}
\colorlet{popgoldL}{popgold!14!white}

\pagecolor{popcream}

% ============================================================
% Typographie
% ============================================================
\defaultfontfeatures{Ligatures=TeX}
\setmainfont{PlusJakartaSans-Medium.ttf}[Path=../../../fonts/,BoldFont=PlusJakartaSans-Bold.ttf]
% Maths en sans-serif (Fira Math) avec chiffres et lettres droites de Plus
% Jakarta, pour que les formules se fondent dans le texte.
\usepackage[math-style=ISO,bold-style=ISO]{unicode-math}
\setmathfont{FiraMath-Regular.otf}[Path=../../../fonts/]
\setmathfont{PlusJakartaSans-Medium.ttf}[Path=../../../fonts/,range={up/num,up/{Latin,latin},\mathup}]
\usepackage{icomma} % virgule décimale sans espace en mode math
% Caractères absents de la police de texte (sinon carré vide) : rendus par
% leur équivalent mathématique, en texte comme en formule.
\usepackage{newunicodechar}
\newunicodechar{α}{\ensuremath{\alpha}}
\newunicodechar{Α}{\ensuremath{\alpha}}
\newunicodechar{β}{\ensuremath{\beta}}
\newunicodechar{δ}{\ensuremath{\delta}}
\newunicodechar{✓}{\popcheck{popgreen}}
\newfontfamily\popdisplay{ArchivoBlack-Regular.ttf}[Path=../../../fonts/]
\newfontfamily\poplogo{SpaceGrotesk-Bold.ttf}[Path=../../../fonts/]
\newfontfamily\popsemi{PlusJakartaSans-SemiBold.ttf}[Path=../../../fonts/]
\newfontfamily\popxbold{PlusJakartaSans-ExtraBold.ttf}[Path=../../../fonts/]
\newfontfamily\popxboldls{PlusJakartaSans-ExtraBold.ttf}[Path=../../../fonts/,LetterSpace=3.0]

\setlist[enumerate]{leftmargin=1.7em,itemsep=5pt,topsep=4pt}
\setlist[itemize]{leftmargin=1.7em,itemsep=4pt,topsep=4pt,label={\color{poporange}\textbullet}}
\setlength{\parindent}{0pt}
\setlength{\parskip}{5pt}
\renewcommand{\arraystretch}{1.25}

% chemfig : liaisons un peu plus épaisses, encre Pop
\setchemfig{atom sep=2em,bond style={line width=0.9pt,popink},cram width=3pt}

% pgfplots : style Pop par défaut
\pgfplotsset{
  every axis/.append style={axis line style={popink,line width=1pt},tick style={popink},
    grid style={popline,line width=0.5pt},label style={font=\small},tick label style={font=\footnotesize}},
  popcurve/.style={poporange,line width=1.8pt,no markers,smooth},
}

% ============================================================
% Pill / squiggle
% ============================================================
\newcommand{\poppill}[3]{%
  \tikz[baseline=-0.55ex]\node[draw=popink,line width=1.2pt,rounded corners=2.6pt,%
    fill=#2,inner xsep=6.5pt,inner ysep=3pt]{%
    \popxboldls\fontsize{7}{7}\selectfont\color{#3}\MakeUppercase{#1}};%
}
\newcommand{\popsquiggle}[1]{%
  \tikz[baseline=-0.4ex]{
    \draw[#1,line width=2.6pt,line cap=round]
      plot [smooth] coordinates {(0,0.09) (0.34,0.01) (0.68,0.21) (1.02,0.01) (1.36,0.10)};
  }%
}
% Mot-clé surligné (marqueur orange sous le texte)
\newcommand{\cle}[1]{{\popxbold\color{popdark}#1}}

% ============================================================
% En-tête / pied
% ============================================================
\pagestyle{fancy}
\fancyhf{}
\renewcommand{\headrulewidth}{0pt}
\renewcommand{\footrulewidth}{0pt}
\lhead{\raisebox{-0.15ex}{\poplogo\fontsize{13}{13}\selectfont\color{popink}OnBuch\color{poporange}+}}
\rhead{\poppill{\DOCMATIERE\ \textbullet\ \DOCNIVEAU\ \textbullet\ Leçon \DOCLECON}{white}{popink}}
\lfoot{\popxbold\fontsize{7.6}{7.6}\selectfont\color{popmuted}\MakeUppercase{Cours signé OnBuch+}\ \textbullet\ {\popsemi\DOCTITRE}}
\rfoot{\tikz[baseline=-0.6ex]\node[rounded corners=8.5pt,fill=popink,minimum height=17pt,minimum width=17pt,inner xsep=5pt,inner sep=0]{\popxbold\fontsize{8}{8}\selectfont\color{popcream}\thepage\,/\,\pageref*{LastPage}};}

% ============================================================
% Titres
% ============================================================
\newcommand{\chaptitre}[2]{%
  \begin{tcolorbox}[enhanced,colback=poporange,colframe=popink,boxrule=2.6pt,arc=11pt,
      drop shadow=popink,left=15pt,right=15pt,top=12pt,bottom=11pt,before skip=3pt,after skip=18pt]
    \poppill{#2}{popink}{popcream}\par\vspace{9pt}
    {\raggedright\hyphenpenalty=10000\exhyphenpenalty=10000\popdisplay\fontsize{22}{24}\selectfont\color{popcream}\MakeUppercase{#1}\par}
  \end{tcolorbox}
}
% Section numérotée : pastille orange avec le numéro + titre Archivo
\newcounter{popsec}
\newcommand{\coursec}[1]{%
  \stepcounter{popsec}\setcounter{popsub}{0}%
  \par\vspace{16pt}\noindent
  \begin{minipage}{\linewidth}
  \tikz[baseline=-0.7ex]\node[draw=popink,line width=1.6pt,fill=poporange,rounded corners=4pt,minimum size=22pt,inner sep=0]{\popdisplay\fontsize{12}{12}\selectfont\color{popcream}\thepopsec};%
  \hspace{8pt}{\popdisplay\fontsize{15}{17}\selectfont\color{popink}\MakeUppercase{#1}}\par
  \vspace{4pt}\noindent{\color{poporange}\rule{36pt}{3.6pt}}
  \end{minipage}\par\nopagebreak\vspace{8pt}
}
\newcounter{popsub}
\newcommand{\courssub}[1]{%
  \stepcounter{popsub}%
  \par\vspace{9pt}\noindent{\popxbold\fontsize{11.5}{13}\selectfont\color{popink}{\color{poporange}\thepopsec.\thepopsub}\hspace{6pt}#1}\par\nopagebreak\vspace{4pt}
}

% ============================================================
% Boîtes pédagogiques — même recette : fond blanc, bordure épaisse colorée,
% ombre dure encre, badge plein. Titre optionnel [#1].
% ============================================================
\tcbset{popbase/.style={breakable,enhanced,colback=white,boxrule=2.3pt,arc=9pt,
  shadow={3pt}{-3pt}{0pt}{popink},left=14pt,right=14pt,top=11pt,bottom=11pt,before skip=11pt,after skip=15pt,
  fontupper=\popsemi\fontsize{10.6}{14.5}\selectfont\color{popink},
  fontlower=\popsemi\fontsize{10.6}{14.5}\selectfont\color{popink}}}
% Coche dessinée (le glyphe ✓ n'existe pas dans les polices du document)
\newcommand{\popcheck}[1]{\tikz[baseline=-0.5ex]\draw[#1,line width=1.6pt,line cap=round,line join=round](-0.13,0.0)--(-0.04,-0.09)--(0.13,0.1);}
\newcommand{\popboxhead}[3]{\poppill{#1}{#2}{white}\ifx\relax#3\relax\else\hspace{6pt}{\popxbold\fontsize{10.6}{12}\selectfont\color{popink}#3}\fi\par\vspace{7pt}}

\newtcolorbox{aretenir}[1][]{popbase,colframe=popgreen,before upper={\popboxhead{À retenir}{popgreen}{#1}}}
\newtcolorbox{exemplebox}[1][]{popbase,colframe=poporange,before upper={\popboxhead{Exemple}{poporange}{#1}}}
\newtcolorbox{definition}[1][]{popbase,colframe=popblue,before upper={\popboxhead{Définition}{popblue}{#1}}}
\newtcolorbox{propriete}[1][]{popbase,colframe=poppurple,before upper={\popboxhead{Propriété}{poppurple}{#1}}}
\newtcolorbox{methode}[1][]{popbase,colframe=popgold,before upper={\popboxhead{Méthode}{popgold}{#1}}}
\newtcolorbox{attention}[1][]{popbase,colframe=poppink,before upper={\popboxhead{Attention}{poppink}{#1}}}
\newtcolorbox{experience}[1][]{popbase,colframe=popblue,colback=popblueL,before upper={\popboxhead{Expérience}{popblue}{#1}}}
% « Le savais-tu ? » : carte violette pleine (contexte, culture, Cameroun)
\newtcolorbox{savaistu}[1][]{popbase,colback=poppurple,colframe=popink,
  fontupper=\popsemi\fontsize{10.4}{14.2}\selectfont\color{white},
  before upper={\poppill{Le savais-tu\,?}{white}{poppurple}\ifx\relax#1\relax\else\hspace{6pt}{\popxbold\color{white}#1}\fi\par\vspace{7pt}}}
% Exercice résolu : énoncé en haut, \tcblower puis la solution sur fond crème
\newcounter{popexo}\newcounter{popexr}
\newtcolorbox{exoresolu}[1][]{popbase,colframe=poporange,colbacklower=popcream,
  segmentation style={popink,line width=1.2pt,dashed},
  before upper={\stepcounter{popexr}\popboxhead{Exercice résolu \thepopexr}{poporange}{#1}},
  before lower={\poppill{Solution}{popink}{popcream}\par\vspace{6pt}}}
% Exercice (non résolu en place) — la correction est dans la partie Corrigés
\newtcolorbox{exercice}[1][]{popbase,colframe=popink,
  before upper={\stepcounter{popexo}\popboxhead{Exercice \thepopexo}{popink}{#1}}}
% Corrigé
\newtcolorbox{corrige}[1][]{popbase,colframe=popgreen,colback=popgreenL,
  before upper={\popboxhead{Corrigé}{popgreen}{#1}}}
% Objectifs / prérequis / bilan
\newtcolorbox{objectifs}{popbase,colframe=popink,colback=poporangeL,
  before upper={\popboxhead{Objectifs de la leçon}{popink}{}}}
\newtcolorbox{prerequis}{popbase,colframe=popink,colback=popblueL,
  before upper={\popboxhead{Prérequis}{popblue}{}}}
\newtcolorbox{bilan}{popbase,colframe=popink,colback=white,boxrule=2.8pt,
  before upper={\popboxhead{Fiche bilan}{poporange}{L'essentiel à connaître pour l'examen}}}
% Figure encadrée avec légende
\newtcolorbox{popfigure}[1][]{enhanced,breakable=false,colback=white,colframe=popline,boxrule=1.4pt,arc=8pt,
  left=10pt,right=10pt,top=10pt,bottom=8pt,before skip=10pt,after skip=12pt,halign=center,
  lower separated=false,halign lower=center,
  fontlower=\popsemi\fontsize{9}{11.5}\selectfont\color{popmuted},
  #1}
\newcommand{\legende}[1]{\par\vspace{6pt}{\popsemi\fontsize{9}{11.5}\selectfont\color{popmuted}#1}}

% ============================================================
% Signature OnBuch+
% ============================================================
\newcommand{\poplogoinline}[1]{{\poplogo\fontsize{#1}{#1}\selectfont\color{popink}OnBuch\color{poporange}+}}
\newcommand{\signatureonbuch}{%
  \begin{tcolorbox}[enhanced,colback=popink,colframe=popink,boxrule=2.2pt,arc=10pt,
      drop shadow=poporange,left=14pt,right=14pt,top=10pt,bottom=10pt,before skip=0pt,after skip=0pt]
    \tikz[baseline=-0.7ex]\node[circle,fill=popgreen,draw=popcream,line width=1.3pt,minimum size=22pt,inner sep=0]{\popcheck{white}};%
    \hspace{9pt}\begin{minipage}[c]{0.62\linewidth}
      {\popxbold\fontsize{12}{13}\selectfont\color{popcream}Signé\ \textbullet\ {\poplogo OnBuch\color{poporange}+}}\par\vspace{2pt}
      {\raggedright\popsemi\fontsize{8}{10.5}\selectfont\color{popline}\MakeUppercase{Équipe pédagogique OnBuch+}\\\MakeUppercase{Conforme au programme officiel MINESEC}\par}
    \end{minipage}\hfill
    {\poplogo\fontsize{22}{22}\selectfont\color{popcream}OnBuch\color{poporange}+}
  \end{tcolorbox}%
}

% ============================================================
% Page de garde
% #1 titre · #2 matière · #3 niveau · #4 module · #5 n° leçon · #6 durée ·
% #7 description / objectifs (texte ou itemize)
% ============================================================
\newcommand{\pagegarde}[7]{%
  \thispagestyle{empty}
  \newgeometry{a4paper,margin=1.7cm,top=1.6cm,bottom=1.4cm}%
  \begin{tikzpicture}[remember picture,overlay]
    \fill[poporange!18] ([xshift=-3.2cm,yshift=-3.6cm]current page.north east) circle (4.6cm);
    \fill[poppurple!14] ([xshift=2.4cm,yshift=7.6cm]current page.south west) circle (3.8cm);
  \end{tikzpicture}%
  {\poplogo\fontsize{30}{30}\selectfont\color{popink}OnBuch\color{poporange}+}\hfill
  \raisebox{0.6ex}{\poppill{Cours signé \textbullet\ Premium}{popink}{popcream}}\par
  \vspace{1pt}\popsquiggle{poppurple}\par
  \vspace{8pt}
  \begin{tcolorbox}[enhanced,colback=poporange,colframe=popink,boxrule=2.8pt,arc=13pt,
      drop shadow=popink,left=18pt,right=18pt,top=12pt,bottom=13pt,after skip=10pt]
    \poppill{#2 \textbullet\ #3}{popink}{popcream}\hspace{5pt}\poppill{#4}{white}{popink}\par\vspace{8pt}
    {\popdisplay\fontsize{14}{14}\selectfont\color{popink}LEÇON #5}\par\vspace{4pt}
    {\raggedright\hyphenpenalty=10000\exhyphenpenalty=10000\popdisplay\fontsize{25}{27}\selectfont\color{popcream}\MakeUppercase{#1}\par}
    \vspace{8pt}
    {\popxbold\fontsize{9.5}{9.5}\selectfont\color{popink}\MakeUppercase{Durée conseillée : #6}}
  \end{tcolorbox}
  \begin{tcolorbox}[enhanced,colback=white,colframe=popink,boxrule=2.2pt,arc=10pt,
      drop shadow=popink,left=16pt,right=16pt,top=10pt,bottom=10pt,after skip=8pt]
    \poppill{Dans cette leçon}{poppurple}{white}\par\vspace{5pt}
    \popsemi\fontsize{9.3}{12.6}\selectfont\color{popink}#7
  \end{tcolorbox}
  \noindent\begin{minipage}[t]{0.487\linewidth}
    \begin{tcolorbox}[enhanced,colback=popgreen,colframe=popink,boxrule=2.2pt,arc=10pt,
        drop shadow=popink,left=11pt,right=11pt,top=10pt,bottom=10pt,height=3.1cm,valign=center]
      \tcbox[colback=white,colframe=popink,boxrule=1.3pt,arc=5pt,boxsep=0pt,nobeforeafter,
        left=3pt,right=3pt,top=3pt,bottom=3pt,valign=center]{\qrcode[height=1.9cm]{https://play.google.com/store/apps/details?id=com.luvvix.onbuch&hl=fr}}%
      \hspace{8pt}\begin{minipage}[c]{0.5\linewidth}
        {\popdisplay\fontsize{11}{12.5}\selectfont\color{white}\MakeUppercase{Télécharge OnBuch}}\par\vspace{4pt}
        {\raggedright\popsemi\fontsize{8.4}{11}\selectfont\color{white}Gratuit \textbullet\ Google Play\\Tuteur IA, annales, concours, résultats\par}
      \end{minipage}
    \end{tcolorbox}
  \end{minipage}\hfill
  \begin{minipage}[t]{0.487\linewidth}
    \begin{tcolorbox}[enhanced,colback=poppurple,colframe=popink,boxrule=2.2pt,arc=10pt,
        drop shadow=popink,left=11pt,right=11pt,top=10pt,bottom=10pt,height=3.1cm,valign=center]
      \tikz[baseline=-0.7ex]\node[circle,fill=white,draw=popink,line width=1.3pt,minimum size=1.6cm,inner sep=0]{\popdisplay\fontsize{24}{24}\selectfont\color{poppurple}+};%
      \hspace{8pt}\begin{minipage}[c]{0.6\linewidth}
        {\popdisplay\fontsize{11}{12.5}\selectfont\color{white}\MakeUppercase{Passe à OnBuch+}}\par\vspace{4pt}
        {\raggedright\popsemi\fontsize{8.4}{11}\selectfont\color{white}Tous les cours signés, Tuteur IA illimité, fiches PDF — onglet OnBuch+ de l'app\par}
      \end{minipage}
    \end{tcolorbox}
  \end{minipage}
  \vspace{8pt}
  \popcontenu
  \vfill
  \signatureonbuch
  \restoregeometry
  \newpage
}

% Rangée « Ce cours contient » (page de garde)
\newcommand{\poptile}[3]{% #1 couleur #2 gros texte #3 légende
  \begin{tcolorbox}[enhanced,colback=#1,colframe=popink,boxrule=2pt,arc=9pt,shadow={2.5pt}{-2.5pt}{0pt}{popink},
      left=6pt,right=6pt,top=9pt,bottom=9pt,halign=center,valign=center,height=1.75cm,width=\linewidth,nobeforeafter]
    {\popdisplay\fontsize{12.5}{13}\selectfont\color{white}\MakeUppercase{#2}}\par\vspace{3pt}
    {\centering\popsemi\fontsize{7.8}{9.5}\selectfont\color{white}#3\par}
  \end{tcolorbox}}
\newcommand{\popcontenu}{%
  \par\noindent{\popxboldls\fontsize{7.5}{7.5}\selectfont\color{popmuted}CE COURS SIGNÉ CONTIENT}\par\vspace{7pt}
  \noindent\begin{minipage}[t]{0.235\linewidth}\poptile{popblue}{Cours}{complet, illustré, pas à pas}\end{minipage}\hfill
  \begin{minipage}[t]{0.235\linewidth}\poptile{poporange}{Méthodes}{et exercices résolus}\end{minipage}\hfill
  \begin{minipage}[t]{0.235\linewidth}\poptile{poppink}{Exercices}{type examen + corrigés}\end{minipage}\hfill
  \begin{minipage}[t]{0.235\linewidth}\poptile{popgreen}{Fiche bilan}{l'essentiel à retenir}\end{minipage}\par}

% Sommaire
\newtcolorbox{sommaire}{enhanced,colback=white,colframe=popink,boxrule=2.2pt,arc=10pt,
  drop shadow=popink,left=18pt,right=18pt,top=13pt,bottom=13pt,before skip=6pt,after skip=20pt,
  before upper={\poppill{Sommaire}{poppurple}{white}\par\vspace{9pt}\popsemi\fontsize{10.6}{14.5}\selectfont\color{popink}}}

% Signature de fin de cours — poussée tout en bas de la dernière page
\newcommand{\findecours}{%
  \par\vfill\nopagebreak
  \begin{center}\popsquiggle{poporange}\end{center}\vspace{4pt}
  \signatureonbuch
}
% Glyphes absents de Fira Math : redéfinis à partir de glyphes disponibles
\AtBeginDocument{%
  \def\setminus{\mathbin{\backslash}}\let\smallsetminus\setminus
  \def\vdots{\mathord{\vbox{\baselineskip3.2pt\lineskiplimit0pt\kern4pt\hbox{.}\hbox{.}\hbox{.}}}}%
  \def\ddots{\mathinner{\mkern1mu\raise6.5pt\hbox{.}\mkern2mu\raise3.6pt\hbox{.}\mkern2mu\raise0.7pt\hbox{.}\mkern1mu}}%
  \def\triangle{\mathord{\scalebox{0.9}{$\symup{\Delta}$}}}%
  \def\nabla{\mathord{\raisebox{\depth}{\scalebox{1}[-1]{$\symup{\Delta}$}}}}%
  \def\checkmark{\popcheck{popdark}}%
  \def\star{\mathbin{\ast}}\def\bigstar{\mathord{\ast}}%
  \def\diamond{\mathbin{\scalebox{0.6}{$\Diamond$}}}%
  \def\bigcup{\mathop{\mathchoice{\raisebox{-0.3em}{\scalebox{1.7}{$\cup$}}}{\scalebox{1.25}{$\cup$}}{\cup}{\cup}}\displaylimits}%
  \def\bigcap{\mathop{\mathchoice{\raisebox{-0.3em}{\scalebox{1.7}{$\cap$}}}{\scalebox{1.25}{$\cap$}}{\cap}{\cap}}\displaylimits}%
  \def\lesseqqgtr{\mathrel{\lessgtr}}\def\bigoplus{\mathop{\oplus}\displaylimits}%
}
\providecommand{\ssi}{si et seulement si} % abréviation fréquente

%% FIGFIX-BEGIN v3 — correctifs globaux des figures (docs/onbuch-generation/tools/figfix)
\usepackage{adjustbox}\usepackage{etoolbox}
% toute figure TikZ/pgfplots plus large que la ligne est réduite à la largeur disponible
\BeforeBeginEnvironment{tikzpicture}{\begin{adjustbox}{max width=\linewidth}}
\AfterEndEnvironment{tikzpicture}{\end{adjustbox}}
% légendes pgfplots lisibles par-dessus les courbes
\pgfplotsset{every axis/.append style={legend style={fill=white,fill opacity=0.92,text opacity=1,draw=black!15,inner sep=3pt}}}
% étiquettes d'arêtes (syntaxe edge["..."]) avec fond blanc
\tikzset{every edge quotes/.append style={fill=white,inner sep=1.2pt}}
% bulles de cartes mentales : texte à la ligne dans la bulle, bulle un peu plus grande
\tikzset{every concept/.append style={text width=2.0cm,align=center,minimum size=2.3cm,inner sep=2pt}}
%% FIGFIX-END
\begin{document}

\pagegarde{Transformations chimiques forcées : l'électrolyse}{Chimie}{1ère D}{Module 2 : Oxydoréduction}{09}{2 h}{Cette leçon étudie les transformations chimiques forcées réalisées grâce à l'énergie électrique : l'électrolyse. Elle présente les réactions aux électrodes, les règles de prévision, le phénomène de surtension, les aspects quantitatifs, ainsi que la corrosion des métaux et les moyens de s'en protéger.
\vspace{4pt}
\begin{multicols}{2}\small
\begin{itemize}
\item À la fin de cette leçon, je sais définir une électrolyse et la distinguer d'une transformation spontanée (pile)
\item identifier l'anode et la cathode et indiquer le type de réaction à chaque électrode (oxydation anodique, réduction cathodique)
\item écrire les demi-équations électroniques aux électrodes et l'équation bilan d'une électrolyse
\item appliquer les règles de prévision pour déterminer les produits d'une électrolyse en solution aqueuse
\item expliquer le phénomène de surtension et son rôle dans le choix des réactions aux électrodes
\item réaliser et interpréter les électrolyses de HCl, NaCl, NaOH et CuSO4 en solution aqueuse
\end{itemize}\end{multicols}}

\begin{sommaire}
\begin{multicols}{2}
\begin{enumerate}
\item Généralités sur l'électrolyse
\item Réactions aux électrodes
\item Prévision des réactions d'électrolyse
\item Le phénomène de surtension
\item Aspects quantitatifs de l'électrolyse
\item Corrosion et protection des métaux
\item Travaux pratiques
\item Méthodes et exercices résolus
\item Exercices
\item Corrigés des exercices
\item Fiche bilan
\end{enumerate}
\end{multicols}
\end{sommaire}

\begin{objectifs}
\begin{itemize}
\item À la fin de cette leçon, je sais définir une électrolyse et la distinguer d'une transformation spontanée (pile)
\item identifier l'anode et la cathode et indiquer le type de réaction à chaque électrode (oxydation anodique, réduction cathodique)
\item écrire les demi-équations électroniques aux électrodes et l'équation bilan d'une électrolyse
\item appliquer les règles de prévision pour déterminer les produits d'une électrolyse en solution aqueuse
\item expliquer le phénomène de surtension et son rôle dans le choix des réactions aux électrodes
\item réaliser et interpréter les électrolyses de HCl, NaCl, NaOH et CuSO4 en solution aqueuse
\item exploiter les aspects quantitatifs de l'électrolyse (quantité d'électricité, masses et volumes de produits)
\item expliquer le phénomène de corrosion des métaux et citer des méthodes de protection
\end{itemize}
\end{objectifs}

\begin{prerequis}
\begin{itemize}
\item Couples oxydant/réducteur et demi-équations électroniques (leçon sur l'oxydoréduction)
\item Réactions d'oxydoréduction spontanées et fonctionnement d'une pile
\item Classification électrochimique des métaux et potentiels standards
\item Quantité de matière, concentration molaire et relation Q = I × t (électricité, Seconde)
\item Ions en solution aqueuse et électroneutralité (Première)
\end{itemize}
\end{prerequis}

\newpage

\coursec{Généralités sur l'électrolyse}

\noindent
Imagine un instant : tu es au marché de \emph{Nkololoun} à Douala ou au \emph{Marché Central} de Yaoundé. Sur les étals, des bidons d'eau de Javel attendent les acheteurs. Cette eau de Javel, si précieuse pour désinfecter l'eau de boisson, blanchir le linge ou nettoyer les sols, est produite industriellement à partir d'une simple solution de sel de cuisine (chlorure de sodium) dans laquelle on fait passer un courant électrique. Comme par magie, le sel et l'eau se transforment en hypochlorite de sodium (le principe actif de l'eau de Javel), en soude et en dihydrogène.

\medskip\noindent
Changeons de décor. Direction \textbf{Édéa}, au bord du fleuve Sanaga. Là, l'usine \textbf{ALUCAM} (Aluminium du Cameroun) fonctionne jour et nuit. D'immenses cuves remplies d'alumine dissoute dans un bain fondu reçoivent un courant électrique d'une intensité phénoménale (de l'ordre de \SI{100}{\kilo\ampere} et plus par cuve). Résultat : de l'aluminium métal liquide s'accumule au fond des cuves, prêt à être coulé en lingots. Pourtant, si tu laisses de l'alumine ($\ce{Al2O3}$) dans un coin, elle ne se transformera \emph{jamais} en aluminium toute seule.

\medskip\noindent
À l'opposé, regarde les tôles ondulées des cases, les carcasses de voitures abandonnées au bord de la route ou les pylônes électriques : la rouille les grignote inexorablement, \emph{sans que personne n'ait branché le moindre fil électrique}. Cette transformation (fer $\rightarrow$ oxyde de fer hydraté) se fait \emph{toute seule}, spontanément.

\medskip\noindent
\textbf{Problématique :} Pourquoi certaines réactions (fabrication de l'eau de Javel, production d'aluminium) refusent-elles de se produire naturellement et ont-elles besoin qu'on les « force » avec de l'électricité, alors que d'autres (rouille, fonctionnement d'une pile) se font toutes seules ? Comment un courant électrique impose-t-il son sens à une transformation chimique ? C'est tout l'objet de l'\textbf{électrolyse}, que nous allons décortiquer dans cette première section.

\courssub{De la pile à l'électrolyse : spontanéité contre contrainte}

\noindent
Tu as déjà découvert la \textbf{pile} : un dispositif où une transformation chimique \textbf{spontanée} (une réaction d'oxydoréduction qui se produit d'elle-même) fournit de l'énergie électrique. Les électrons circulent \emph{d'eux-mêmes} de l'anode (pôle $-$) vers la cathode (pôle $+$) dans le circuit extérieur. La pile est un \textbf{générateur} d'énergie électrique.

\medskip\noindent
L'\textbf{électrolyse} est l'exact inverse : on utilise de l'énergie électrique (fournie par un générateur extérieur : pile, batterie, alimentation) pour \textbf{forcer} une transformation chimique \textbf{non spontanée} à se produire. L'électrolyseur est un \textbf{récepteur} d'énergie électrique.

\begin{definition}[Électrolyse]
L'\textbf{électrolyse} est une transformation chimique d'oxydoréduction \textbf{forcée}, provoquée par le passage d'un courant électrique imposé par un générateur extérieur, dans un milieu conducteur ionique (solution aqueuse ionique ou sel fondu) appelé \textbf{électrolyte}.
\end{definition}

\medskip\noindent
Pour qu'il y ait électrolyse, il faut trois ingrédients indispensables :
\begin{enumerate}
    \item Un \textbf{électrolyte} : un milieu conducteur ionique (solution de sel, d'acide ou de base, ou sel fondu) contenant des ions mobiles (cations et anions).
    \item Deux \textbf{électrodes} (conducteurs électroniques : graphite, platine, cuivre, acier, etc.) plongées dans l'électrolyte sans se toucher.
    \item Un \textbf{générateur de tension continue} capable d'imposer une tension \textbf{supérieure à une tension minimale} appelée \textbf{tension de seuil} (ou tension de décomposition).
\end{enumerate}

\begin{attention}[Piège classique : « L'électrolyse, c'est quand ça bout »]
Non ! L'électrolyse n'est pas l'ébullition. L'ébullition est un changement d'état \emph{physique} (liquide $\rightarrow$ gaz) dû à la chaleur. L'électrolyse est une \textbf{transformation chimique} (réaction d'oxydoréduction) due au passage du courant. Certes, le courant chauffe la solution (effet Joule), mais ce sont les réactions chimiques aux électrodes qui constituent l'objectif.
\end{attention}

\courssub{Le montage d'électrolyse : description et schématisation}

\noindent
Le montage type comprend :
\begin{itemize}
    \item Un \textbf{générateur de tension continue} ($G$), avec son pôle $+$ et son pôle $-$.
    \item Un \textbf{électrolyseur} : un bécher ou un tube en U contenant l'électrolyte.
    \item Deux \textbf{électrodes} : l'\textbf{anode} (reliée au pôle $+$ du générateur) et la \textbf{cathode} (reliée au pôle $-$ du générateur).
    \item Un \textbf{interrupteur} ($K$) pour ouvrir ou fermer le circuit.
    \item Un \textbf{ampèremètre} ($A$) monté en série pour mesurer l'intensité du courant.
    \item (Optionnel) Un \textbf{voltmètre} ($V$) monté en dérivation aux bornes de l'électrolyseur.
\end{itemize}

\begin{popfigure}
\begin{tikzpicture}[>=Stealth, thick, font=\small]
% Cuve (tube en U)
\draw[very thick, popdark] (0,3) -- (0,0) -- (4,0) -- (4,3);
% Électrolyte
\fill[popblueL, opacity=0.6] (0.05,0.05) rectangle (3.95,2.2);
\draw[popblue, dashed] (0,2.2) -- (4,2.2);
\node[popblue, anchor=west] at (4.1,2.2) {surface de l'électrolyte};
% Électrodes
\fill[popdark] (0.9,1.0) rectangle (1.2,3.6);
\fill[popdark] (2.8,1.0) rectangle (3.1,3.6);
\node[above, poporange, font=\small\bfseries] at (1.05,3.6) {Anode ($+$)};
\node[above, poppurple, font=\small\bfseries] at (2.95,3.6) {Cathode ($-$)};
% Générateur
\draw (1.05,3.6) -- (1.05,4.6) -- (2.0,4.6);
\draw (2.95,3.6) -- (2.95,4.6) -- (2.6,4.6);
\draw[fill=popcream] (2.0,4.3) rectangle (2.6,4.9);
\node at (2.3,4.6) {$G$};
\node[left, poporange] at (2.0,4.6) {$+$};
\node[right, poppurple] at (2.6,4.6) {$-$};
% Ampèremètre et interrupteur sur le fil
\draw[fill=popcream] (0.2,4.3) circle (0.3);
\node at (0.2,4.3) {$A$};
\draw (1.05,4.6) -- (0.5,4.6);
\draw (-0.1,4.6) -- (-1.2,4.6) -- (-1.2,3.6);
\draw[fill=popdark] (-1.35,3.4) rectangle (-1.05,3.8);
\node[left, align=right] at (-1.4,3.6) {électrode\\(anode)};
% Ions dans la solution
\node[popink] at (0.7,1.7) {$\oplus$};
\node[popink] at (1.6,0.5) {$\oplus$};
\node[popink] at (2.3,1.9) {$\oplus$};
\node[popblue] at (1.5,1.3) {$\ominus$};
\node[popblue] at (2.4,0.6) {$\ominus$};
\node[popblue] at (3.4,1.6) {$\ominus$};
% Migrations ioniques
\draw[->, popink, very thick] (2.3,1.7) -- (2.85,1.4);
\draw[->, popblue, very thick] (1.6,1.1) -- (1.15,1.35);
\node[popink, align=center] at (3.55,0.5) {cations\\$\rightarrow$ cathode};
\node[popblue, align=center] at (0.55,0.55) {anions\\$\rightarrow$ anode};
% Sens du courant extérieur
\draw[->, poporange, very thick] (0.6,5.2) -- (1.9,5.2);
\node[poporange, above] at (1.25,5.2) {courant $I$};
\draw[->, poppurple, very thick] (3.4,5.2) -- (2.7,5.2);
\node[poppurple, above] at (3.05,5.2) {$e^-$};
\end{tikzpicture}
\legende{Schéma d'un montage d'électrolyse : générateur $G$, ampèremètre $A$, électrolyseur en U contenant l'électrolyte et deux électrodes. L'anode est reliée au pôle $+$ du générateur, la cathode au pôle $-$. Dans la solution, les cations migrent vers la cathode et les anions vers l'anode.}
\end{popfigure}

\courssub{Polarité des électrodes et sens des courants}

\noindent
C'est un point \textbf{crucial} où beaucoup d'élèves se trompent. Retiens cette règle d'or, valable \textbf{toujours} (pour la pile \textbf{et} pour l'électrolyse) :

\begin{aretenir}[Règle A-O / C-R]
\begin{itemize}
    \item \textbf{À l'anode}, il se produit \textbf{toujours une oxydation} (perte d'électrons). Moyen mnémotechnique : \textbf{A}node = \textbf{O}xydation (\textbf{A-O}).
    \item \textbf{À la cathode}, il se produit \textbf{toujours une réduction} (gain d'électrons). Moyen mnémotechnique : \textbf{C}athode = \textbf{R}éduction (\textbf{C-R}).
\end{itemize}
\end{aretenir}

\medskip\noindent
Ce qui \textbf{change} entre la pile et l'électrolyse, c'est la \textbf{polarité} (le signe $+$ ou $-$) des électrodes :

\begin{center}
\begin{tabularx}{\linewidth}{|>{\bfseries}l|X|X|}
\hline
\rowcolor{popblueL}
\textbf{Caractéristique} & \textbf{Pile (spontanée)} & \textbf{Électrolyse (forcée)} \\
\hline
Rôle du dispositif & Générateur d'énergie électrique & Récepteur d'énergie électrique \\
\hline
Anode (oxydation) & Pôle $-$ & Pôle $+$ \\
\hline
Cathode (réduction) & Pôle $+$ & Pôle $-$ \\
\hline
Sens du courant $I$ (circuit extérieur) & de la cathode ($+$) vers l'anode ($-$) & de l'anode ($+$) vers la cathode ($-$) \\
\hline
Sens des électrons (circuit extérieur) & de l'anode vers la cathode & de la cathode vers l'anode \\
\hline
\end{tabularx}
\end{center}

\begin{methode}[Identifier l'anode et la cathode en électrolyse]
\begin{enumerate}
    \item Repère le \textbf{générateur} : son pôle $+$ et son pôle $-$.
    \item Suis le fil qui part du \textbf{pôle $+$} du générateur : l'électrode à laquelle il aboutit est l'\textbf{anode}.
    \item L'autre électrode (reliée au pôle $-$ du générateur) est la \textbf{cathode}.
    \item \textbf{Vérification :} à l'anode, tu dois écrire une oxydation ; à la cathode, une réduction.
\end{enumerate}
\end{methode}

\medskip\noindent
\textbf{Dans le circuit extérieur (fils métalliques) :}
\begin{itemize}
    \item Le \textbf{courant conventionnel} $I$ sort du pôle $+$ du générateur, arrive à l'anode, traverse l'électrolyseur, ressort par la cathode et revient au pôle $-$ du générateur.
    \item Les \textbf{électrons} (porteurs de charge réels dans les métaux) font le chemin inverse : ils sortent du pôle $-$ du générateur, arrivent à la \textbf{cathode} (où ils sont consommés par la réduction), tandis que l'oxydation à l'\textbf{anode} libère des électrons qui remontent vers le pôle $+$ du générateur.
\end{itemize}

\medskip\noindent
\textbf{Dans la solution (électrolyte) :} le courant est porté par les \textbf{ions} :
\begin{itemize}
    \item les \textbf{cations} (ions positifs : \ce{Cu^2+}, \ce{H+}, \ce{Na+}, \ce{Al^3+}\ldots) sont attirés par la \textbf{cathode} (pôle $-$) ;
    \item les \textbf{anions} (ions négatifs : \ce{Cl-}, \ce{OH-}, \ce{SO4^2-}\ldots) sont attirés par l'\textbf{anode} (pôle $+$).
\end{itemize}
Cette double migration assure le passage du courant \emph{dans} la solution.

\courssub{Condition de démarrage : la tension de seuil}

\noindent
Tu as branché le montage, tu fermes l'interrupteur\ldots et \textbf{rien ne se passe} si la tension du générateur est trop faible ! Pourquoi ?

\medskip\noindent
Pour qu'une réaction d'oxydoréduction non spontanée se produise, il faut lui apporter de l'énergie. Le générateur doit « pousser » les électrons avec une force suffisante pour vaincre l'opposition de la réaction. La tension minimale nécessaire s'appelle la \textbf{tension de seuil} (ou \textbf{tension de décomposition}).

\begin{propriete}[Tension de seuil]
L'électrolyse ne démarre réellement (courant mesurable, réactions observables aux électrodes) que si la tension $U$ imposée par le générateur aux bornes de l'électrolyseur vérifie :
\[
U \geq U_{\text{seuil}}
\]
$U_{\text{seuil}}$ dépend de la nature de l'électrolyte, des concentrations, de la température et de la nature des électrodes (phénomène de \textbf{surtension}, étudié dans la section 4).
\end{propriete}

\medskip\noindent
\textbf{Exemple :} pour décomposer l'eau ($\ce{2 H2O -> 2 H2 + O2}$), la tension minimale théorique est d'environ \SI{1,23}{V} à \SI{25}{\celsius}. En pratique, à cause des surtensions aux électrodes, il faut souvent appliquer une tension nettement supérieure.

\courssub{Exemple fondateur : électrolyse d'une solution de chlorure de cuivre(II)}

\begin{experience}[Électrolyse d'une solution aqueuse de \ce{CuCl2} avec électrodes de graphite]
\textbf{Matériel :} tube en U, solution aqueuse de chlorure de cuivre(II) \ce{CuCl2} (bleue), deux électrodes de graphite (mines de crayon), générateur de tension continue, fils, interrupteur.
\textbf{Protocole :} plonger les électrodes sans qu'elles se touchent, fermer le circuit, observer chaque électrode.
\textbf{Observations :}
\begin{itemize}
    \item \textbf{À la cathode (pôle $-$)} : un dépôt \textbf{rouge-orangé} se forme sur le graphite : c'est du \textbf{cuivre métal} \ce{Cu(s)}.
    \item \textbf{À l'anode (pôle $+$)} : un dégagement gazeux \textbf{vert-jaune} à odeur âcre, qui décolore un papier imbibé d'encre : c'est du \textbf{dichlore} \ce{Cl2(g)}.
    \item La solution \textbf{se décolore} progressivement : les ions \ce{Cu^2+} (bleus) disparaissent.
\end{itemize}
\textbf{Interprétation :} l'électrolyte contient les ions \ce{Cu^2+} et \ce{Cl-}, ainsi que l'eau.
\begin{itemize}
    \item \textbf{Cathode (réduction, C-R) :} le dépôt de cuivre montre que ce sont les ions \ce{Cu^2+} qui sont réduits :
    \[ \ce{Cu^2+(aq) + 2 e- -> Cu(s)} \]
    \item \textbf{Anode (oxydation, A-O) :} le dégagement de dichlore montre que ce sont les ions \ce{Cl-} qui sont oxydés :
    \[ \ce{2 Cl-(aq) -> Cl2(g) + 2 e-} \]
\end{itemize}
\textbf{Bilan global} (somme des deux demi-équations, les électrons se simplifient) :
\[ \ce{Cu^2+(aq) + 2 Cl-(aq) -> Cu(s) + Cl2(g)} \]
Le chlorure de cuivre(II) est décomposé en ses éléments : la transformation, non spontanée, a été \emph{forcée} par le courant.
\end{experience}

\begin{exoresolu}[Électrolyse d'une solution de sulfate de cuivre(II)]
\textbf{Énoncé :} On électrolyse une solution aqueuse de sulfate de cuivre(II) \ce{CuSO4} avec des électrodes de graphite. On observe un dépôt rouge à la cathode et, à l'anode, un dégagement de gaz incolore qui ravive la combustion d'une bûchette incandescente. Écrire les demi-équations aux électrodes et l'équation bilan.
\tcblower
\textbf{Solution détaillée :}
\begin{enumerate}
    \item \textbf{Espèces présentes en solution :} \ce{Cu^2+}, \ce{SO4^2-} et \ce{H2O}.
    \item \textbf{Cathode (pôle $-$, réduction) :} le dépôt rouge est du cuivre métal ; les ions \ce{Cu^2+} sont réduits :
    \[ \ce{Cu^2+(aq) + 2 e- -> Cu(s)} \]
    \item \textbf{Anode (pôle $+$, oxydation) :} le gaz qui ravive une bûchette incandescente est le dioxygène \ce{O2}. Les ions sulfate ne sont pas oxydables dans ces conditions : c'est l'\textbf{eau} qui est oxydée :
    \[ \ce{2 H2O(l) -> O2(g) + 4 H+(aq) + 4 e-} \]
    \item \textbf{Bilan :} on multiplie la demi-équation cathodique par 2 pour égaliser les électrons échangés ($4\,e^-$) :
    \[ \ce{2 Cu^2+(aq) + 2 H2O(l) -> 2 Cu(s) + O2(g) + 4 H+(aq)} \]
\end{enumerate}
\textit{Remarques :} la solution se décolore (disparition de \ce{Cu^2+}) et s'acidifie (apparition de \ce{H+}).
\end{exoresolu}

\courssub{Application industrielle majeure : ALUCAM et l'aluminium}

\begin{savaistu}[L'aluminium camerounais, enfant du Sanaga]
L'usine \textbf{ALUCAM} (Aluminium du Cameroun), implantée à \textbf{Édéa} depuis 1957, produit de l'aluminium primaire par le procédé \textbf{Hall--Héroult} : une électrolyse à haute température (environ \SI{960}{\celsius}) de l'\textbf{alumine} \ce{Al2O3} dissoute dans de la \textbf{cryolithe} fondue (\ce{Na3AlF6}), qui abaisse la température de fusion du mélange. Les cuves sont tapissées de carbone (cathode) et traversées par des anodes en carbone.
\begin{itemize}
    \item \textbf{Cathode (réduction) :} $\ce{Al^3+ + 3 e- -> Al(l)}$ : l'aluminium liquide s'accumule au fond de la cuve.
    \item \textbf{Anode (oxydation) :} les ions oxyde réagissent avec le carbone des anodes : $\ce{2 O^2- + C(s) -> CO2(g) + 4 e-}$. Les anodes \textbf{se consomment} et doivent être remplacées régulièrement.
    \item \textbf{Bilan :} $\ce{2 Al2O3 + 3 C -> 4 Al + 3 CO2}$.
\end{itemize}
L'énergie électrique provient essentiellement du barrage hydroélectrique d'\textbf{Édéa} sur le \textbf{Sanaga} : l'électrolyse de l'aluminium est l'un des procédés industriels les plus gourmands en électricité, ce qui explique l'implantation de l'usine à proximité immédiate de la centrale.
\end{savaistu}

\courssub{Synthèse des idées forces de la section}

\begin{aretenir}[Résumé : généralités sur l'électrolyse]
\begin{enumerate}
    \item \textbf{Nature :} transformation d'oxydoréduction \textbf{forcée} (non spontanée), imposée par un courant électrique.
    \item \textbf{Montage :} générateur de tension continue, électrolyseur (électrolyte + deux électrodes), ampèremètre en série, interrupteur. L'électrolyseur est un \textbf{récepteur}.
    \item \textbf{Condition de démarrage :} $U_{\text{générateur}} \geq U_{\text{seuil}}$.
    \item \textbf{Polarité :} l'anode est reliée au pôle $+$ du générateur ; la cathode au pôle $-$.
    \item \textbf{Règle universelle :} \textbf{anode = oxydation (A-O)} ; \textbf{cathode = réduction (C-R)} (valable pour la pile et pour l'électrolyse).
    \item \textbf{Conduction :} électrons dans les fils ; ions dans la solution (cations vers la cathode, anions vers l'anode).
    \item \textbf{Exemple type :} \ce{CuCl2} aqueux avec électrodes de graphite $\rightarrow$ \ce{Cu(s)} à la cathode et \ce{Cl2(g)} à l'anode.
\end{enumerate}
\end{aretenir}

\begin{exercice}[Vérification des acquis -- Section 1]
\begin{enumerate}
    \item Reproduis et complète le tableau comparatif pile / électrolyse (polarité des électrodes, rôle du dispositif, sens du courant et des électrons).
    \item En électrolyse, l'anode est-elle toujours le pôle $+$ ? Justifie.
    \item Pourquoi une solution aqueuse de \ce{NaCl} ne se décompose-t-elle pas spontanément en sodium et en dichlore ? Quel dispositif permet de forcer une telle transformation ?
    \item À la cathode d'une électrolyse, les électrons arrivent-ils par le fil ou par la solution ? Quelle espèce chimique les reçoit ?
    \item À l'usine ALUCAM, les anodes en carbone se consomment. Écris la demi-équation anodique et explique pourquoi il faut remplacer les anodes périodiquement.
\end{enumerate}
\end{exercice}

\begin{corrige}[Exercice Section 1]
\begin{enumerate}
    \item Voir le tableau du cours (paragraphe « Polarité des électrodes »). Points clés : en électrolyse, anode = pôle $+$ et cathode = pôle $-$ ; c'est l'inverse dans une pile. Dans les deux cas : oxydation à l'anode, réduction à la cathode.
    \item \textbf{Oui.} En électrolyse, l'anode est par construction reliée au pôle $+$ du générateur : c'est donc toujours le pôle $+$ de l'électrolyseur. (Dans une pile, c'est l'inverse : l'anode est le pôle $-$.)
    \item La réaction $\ce{2 NaCl -> 2 Na + Cl2}$ n'est pas spontanée : elle ne peut se produire que si un générateur extérieur impose une tension supérieure à la tension de seuil (électrolyse). \textit{Précision :} en solution aqueuse, ce n'est d'ailleurs pas \ce{Na+} qui est réduit à la cathode, mais l'eau (dégagement de \ce{H2}) ; pour obtenir du sodium métal, il faut électrolyser du \ce{NaCl} \textbf{fondu}.
    \item Les électrons arrivent \textbf{par le fil métallique}, depuis le pôle $-$ du générateur. Ils sont captés à la surface de la cathode par l'espèce \textbf{oxydante} présente (un cation comme \ce{Cu^2+}, ou l'eau), qui subit alors une \textbf{réduction}.
    \item Demi-équation anodique : $\ce{2 O^2- + C(s) -> CO2(g) + 4 e-}$. Le carbone de l'anode est consommé par la réaction avec les ions oxyde pour former du dioxyde de carbone gazeux : l'anode s'use progressivement et doit être remplacée.
\end{enumerate}
\end{corrige}

\coursec{Réactions aux électrodes}

Dans la section précédente, nous avons découvert l'électrolyse comme une \cle{transformation forcée} : le générateur impose le sens de circulation des électrons et « oblige » la réaction à se produire dans le sens inverse de sa spontanéité. Nous avons aussi appris à repérer les deux électrodes : la \cle{cathode}, reliée au pôle négatif du générateur, siège d'une réduction ; l'\cle{anode}, reliée au pôle positif, siège d'une oxydation. Mais concrètement, \emph{quelles espèces chimiques réagissent à chaque électrode ?} Et comment passer des deux demi-équations électroniques à l'équation bilan de l'électrolyse ? C'est toute l'ambition de cette section : décrire précisément ce qui se joue à la surface de chaque électrode, distinguer le cas des électrodes inertes de celui des électrodes attaquables, et apprendre à écrire proprement les équations. Tu verras que tout repose sur une idée simple : la cathode est une « pompe à électrons » qui les distribue, l'anode est un « aspirateur à électrons » qui les prélève.

\courssub{Réduction cathodique : la cathode donne des électrons}

\begin{definition}[Réduction cathodique]
À la \cle{cathode} (électrode reliée au pôle $-$ du générateur), il se produit toujours une \cle{réduction} : une espèce chimique présente au voisinage de l'électrode \textbf{gagne un ou plusieurs électrons} fournis par le générateur. La demi-équation électronique s'écrit toujours sous la forme :
\[ \text{Ox} + n\,e^- \ce{->} \text{Red} \]
L'oxydant du couple est réduit. L'électrode elle-même ne fait que transporter les électrons : elle n'apparaît jamais dans la demi-équation.
\end{definition}

L'intuition est la suivante : le pôle négatif du générateur « refoule » des électrons vers la cathode, qui s'enrichit donc en électrons. Toute espèce oxydante qui s'approche de sa surface peut capter ces électrons disponibles. En solution aqueuse, les candidats à la réduction sont :
\begin{itemize}
\item les \textbf{cations métalliques} : $\ce{Cu^2+} + 2e^- \ce{->} \ce{Cu}$ (dépôt métallique rougeâtre), $\ce{Zn^2+} + 2e^- \ce{->} \ce{Zn}$, $\ce{Ag+} + e^- \ce{->} \ce{Ag}$ ;
\item les \textbf{ions oxonium} (ou protons) : $\ce{2H+} + 2e^- \ce{->} \ce{H2}$ (dégagement gazeux, effervescence) ;
\item les molécules oxydantes dissoutes, le cas échéant.
\end{itemize}

\begin{attention}[La cathode n'est jamais consommée]
La cathode fournit des électrons mais ne participe pas chimiquement à la réaction : sa masse ne diminue jamais par oxydation. En revanche, sa masse peut \textbf{augmenter} si un métal se dépose à sa surface (c'est le principe de la galvanoplastie et du raffinage). Dire « la cathode s'oxyde » est une erreur rédhibitoire au Probatoire : à la cathode, c'est toujours une \textbf{réduction}.
\end{attention}

\courssub{Oxydation anodique : l'anode prélève des électrons}

\begin{definition}[Oxydation anodique]
À l'\cle{anode} (électrode reliée au pôle $+$ du générateur), il se produit toujours une \cle{oxydation} : une espèce chimique \textbf{perd un ou plusieurs électrons}, captés par le circuit électrique. La demi-équation s'écrit :
\[ \text{Red} \ce{->} \text{Ox} + n\,e^- \]
Le réducteur du couple est oxydé.
\end{definition}

Le pôle positif du générateur « aspire » les électrons : l'anode se trouve appauvrie en électrons et les prélève à toute espèce réductrice qui touche sa surface. Mais ici apparaît une subtilité capitale : \textbf{l'anode elle-même peut être cette espèce réductrice}. Tout dépend de la nature du matériau dont elle est faite. C'est l'objet des deux paragraphes suivants.

\courssub{Électrodes inertes : seules les espèces de la solution réagissent}

\begin{definition}[Électrode inerte]
Une électrode est dite \cle{inerte} (ou \emph{inattaquable}) lorsque son matériau ne participe à \textbf{aucune} réaction électrochimique : il se contente d'assurer le passage des électrons entre le circuit et la solution. Les deux électrodes inertes classiques sont le \cle{graphite} (carbone) et le \cle{platine}.
\end{definition}

Avec des électrodes inertes, les seules espèces qui peuvent réagir sont celles présentes dans la solution : les ions du soluté, l'eau, éventuellement le solvant et les molécules dissoutes. L'électrode est un simple « guichet » d'échange d'électrons.

\begin{experience}[Électrolyse d'une solution aqueuse d'acide chlorhydrique]
\textbf{Dispositif :} un tube en U contient une solution d'acide chlorhydrique ($\ce{H+} + \ce{Cl-}$) de concentration voisine de \SI{1}{mol.L^{-1}} ; deux électrodes de graphite plongent dans la solution et sont reliées à un générateur de tension continue (environ \SI{6}{V}). Deux tubes retournés sur les électrodes recueillent les gaz.

\textbf{Observations :}
\begin{itemize}
\item à la \textbf{cathode} : dégagement gazeux abondant ; le gaz recueilli produit une petite détonation à l'approche d'une flamme : c'est le dihydrogène $\ce{H2}$ ;
\item à l'\textbf{anode} : dégagement d'un gaz jaune-vert, âcre, qui bleuit un papier imbibé d'iodure de potassium amidonné : c'est le dichlore $\ce{Cl2}$.
\end{itemize}

\textbf{Interprétation :} le graphite est inerte, donc ce sont les ions de la solution qui réagissent :
\begin{itemize}
\item réduction cathodique des ions oxonium : $\ce{2H+} + 2e^- \ce{->} \ce{H2}$ ;
\item oxydation anodique des ions chlorure : $\ce{2Cl-} \ce{->} \ce{Cl2} + 2e^-$.
\end{itemize}
\end{experience}

\begin{popfigure}
\begin{center}
\begin{tikzpicture}[scale=0.95, every node/.style={font=\small}]
% Tube en U simplifié
\draw[thick, popblue] (-2.2,0) -- (-2.2,2.6);
\draw[thick, popblue] (2.2,0) -- (2.2,2.6);
\draw[thick, popblue] (-2.2,0) .. controls (-2.2,-1.3) and (2.2,-1.3) .. (2.2,0);
% Solution
\fill[popblue!20] (-2.1,0.1) .. controls (-2.1,-1.2) and (2.1,-1.2) .. (2.1,0.1) -- (2.1,1.9) -- (-2.1,1.9) -- cycle;
\draw[popblue!60] (-2.1,1.9) -- (2.1,1.9);
\node[popdark] at (0,0.6) {$\ce{H+} + \ce{Cl-}$ (solution aqueuse)};
% Électrodes graphite
\fill[popdark!80] (-1.5,1.2) rectangle (-1.1,3.4);
\fill[popdark!80] (1.1,1.2) rectangle (1.5,3.4);
\node[above, popdark] at (-1.3,3.4) {Cathode (graphite)};
\node[above, popdark] at (1.3,3.4) {Anode (graphite)};
% Générateur
\draw[thick] (-1.3,3.4) -- (-1.3,4.3) -- (-0.45,4.3);
\draw[thick] (1.3,3.4) -- (1.3,4.3) -- (0.45,4.3);
\draw[thick] (-0.45,4.15) rectangle (0.45,4.45);
\node at (0,4.3) {\footnotesize G};
\node[popdark] at (-0.75,4.6) {$-$};
\node[popdark] at (0.75,4.6) {$+$};
% Bulles et demi-équations
\foreach \y in {1.5,1.9,2.3} {\draw[popgreen!70, fill=popgreen!30] (-1.3,\y) circle (0.06);}
\foreach \y in {1.5,1.9,2.3} {\draw[popgold!80, fill=popgold!40] (1.3,\y) circle (0.06);}
\node[align=center, popgreen!60!black, below] at (-3.6,1.0) {$\ce{2H+} + 2e^- \ce{->} \ce{H2}$\\ \footnotesize réduction};
\draw[-{Stealth}, popgreen!60!black] (-2.6,1.15) -- (-1.55,1.5);
\node[align=center, popgold!60!black, below] at (3.7,1.0) {$\ce{2Cl-} \ce{->} \ce{Cl2} + 2e^-$\\ \footnotesize oxydation};
\draw[-{Stealth}, popgold!60!black] (2.6,1.15) -- (1.55,1.5);
% Sens du courant ionique
\draw[-{Stealth}, thick, poporange] (0.2,-1.6) -- (-0.8,-1.6);
\node[poporange, below] at (-0.3,-1.6) {\footnotesize migration des ions};
\end{tikzpicture}
\end{center}
\legende{Électrolyse de l'acide chlorhydrique avec électrodes de graphite inertes : dégagement de $\ce{H2}$ à la cathode et de $\ce{Cl2}$ à l'anode.}
\end{popfigure}

\courssub{Électrodes attaquables : le phénomène d'anode soluble}

Changeons maintenant de matériau. Si l'anode est constituée d'un métal \textbf{réducteur} comme le cuivre, le zinc ou le fer, une nouvelle possibilité s'offre au système : au lieu d'oxyder une espèce de la solution, le générateur peut oxyder \emph{le métal de l'anode lui-même}. L'anode se dissout alors progressivement : on parle d'\cle{anode soluble}.

\begin{definition}[Anode soluble]
Une \cle{anode soluble} (ou \emph{anode attaquable}) est une anode constituée d'un métal qui \textbf{participe à la réaction d'électrolyse en s'oxydant} : les atomes métalliques perdent des électrons et passent en solution sous forme de cations. La masse de l'anode \textbf{diminue} au cours de l'électrolyse. Exemple avec une anode de cuivre :
\[ \ce{Cu} \ce{->} \ce{Cu^2+} + 2e^- \]
\end{definition}

\begin{experience}[Électrolyse d'une solution de sulfate de cuivre avec électrodes de cuivre]
\textbf{Dispositif :} un bécher contient une solution bleue de sulfate de cuivre ($\ce{Cu^2+} + \ce{SO4^2-}$) ; deux électrodes de \textbf{cuivre} sont reliées au générateur.

\textbf{Observations :}
\begin{itemize}
\item à la \textbf{cathode} : dépôt rougeâtre de cuivre métallique ; la masse de la cathode \textbf{augmente} ;
\item à l'\textbf{anode} : la lame de cuivre s'amincit ; sa masse \textbf{diminue} ;
\item la teinte bleue de la solution reste \textbf{inchangée} : la concentration en ions $\ce{Cu^2+}$ est constante, car chaque ion consommé à la cathode est remplacé par un ion produit à l'anode.
\end{itemize}

\textbf{Interprétation :}
\begin{itemize}
\item cathode (réduction) : $\ce{Cu^2+} + 2e^- \ce{->} \ce{Cu}$ ;
\item anode (oxydation du métal) : $\ce{Cu} \ce{->} \ce{Cu^2+} + 2e^-$.
\end{itemize}
Tout se passe comme si le cuivre était \emph{transporté} de l'anode vers la cathode à travers la solution.
\end{experience}

\begin{popfigure}
\begin{center}
\begin{tikzpicture}[scale=0.95, every node/.style={font=\small}]
% Bécher
\draw[thick, popblue] (-3,0) -- (-3,3.2);
\draw[thick, popblue] (3,0) -- (3,3.2);
\draw[thick, popblue] (-3,0) -- (3,0);
% Solution
\fill[popblue!20] (-2.9,0.05) rectangle (2.9,2.3);
\draw[popblue!60] (-2.9,2.3) -- (2.9,2.3);
\node[popdark] at (0,0.45) {Solution de $\ce{CuSO4}$ : $\ce{Cu^2+} + \ce{SO4^2-}$};
% Anode (cuivre, qui s'amincit)
\fill[poporange!70] (1.5,1.0) rectangle (1.75,3.9);
\draw[dashed, popdark] (1.4,1.0) rectangle (1.85,3.9);
\node[above, poporange!80!black] at (1.62,3.95) {Anode (Cu)};
\node[poporange!80!black, align=center] at (4.4,2.6) {\footnotesize l'anode\\ \footnotesize s'amincit};
\draw[-{Stealth}, poporange!80!black] (3.7,2.5) -- (1.9,2.3);
% Cathode (dépôt)
\fill[popdark!60] (-1.75,1.0) rectangle (-1.5,3.9);
\fill[popgold!80] (-1.85,1.0) rectangle (-1.4,2.1);
\node[above, popdark] at (-1.62,3.95) {Cathode (Cu)};
\node[popgold!60!black, align=center] at (-4.3,2.6) {\footnotesize dépôt de Cu\\ \footnotesize (masse $\nearrow$)};
\draw[-{Stealth}, popgold!60!black] (-3.5,2.4) -- (-1.95,1.7);
% Générateur
\draw[thick] (-1.62,3.9) -- (-1.62,4.6) -- (-0.45,4.6);
\draw[thick] (1.62,3.9) -- (1.62,4.6) -- (0.45,4.6);
\draw[thick] (-0.45,4.45) rectangle (0.45,4.75);
\node at (0,4.6) {\footnotesize G};
\node[popdark] at (-0.72,4.9) {$-$};
\node[popdark] at (0.72,4.9) {$+$};
% Migration des ions Cu2+
\draw[-{Stealth}, very thick, poppurple] (0.9,1.4) .. controls (0,1.9) .. (-0.9,1.5);
\node[poppurple, above] at (0,1.95) {\footnotesize $\ce{Cu^2+}$};
% Demi-équations
\node[align=center, popdark] at (-4.4,0.7) {$\ce{Cu^2+} + 2e^- \ce{->} \ce{Cu}$};
\node[align=center, popdark] at (4.4,0.7) {$\ce{Cu} \ce{->} \ce{Cu^2+} + 2e^-$};
\end{tikzpicture}
\end{center}
\legende{Électrolyse de $\ce{CuSO4}$ à anode soluble : le cuivre se dissout à l'anode et se dépose à la cathode ; la concentration en $\ce{Cu^2+}$ reste constante.}
\end{popfigure}

\begin{savaistu}[Le raffinage électrolytique du cuivre]
Le cuivre issu de la métallurgie contient des impuretés (fer, zinc, argent, or\ldots) qui dégradent sa conductivité. Pour le purifier, l'industrie utilise précisément l'électrolyse à anode soluble : l'anode est une plaque de \textbf{cuivre impur} qui se dissout, tandis que du \textbf{cuivre pur à \SI{99,99}{\percent}} se dépose sur la cathode. Les impuretés précieuses (or, argent) tombent au fond de la cuve en « boues anodiques » que l'on récupère ! C'est ce cuivre ultra-pur qui est indispensable à la fabrication des câbles électriques, car même de traces d'impuretés augmentent fortement la résistance électrique. Au Cameroun, les réseaux de distribution d'ENEO et les installations artisanales de câblage dépendent de ce cuivre raffiné.
\end{savaistu}

\courssub{Écrire l'équation bilan de l'électrolyse}

Une électrolyse met en jeu \textbf{deux} demi-équations électroniques simultanées, une à chaque électrode. L'équation bilan rend compte de la transformation globale : elle s'obtient en additionnant les deux demi-équations après avoir égalisé le nombre d'électrons échangés.

\begin{methode}[Écrire l'équation bilan d'une électrolyse]
\begin{enumerate}
\item \textbf{Identifier} les espèces présentes (ions du soluté, eau, métal des électrodes) et repérer quelle électrode est la cathode (réduction) et quelle est l'anode (oxydation).
\item \textbf{Écrire} la demi-équation de réduction à la cathode : $\text{Ox}_1 + n_1 e^- \ce{->} \text{Red}_1$.
\item \textbf{Écrire} la demi-équation d'oxydation à l'anode : $\text{Red}_2 \ce{->} \text{Ox}_2 + n_2 e^-$.
\item \textbf{Multiplier} chaque demi-équation par un coefficient tel que le nombre d'électrons cédés à l'anode soit égal au nombre d'électrons captés à la cathode (plus petit multiple commun de $n_1$ et $n_2$).
\item \textbf{Additionner} membre à membre et \textbf{simplifier} : les électrons doivent disparaître totalement de l'équation bilan.
\item \textbf{Vérifier} la conservation des éléments et des charges électriques dans le bilan final.
\end{enumerate}
\end{methode}

\begin{exemplebox}[Électrolyse d'une solution de chlorure de cuivre(II)]
On électrolyse une solution de chlorure de cuivre(II) $\ce{CuCl2}$ ($\ce{Cu^2+ + 2Cl-}$) entre deux électrodes de graphite inertes.

\textbf{Étape 1 — Espèces présentes :} $\ce{Cu^2+}$, $\ce{Cl-}$, $\ce{H2O}$. Le graphite ne réagit pas.

\textbf{Étape 2 — Cathode (réduction) :} les ions cuivre(II) captent les électrons :
\[ \ce{Cu^2+ + 2e- -> Cu} \]
On observe un dépôt rougeâtre de cuivre.

\textbf{Étape 3 — Anode (oxydation) :} les ions chlorure cèdent des électrons :
\[ \ce{2Cl- -> Cl2 + 2e-} \]
On observe un dégagement de dichlore, gaz jaune-vert suffocant.

\textbf{Étape 4 — Électrons :} les deux demi-équations échangent déjà le même nombre d'électrons ($n = 2$) : aucune multiplication n'est nécessaire.

\textbf{Étape 5 — Addition :} on additionne les deux demi-équations ; les électrons se simplifient :
\[ \ce{Cu^2+ + 2e- + 2Cl- -> Cu + Cl2 + 2e-} \]
d'où le \textbf{bilan de l'électrolyse} (équation ionique nette) :
\[ \ce{Cu^2+ + 2Cl- -> Cu + Cl2} \]

\textbf{Étape 6 — Vérification :} à gauche, charge totale $= +2 + 2\times(-1) = 0$ ; à droite, $0 + 0 = 0$. Les éléments Cu et Cl sont conservés. L'équation est correcte.
\end{exemplebox}

\begin{attention}[Les électrons ne circulent pas dans la solution !]
C'est l'erreur la plus fréquente au Probatoire. Les électrons circulent \textbf{uniquement dans les fils métalliques et les électrodes}, de l'anode vers le pôle $+$ du générateur, puis du pôle $-$ vers la cathode. \textbf{Dans la solution, ce sont les ions qui assurent le passage du courant} : les cations migrent vers la cathode, les anions migrent vers l'anode. On dit que la conduction est \emph{métallique} (électronique) dans les fils et \emph{ionique} dans l'électrolyte. Autre conséquence : les électrons ne doivent \textbf{jamais} apparaître dans l'équation bilan finale ; s'il en reste, c'est que les demi-équations n'ont pas été correctement multipliées.
\end{attention}

\begin{exoresolu}[Électrolyse de l'acide chlorhydrique : bilan complet]
On réalise l'électrolyse d'une solution aqueuse d'acide chlorhydrique de concentration \SI{1,0}{mol.L^{-1}} entre deux électrodes de graphite. Pendant \SI{20}{minutes}, on recueille \SI{240}{mL} de dichlore à l'anode dans les conditions où le volume molaire vaut $V_m = \SI{24}{L.mol^{-1}}$.
\begin{enumerate}
\item Écrire les demi-équations à chaque électrode et le bilan de l'électrolyse.
\item Calculer la quantité de matière de $\ce{Cl2}$ recueillie, puis la quantité de matière de $\ce{H2}$ dégagée à la cathode et le volume correspondant.
\end{enumerate}
\tcblower
\textbf{1. Demi-équations et bilan.}

Le graphite est inerte : seules les espèces de la solution réagissent.

À la \textbf{cathode}, réduction des ions oxonium :
\[ \ce{2H+} + 2e^- \ce{->} \ce{H2} \]

À l'\textbf{anode}, oxydation des ions chlorure :
\[ \ce{2Cl-} \ce{->} \ce{Cl2} + 2e^- \]

Les deux demi-équations échangent 2 électrons : on les additionne directement. Bilan :
\[ \ce{2H+} + \ce{2Cl-} \ce{->} \ce{H2} + \ce{Cl2} \]

\textbf{2. Quantités de matière.}

Quantité de dichlore recueillie :
\[ n(\ce{Cl2}) = \frac{V(\ce{Cl2})}{V_m} = \frac{\SI{0,240}{L}}{\SI{24}{L.mol^{-1}}} = \SI{0,010}{mol} \]

Dressons le tableau d'avancement :
\begin{center}
\begin{tabular}{lcccc}
\toprule
Équation & \multicolumn{4}{c}{$\ce{2H+} + \ce{2Cl-} \ce{->} \ce{H2} + \ce{Cl2}$} \\
\midrule
État initial (mol) & $n_0(\ce{H+})$ & $n_0(\ce{Cl-})$ & $0$ & $0$ \\
En cours (mol) & $n_0(\ce{H+}) - 2x$ & $n_0(\ce{Cl-}) - 2x$ & $x$ & $x$ \\
\bottomrule
\end{tabular}
\end{center}

D'après les coefficients stœchiométriques, $n(\ce{H2}) = n(\ce{Cl2}) = x$, donc :
\[ n(\ce{H2}) = \SI{0,010}{mol} \]
et le volume de dihydrogène dégagé à la cathode vaut :
\[ V(\ce{H2}) = n(\ce{H2}) \times V_m = 0{,}010 \times 24 = \SI{0,24}{L} = \SI{240}{mL} \]

\textbf{Remarque :} les volumes des deux gaz sont égaux, ce que confirme l'expérience lorsqu'on recueille les gaz dans deux tubes identiques : les niveaux descendent de la même hauteur. C'est une vérification expérimentale élégante de la stœchiométrie du bilan.
\end{exoresolu}

\begin{aretenir}[L'essentiel de la section]
\begin{itemize}
\item À la \textbf{cathode} (pôle $-$) : toujours une \textbf{réduction}, $\text{Ox} + ne^- \ce{->} \text{Red}$.
\item À l'\textbf{anode} (pôle $+$) : toujours une \textbf{oxydation}, $\text{Red} \ce{->} \text{Ox} + ne^-$.
\item Électrode \textbf{inerte} (graphite, platine) : seules les espèces de la solution réagissent.
\item Électrode \textbf{attaquable} (Cu, Zn, Fe) : l'anode peut s'oxyder et se dissoudre — c'est l'\textbf{anode soluble}, base du raffinage des métaux.
\item L'\textbf{équation bilan} est la somme des deux demi-équations après égalisation des électrons ; les électrons n'y apparaissent plus.
\item Dans les fils, le courant est porté par les \textbf{électrons} ; dans la solution, par les \textbf{ions}.
\end{itemize}
\end{aretenir}

Maintenant que tu sais écrire les réactions à chaque électrode, une question se pose naturellement : lorsque \emph{plusieurs} espèces sont susceptibles de réagir à la même électrode (par exemple $\ce{Na+}$ et $\ce{H2O}$ à la cathode d'une solution de sel), laquelle réagit réellement ? C'est l'objet de la section suivante : la prévision des réactions d'électrolyse.

\coursec{Prévision des réactions d'électrolyse}

Dans la section précédente, nous avons appris à écrire les demi-équations qui se déroulent à chaque électrode : une \cle{oxydation} à l'\cle{anode} et une \cle{réduction} à la \cle{cathode}. Mais une question fondamentale reste en suspens : lorsqu'on plonge deux électrodes dans une solution aqueuse, \emph{plusieurs} espèces chimiques sont présentes en même temps — les ions du soluté, les molécules d'eau, parfois le métal de l'électrode elle-même. Laquelle va réagir ? C'est tout l'objet de cette section : apprendre à \cle{prévoir} les réactions d'une électrolyse avant même de brancher le générateur.

\courssub{Le problème posé : plusieurs candidats à chaque électrode}

Prenons une situation concrète. Tu verses du chlorure de sodium \ce{NaCl} (du sel de cuisine, comme celui qu'on récolte à la lagune de Kribi) dans de l'eau distillée, et tu y plonges deux électrodes inertes en graphite. Quelles espèces chimiques se trouvent réellement dans le bécher ?

\begin{itemize}
\item les ions \ce{Na+} et \ce{Cl-} issus de la dissolution du sel : $\ce{NaCl -> Na+ + Cl-}$ ;
\item les molécules d'eau \ce{H2O}, le solvant, présentes en très grande quantité ;
\item le carbone des électrodes de graphite (inerte, il ne réagit pas).
\end{itemize}

À la \textbf{cathode} (électrode reliée au pôle $-$ du générateur), une réduction doit avoir lieu. Deux oxydants sont candidats : \ce{Na+} (qui pourrait donner du sodium métallique \ce{Na}) et \ce{H2O} (qui pourrait donner du dihydrogène \ce{H2}). Lequel gagne ?

À l'\textbf{anode} (électrode reliée au pôle $+$), une oxydation doit avoir lieu. Deux réducteurs sont candidats : \ce{Cl-} (qui pourrait donner du dichlore \ce{Cl2}) et \ce{H2O} (qui pourrait donner du dioxygène \ce{O2}). Lequel gagne ?

On ne peut pas répondre au hasard : il faut un critère. Ce critère, c'est le \cle{potentiel standard} $E^{\circ}$ des couples redox en présence.

\courssub{Les règles de prévision}

Rappelons une idée intuitive déjà rencontrée avec les réactions spontanées : dans une « compétition » entre oxydants, c'est le \emph{plus avide d'électrons} qui les capte ; entre réducteurs, c'est le \emph{plus généreux en électrons} qui les cède. Le potentiel standard mesure exactement cette « avidité » : plus $E^{\circ}$ est \textbf{élevé}, plus l'oxydant du couple est \textbf{fort} ; plus $E^{\circ}$ est \textbf{faible}, plus le réducteur du couple est \textbf{fort}.

\begin{propriete}[Règles de prévision d'une électrolyse]
\begin{itemize}
\item \textbf{À la cathode} (réduction) : se réduit l'\cle{oxydant le plus fort} parmi les espèces présentes, c'est-à-dire celui dont le couple a le potentiel standard $E^{\circ}$ \textbf{le plus élevé}.
\item \textbf{À l'anode} (oxydation) : s'oxyde le \cle{réducteur le plus fort} parmi les espèces présentes, c'est-à-dire celui dont le couple a le potentiel standard $E^{\circ}$ \textbf{le plus faible}.
\end{itemize}
\end{propriete}

\begin{attention}[L'eau, candidate discrète mais omniprésente]
En solution \textbf{aqueuse}, l'eau participe toujours à la compétition, via deux couples redox dont les potentiels standards $E^{\circ}$ (toutes espèces à l'état standard, \ce{H2O} liquide, gaz à \SI{1}{bar}, ions à \SI{1}{mol.L^{-1}}) sont :
\begin{itemize}
\item à la cathode (réduction de l'eau) : $\ce{2 H2O + 2 e- <=> H2 + 2 HO-}$ \quad ($E^{\circ} = \SI{-0,83}{V}$) ;
\item à l'anode (oxydation de l'eau) : $\ce{2 H2O <=> O2 + 4 H+ + 4 e-}$ \quad ($E^{\circ} = \SI{+1,23}{V}$).
\end{itemize}
Ne l'oublie jamais dans ta liste des espèces présentes : c'est l'erreur la plus fréquente au Probatoire.
\end{attention}

\begin{methode}[Prévoir les réactions d'une électrolyse]
\begin{enumerate}
\item \textbf{Lister toutes les espèces présentes} : les ions du soluté, les molécules d'eau \ce{H2O}, et le métal des électrodes si celles-ci ne sont pas inertes (cuivre, zinc, fer...).
\item \textbf{Écrire les couples redox} concernés et \textbf{comparer leurs potentiels standards} $E^{\circ}$ (donnés dans l'énoncé ou dans une table de référence).
\item \textbf{À la cathode} : retenir l'oxydant du couple de $E^{\circ}$ le plus élevé et écrire sa demi-équation de réduction.
\item \textbf{À l'anode} : retenir le réducteur du couple de $E^{\circ}$ le plus faible et écrire sa demi-équation d'oxydation.
\item \textbf{Combiner} les deux demi-équations pour obtenir le bilan de l'électrolyse (le nombre d'électrons échangés doit être le même).
\end{enumerate}
\end{methode}

\courssub{Le diagramme en escalier des potentiels}

Une image vaut mille tableaux. Représentons sur un axe vertical les potentiels standards des couples en présence : les espèces qui se réduisent à la cathode se lisent « en haut » (oxydants forts), celles qui s'oxydent à l'anode « en bas » (réducteurs forts).

\begin{popfigure}
\begin{center}
\begin{tikzpicture}[scale=1.1]
% Axe vertical des potentiels : y = 1.3 * E(V)
\draw[->, thick, popdark] (0,-4.8) -- (0,3.0) node[above] {\small $E^{\circ}$ (V)};
\foreach \y/\v in {-3.9/-3, -2.6/-2, -1.3/-1, 0/0, 1.3/1, 2.6/2}
  \draw[popdark] (-0.15,\y) -- (0.15,\y) node[left=3pt, font=\small] {\v};
% Couples et espèces (positions y = 1.3 * E)
% Cl2/Cl- : 1.36 -> 1.77
\node[anchor=west, font=\small] at (1.0,1.77) {\ce{Cl2}/\ce{Cl-} \quad $E^{\circ}=\SI{+1,36}{V}$};
\fill[poporange] (0.6,1.77) circle (2pt);
% O2/H2O : 1.23 -> 1.60
\node[anchor=west, font=\small] at (1.0,1.60) {\ce{O2}/\ce{H2O} \quad $E^{\circ}=\SI{+1,23}{V}$};
\fill[poporange] (0.6,1.60) circle (2pt);
% Cu2+/Cu : 0.34 -> 0.44
\node[anchor=west, font=\small] at (1.0,0.44) {\ce{Cu^2+}/\ce{Cu} \quad $E^{\circ}=\SI{+0,34}{V}$};
\fill[popblue] (0.6,0.44) circle (2pt);
% 2H+/H2 : 0.00 -> 0.00
\node[anchor=west, font=\small] at (1.0,0.00) {\ce{2H+}/\ce{H2} \quad $E^{\circ}=\SI{0,00}{V}$};
\fill[popblue] (0.6,0.00) circle (2pt);
% H2O/H2 : -0.83 -> -1.08
\node[anchor=west, font=\small] at (1.0,-1.08) {\ce{H2O}/\ce{H2} \quad $E^{\circ}=\SI{-0,83}{V}$};
\fill[popblue] (0.6,-1.08) circle (2pt);
% Na+/Na : -2.71 -> -3.52
\node[anchor=west, font=\small] at (1.0,-3.52) {\ce{Na+}/\ce{Na} \quad $E^{\circ}=\SI{-2,71}{V}$};
\fill[popblue] (0.6,-3.52) circle (2pt);
% Zones colorées
% Zone cathodique (Haut : forts oxydants)
\draw[rounded corners, popblue!60, fill=popblueL] (0.3,0.1) rectangle (8.5,2.0);
\node[font=\small\bfseries, popblue, align=center] at (7.5,1.7) {Zone cathodique : réduction\\ (oxydant le plus haut)};
% Zone anodique (Bas : forts réducteurs)
\draw[rounded corners, poporange!60, fill=poporangeL] (0.3,-4.2) rectangle (8.5,0.05);
\node[font=\small\bfseries, poporange, align=center] at (7.5,-3.5) {Zone anodique : oxydation\\ (réducteur le plus bas)};
% Flèches indicatives
\draw[->, very thick, popblue] (9.0,1.5) -- (9.0,0.3) node[midway, right, font=\small, align=left] {Cathode ($-$) :\\ on choisit en \textbf{haut}};
\draw[->, very thick, poporange] (9.0,-0.5) -- (9.0,-3.5) node[midway, right, font=\small, align=left] {Anode ($+$) :\\ on choisit en \textbf{bas}};
\end{tikzpicture}
\end{center}
\legende{Diagramme en escalier des potentiels standards. L'axe vertical est gradué en volts. À la cathode (zone bleue, hauts potentiels), l'oxydant du couple de potentiel le plus élevé est réduit. À l'anode (zone orange, bas potentiels), le réducteur du couple de potentiel le plus faible est oxydé.}
\end{popfigure}

\courssub{Application 1 : électrolyse d'une solution de chlorure de sodium}

Appliquons la méthode pas à pas à la solution de \ce{NaCl}, avec des électrodes \textbf{inertes} (graphite).

\textbf{Étape 1 — Espèces présentes :} \ce{Na+}, \ce{Cl-}, \ce{H2O}.

\textbf{Étape 2 — Couples candidats et potentiels standards :}
\begin{itemize}
\item à la cathode : $\ce{Na+}/\ce{Na}$ ($E^{\circ} = \SI{-2,71}{V}$) et $\ce{H2O}/\ce{H2}$ ($E^{\circ} = \SI{-0,83}{V}$) ;
\item à l'anode : $\ce{Cl2}/\ce{Cl-}$ ($E^{\circ} = \SI{+1,36}{V}$) et $\ce{O2}/\ce{H2O}$ ($E^{\circ} = \SI{+1,23}{V}$).
\end{itemize}

\textbf{Étape 3 — Cathode :} l'oxydant le plus fort est celui du couple de potentiel le plus élevé : c'est \ce{H2O} ($\SI{-0,83}{V} > \SI{-2,71}{V}$). Ce n'est donc \textbf{pas} le sodium qui se dépose, mais le dihydrogène qui se dégage :
\[ \ce{2 H2O + 2 e- -> H2 + 2 HO-} \]

\textbf{Étape 4 — Anode :} le réducteur le plus fort est celui du couple de potentiel le plus faible. Thermodynamiquement, c'est \ce{H2O} ($E^{\circ}=\SI{+1,23}{V} < \SI{+1,36}{V}$). Cependant, l'oxydation de l'eau en \ce{O2} souffre d'une forte \cle{surtension} sur le graphite (voir section suivante). En solution concentrée en \ce{Cl-}, c'est donc l'ion chlorure qui s'oxyde préférentiellement :
\[ \ce{2 Cl- -> Cl2 + 2 e-} \]

\textbf{Étape 5 — Bilan :}
\[ \ce{2 Cl- + 2 H2O ->[\text{électrolyse}] Cl2 + H2 + 2 OH-} \]

\begin{attention}[Le piège classique : le sodium métallique]
Beaucoup d'élèves écrivent $\ce{Na+ + e- -> Na}$ à la cathode. C'est \textbf{faux en solution aqueuse} ! Le couple $\ce{Na+}/\ce{Na}$ a un potentiel si bas ($\SI{-2,71}{V}$) que l'eau, oxydant bien plus fort, est réduite à sa place : on obtient \ce{H2}, jamais \ce{Na}. Le sodium métallique ne s'obtient que par électrolyse de \ce{NaCl} \emph{fondu} (sans eau), vers \SI{600}{\celsius}. Retiens : \cle{en solution aqueuse, les cations des métaux très réducteurs} (\ce{Na+}, \ce{K+}, \ce{Ca^2+}, \ce{Mg^2+}, \ce{Al^3+}) \cle{ne sont jamais réduits} ; c'est l'eau qui l'est.
\end{attention}

\courssub{Application 2 : électrolyse d'une solution d'hydroxyde de sodium}

Solution de soude \ce{NaOH}, électrodes inertes.

\textbf{Espèces présentes :} \ce{Na+}, \ce{HO-}, \ce{H2O}.

\textbf{Cathode :} compétition entre \ce{Na+} ($E^{\circ}=\SI{-2,71}{V}$) et \ce{H2O} ($E^{\circ}=\SI{-0,83}{V}$). L'eau gagne encore :
\[ \ce{2 H2O + 2 e- -> H2 + 2 HO-} \]

\textbf{Anode :} compétition entre \ce{HO-} et \ce{H2O}. Dans les deux cas, le produit d'oxydation est le même : le dioxygène. On écrit en milieu basique :
\[ \ce{4 HO- -> O2 + 2 H2O + 4 e-} \]

\textbf{Bilan} (en multipliant la demi-équation cathodique par 2 et en simplifiant) :
\[ \ce{2 H2O ->[\text{électrolyse}] 2 H2 + O2} \]

C'est tout simplement la \textbf{décomposition de l'eau} ! Les ions \ce{Na+} et \ce{HO-} ne participent pas aux réactions : ils servent uniquement à rendre la solution conductrice (l'eau pure conduit très mal le courant). On les appelle des \cle{ions spectateurs} ou \cle{électrolyte de support}. Remarque au passage la stœchiométrie remarquable : le volume de \ce{H2} recueilli est le \textbf{double} de celui de \ce{O2}, conformément à l'expérience du voltamètre de Hoffmann.

\courssub{Application 3 : électrolyse d'une solution de sulfate de cuivre}

Solution bleue de \ce{CuSO4}, électrodes inertes (graphite ou platine).

\textbf{Espèces présentes :} \ce{Cu^2+}, \ce{SO4^2-}, \ce{H2O}.

\textbf{Cathode :} compétition entre \ce{Cu^2+} ($E^{\circ} = \SI{+0,34}{V}$) et \ce{H2O} ($E^{\circ} = \SI{-0,83}{V}$). Cette fois, c'est l'ion cuivre, oxydant bien plus fort, qui est réduit :
\[ \ce{Cu^2+ + 2 e- -> Cu} \]
On observe un dépôt rougeâtre de cuivre métallique sur la cathode, et la solution se décolore progressivement.

\textbf{Anode :} l'ion sulfate \ce{SO4^2-} n'est pratiquement pas oxydable dans ces conditions ; c'est l'eau qui s'oxyde :
\[ \ce{2 H2O -> O2 + 4 H+ + 4 e-} \]

\textbf{Bilan :}
\[ \ce{2 Cu^2+ + 2 H2O ->[\text{électrolyse}] 2 Cu + O2 + 4 H+} \]

La solution devient donc de plus en plus \textbf{acide} au cours de l'électrolyse : c'est un indice expérimental précieux (un papier pH le confirme).

\courssub{Tableau récapitulatif des électrolyses au programme}

\begin{center}
\begin{tabularx}{\linewidth}{|l|X|X|X|}
\hline
\rowcolor{popblueL} \textbf{Solution} & \textbf{Cathode ($-$)} & \textbf{Anode ($+$)} & \textbf{Bilan} \\
\hline
\ce{HCl} aqueux & $\ce{2 H+ + 2 e- -> H2}$ & $\ce{2 Cl- -> Cl2 + 2 e-}$ & $\ce{2 H+ + 2 Cl- -> H2 + Cl2}$ \\
\hline
\ce{NaCl} aqueux & $\ce{2 H2O + 2 e- -> H2 + 2 HO-}$ & $\ce{2 Cl- -> Cl2 + 2 e-}$ & $\ce{2 Cl- + 2 H2O -> H2 + Cl2 + 2 HO-}$ \\
\hline
\ce{NaOH} aqueux & $\ce{2 H2O + 2 e- -> H2 + 2 HO-}$ & $\ce{4 HO- -> O2 + 2 H2O + 4 e-}$ & $\ce{2 H2O -> 2 H2 + O2}$ \\
\hline
\ce{CuSO4} aqueux & $\ce{Cu^2+ + 2 e- -> Cu}$ & $\ce{2 H2O -> O2 + 4 H+ + 4 e-}$ & $\ce{2 Cu^2+ + 2 H2O -> 2 Cu + O2 + 4 H+}$ \\
\hline
\end{tabularx}
\end{center}

\begin{aretenir}[Les trois réflexes de la prévision]
\begin{enumerate}
\item Toujours inclure \ce{H2O} dans la liste des espèces candidates, à \emph{chaque} électrode.
\item Cathode : l'oxydant du couple de $E^{\circ}$ le plus élevé est réduit. Anode : le réducteur du couple de $E^{\circ}$ le plus faible est oxydé.
\item Les cations des métaux alcalins et alcalino-terreux (\ce{Na+}, \ce{K+}, \ce{Ca^2+}...) ne sont jamais réduits en solution aqueuse : c'est \ce{H2} qui se dégage.
\end{enumerate}
\end{aretenir}

\courssub{Tension minimale d'électrolyse et limites des règles}

Une dernière question se pose : à partir de quelle tension le générateur déclenche-t-il l'électrolyse ? Si l'on augmente progressivement la tension $U$ aux bornes de la cuve, on constate expérimentalement qu'aucun courant ne passe tant que $U$ reste inférieure à une valeur seuil.

\begin{definition}[Tension minimale d'électrolyse]
La \cle{tension minimale d'électrolyse} (ou tension de seuil) $U_{\text{min}}$ est la plus petite tension qu'il faut appliquer entre les électrodes pour que l'électrolyse démarre de façon observable. En première approximation thermodynamique :
\[ U_{\text{min}} \approx E^{\circ}_{\text{anode}} - E^{\circ}_{\text{cathode}} \]
différence entre les potentiels standards des couples retenus à chaque électrode.
\end{definition}

Par exemple, pour l'électrolyse de l'eau (soude diluée) : $U_{\text{min}} \approx 1{,}23 - (-0{,}83) \approx \SI{2,06}{V}$ en théorie... mais expérimentalement, il faut souvent \SI{2,5}{V} à \SI{3}{V} ! D'où vient cet écart ?

\begin{attention}[Limites des règles de prévision]
Les règles basées sur les potentiels standards sont thermodynamiques : elles prédisent ce qui \emph{devrait} se passer à l'équilibre. Or l'électrolyse est une réaction \emph{cinétique}, forcée par le courant. Certaines réactions, en particulier les dégagements gazeux (\ce{H2}, \ce{O2}), sont « lentes à démarrer » sur certaines électrodes : il faut une tension supplémentaire appelée \cle{surtension}. Cette surtension peut \textbf{inverser les prévisions thermodynamiques} : c'est elle qui explique que \ce{Cl2} (et non \ce{O2}) se dégage à l'anode dans l'électrolyse de \ce{NaCl} concentré, alors que $E^{\circ}(\ce{O2}/\ce{H2O}) < E^{\circ}(\ce{Cl2}/\ce{Cl-})$. Nous étudierons ce phénomène en détail dans la section suivante.
\end{attention}

\begin{exoresolu}[Prévoir l'électrolyse d'une solution de bromure de potassium]
On électrolyse une solution aqueuse de bromure de potassium \ce{KBr} entre deux électrodes de graphite. On donne : $E^{\circ}(\ce{K+}/\ce{K}) = \SI{-2,92}{V}$ ; $E^{\circ}(\ce{Br2}/\ce{Br-}) = \SI{+1,07}{V}$ ; $E^{\circ}(\ce{O2}/\ce{H2O}) = \SI{+1,23}{V}$ ; $E^{\circ}(\ce{H2O}/\ce{H2}) = \SI{-0,83}{V}$.
\begin{enumerate}
\item Quelles réactions observe-t-on à chaque électrode ?
\item Écrire le bilan de l'électrolyse.
\item Estimer la tension minimale théorique.
\end{enumerate}
\tcblower
\textbf{1. Espèces présentes :} \ce{K+}, \ce{Br-}, \ce{H2O} (graphite inerte).

\textbf{Cathode (réduction) :} candidats \ce{K+} ($\SI{-2,92}{V}$) et \ce{H2O} ($\SI{-0,83}{V}$). L'oxydant le plus fort est \ce{H2O} (potentiel le plus élevé) :
\[ \ce{2 H2O + 2 e- -> H2 + 2 HO-} \]
Le potassium métallique ne se forme jamais en solution aqueuse.

\textbf{Anode (oxydation) :} candidats \ce{Br-} (couple $\ce{Br2}/\ce{Br-}$, $E^{\circ}=\SI{+1,07}{V}$) et \ce{H2O} (couple $\ce{O2}/\ce{H2O}$, $E^{\circ}=\SI{+1,23}{V}$). Le réducteur le plus fort est \ce{Br-} (potentiel le plus faible) :
\[ \ce{2 Br- -> Br2 + 2 e-} \]
On observe une coloration brune-jaune due au dibrome dissous.

\textbf{2. Bilan :}
\[ \ce{2 Br- + 2 H2O ->[\text{électrolyse}] Br2 + H2 + 2 HO-} \]

\textbf{3. Tension minimale théorique :}
\[ U_{\text{min}} \approx 1{,}07 - (-0{,}83) = \SI{1,90}{V} \]
En pratique, la surtension du dégagement de \ce{H2} impose une tension un peu supérieure.
\end{exoresolu}

\begin{savaistu}[De l'électrolyse du sel à l'eau de Javel]
L'électrolyse d'une solution concentrée de chlorure de sodium (la \emph{saumure}) est l'un des plus grands procédés de l'industrie chimique mondiale : le procédé \emph{chlore-soude}. Elle fournit simultanément trois produits stratégiques : le \textbf{dichlore} \ce{Cl2} (désinfectant, fabrication du PVC), la \textbf{soude} \ce{NaOH} (savonneries, papeteries, raffinage de l'alumine) et le \textbf{dihydrogène} \ce{H2}. Mieux : si l'on laisse le dichlore formé à l'anode rencontrer la soude formée à la cathode, dans une cuve sans séparateur, ils réagissent selon :
\[ \ce{Cl2 + 2 HO- -> ClO- + Cl- + H2O} \]
L'ion hypochlorite \ce{ClO-} obtenu est le principe actif de l'\textbf{eau de Javel}, ce désinfectant que l'on retrouve dans tous les marchés de Douala et de Yaoundé pour l'hygiène domestique et la désinfection de l'eau de boisson. Une simple cuve, du sel, de l'eau et du courant électrique : toute la chimie de cette leçon au service de la santé publique !
\end{savaistu}

\begin{exercice}[À toi de jouer]
On réalise l'électrolyse d'une solution aqueuse de nitrate d'argent \ce{AgNO3} entre électrodes inertes. On donne $E^{\circ}(\ce{Ag+}/\ce{Ag}) = \SI{+0,80}{V}$.
\begin{enumerate}
\item Faire la liste des espèces candidates à chaque électrode.
\item Prévoir les réactions aux électrodes et écrire le bilan.
\item Expliquer pourquoi la masse de la cathode augmente au cours de l'expérience.
\end{enumerate}
\end{exercice}

\begin{corrige}[Exercice]
\textbf{1.} Espèces présentes : \ce{Ag+}, \ce{NO3-}, \ce{H2O}. À la cathode : \ce{Ag+} et \ce{H2O}. À l'anode : \ce{NO3-} (non oxydable dans ces conditions) et \ce{H2O}.

\textbf{2.} Cathode : $E^{\circ}(\ce{Ag+}/\ce{Ag}) = \SI{+0,80}{V} \gg E^{\circ}(\ce{H2O}/\ce{H2}) = \SI{-0,83}{V}$, donc $\ce{Ag+ + e- -> Ag}$. Anode : $\ce{2 H2O -> O2 + 4 H+ + 4 e-}$. Bilan : $\ce{4 Ag+ + 2 H2O -> 4 Ag + O2 + 4 H+}$.

\textbf{3.} Les ions \ce{Ag+} sont réduits en argent métallique qui se \emph{dépose} sur la cathode : sa masse augmente d'autant que le métal reste accroché à l'électrode (c'est le principe de l'argenture par électrolyse).
\end{corrige}

\coursec{Le phénomène de surtension}

Dans la section précédente, nous avons appris à prévoir les réactions d'une électrolyse en comparant les potentiels standards des couples en présence : à la cathode, c'est l'oxydant le plus fort qui est réduit ; à l'anode, c'est le réducteur le plus fort qui est oxydé. Cette règle est puissante, mais elle repose sur une hypothèse implicite : que les réactions aux électrodes se produisent dès que la tension appliquée atteint la valeur thermodynamique prévue. L'expérience montre que la réalité est plus riche — et heureusement pour l'industrie ! Certaines électrolyses ne démarrent qu'à une tension nettement supérieure à la tension théorique. C'est ce phénomène, appelé \cle{surtension}, que nous allons maintenant étudier en détail.

\courssub{Constat expérimental : une électrolyse qui « refuse » de démarrer}

\begin{experience}[Une tension de seuil plus élevée que prévu]
\textbf{Dispositif.} On réalise le montage classique d'électrolyse : un générateur de tension continue réglable, un ampèremètre en série, et un électrolyseur contenant une solution aqueuse de sulfate de zinc \ce{ZnSO4} acidifiée (pH $\approx$ 5), avec une cathode de zinc (ou d'acier) et une anode inerte (platinée, graphite ou \ce{PbO2}).

\textbf{Protocole.} On augmente progressivement la tension $U$ aux bornes de l'électrolyseur, en relevant l'intensité $I$ du courant qui traverse le circuit.

\textbf{Observations.}
\begin{itemize}
\item Tant que $U$ reste inférieure à une certaine valeur seuil $U_s$, l'intensité est pratiquement nulle : aucune réaction visible aux électrodes.
\item Au-delà de $U_s$, l'intensité croît fortement : on observe un dépôt métallique gris de zinc sur la cathode et un dégagement gazeux de dioxygène sur l'anode.
\item Le point capital : la tension de seuil expérimentale $U_s$ est \textbf{supérieure} à la tension théorique calculée à partir des potentiels d'équilibre des couples mis en jeu.
\end{itemize}

\textbf{Interprétation.} Le calcul thermodynamique prévoit une tension minimale $U_{\text{th}} = |E_{\text{éq}}(\text{anode}) - E_{\text{éq}}(\text{cathode})|$ (corrigée des concentrations via la formule de Nernst). Or l'électrolyse ne démarre réellement qu'à $U_s > U_{\text{th}}$. L'écart constaté porte le nom de \cle{surtension}.
\end{experience}

Ce constat est général : presque toutes les électrolyses exigent une tension supérieure à la tension thermodynamique prévue. La courbe intensité--tension ci-dessous illustre cette situation.

\begin{popfigure}
\begin{center}
\begin{tikzpicture}
\begin{axis}[
    width=0.9\linewidth,
    height=7cm,
    axis lines=left,
    xlabel={$U$ (V)},
    ylabel={$I$ (A)},
    xmin=0, xmax=5,
    ymin=0, ymax=3.2,
    xtick={1.2,2.6},
    xticklabels={$U_{\text{th}}$,$U_s$},
    ytick=\empty,
    clip=false,
    xlabel style={anchor=west},
    ylabel style={anchor=south}
]
% Courbe I = f(U) : quasi nulle avant Us, puis croissance forte
\addplot[popcurve,domain=0:2.6,samples=60]{0.04*x};
\addplot[popcurve,domain=2.6:4.8,samples=60]{0.104 + 0.55*(x-2.6)^1.4};
% Lignes pointillées vers l'axe des abscisses
\addplot[dashed,popmuted] coordinates {(1.2,0) (1.2,0.048)};
\addplot[dashed,popmuted] coordinates {(2.6,0) (2.6,0.104)};
% Annotation écart (placée près de l'axe x)
\draw[<->,thick,poporange] (axis cs:1.2,0.15) -- (axis cs:2.6,0.15)
    node[midway,above,popdark]{\small surtension $U_s - U_{\text{th}}$};
\node[popdark,align=center] at (axis cs:3.9,1.5) {\small régime\\ \small d'électrolyse};
\end{axis}
\end{tikzpicture}
\end{center}
\legende{Courbe intensité--tension $I = f(U)$ d'un électrolyseur. Le courant ne devient significatif qu'au-delà de la tension de seuil expérimentale $U_s$, supérieure à la tension thermodynamique $U_{\text{th}}$ : l'écart $U_s - U_{\text{th}}$ traduit la somme des surtensions aux deux électrodes (et de la chute ohmique).}
\end{popfigure}

\courssub{Définition rigoureuse de la surtension}

Pour comprendre la surtension, il faut raisonner électrode par électrode, et non plus sur l'électrolyseur entier. À chaque électrode est associé un \cle{potentiel d'équilibre} (ou potentiel thermodynamique) du couple qui y réagit, noté $E_{\text{éq}}$, calculable par la formule de Nernst. Mais lorsque le courant circule, le potentiel réel de l'électrode, noté $E_{\text{réel}}$, s'écarte de cette valeur d'équilibre.

\begin{definition}[Surtension d'une électrode]
La \cle{surtension} $\eta$ (lettre grecque « êta ») d'une électrode est l'écart algébrique entre le potentiel réel de fonctionnement de l'électrode et le potentiel d'équilibre (thermodynamique) du couple mis en jeu :
\[ \eta = E_{\text{réel}} - E_{\text{éq}} \]
\begin{itemize}
\item À l'\textbf{anode} (oxydation), il faut un potentiel plus élevé que prévu : la surtension anodique est positive, $\eta_a > 0$.
\item À la \textbf{cathode} (réduction), il faut un potentiel plus bas (plus négatif) que prévu : la surtension cathodique est négative, $\eta_c < 0$.
\end{itemize}
Dans les deux cas, la surtension \emph{augmente} la tension à appliquer pour que l'électrolyse se produise : $U_{\text{min}} = U_{\text{th}} + |\eta_a| + |\eta_c| + r\,I$ (le terme $r\,I$ étant la chute ohmique dans l'électrolyte).
\end{definition}

\begin{popfigure}
\begin{center}
\begin{tikzpicture}[scale=1.1]
% Axe vertical des potentiels
\draw[->,thick,popdark] (0,0) -- (0,6.5) node[above]{\small Potentiel $E$ (V)};
% Niveaux cathode
\draw[thick,popblue] (1,1.5) -- (4,1.5) node[right,popblue]{\small $E_{\text{éq}}$ (cathode)};
\draw[thick,popblue!60] (1,0.5) -- (4,0.5) node[right,popblue!80]{\small $E_{\text{réel}}$ (cathode)};
\draw[<->,thick,poporange] (3.7,0.5) -- (3.7,1.5) node[midway,right,popdark]{\small $|\eta_c|$};
\node[popdark,align=left] at (4.5,1.0) {$\eta_c < 0$};
% Niveaux anode
\draw[thick,poppink] (1,4.5) -- (4,4.5) node[right,poppink]{\small $E_{\text{éq}}$ (anode)};
\draw[thick,poppink!60] (1,5.5) -- (4,5.5) node[right,poppink!80]{\small $E_{\text{réel}}$ (anode)};
\draw[<->,thick,poporange] (3.7,4.5) -- (3.7,5.5) node[midway,right,popdark]{\small $\eta_a$};
\node[popdark,align=left] at (4.5,5.0) {$\eta_a > 0$};
% Tension théorique vs réelle
\draw[<->,thick,popgreen] (0.5,1.5) -- (0.5,4.5) node[midway,left,popdark,align=center]{\small $U_{\text{th}}$};
\draw[<->,thick,popgreen!60!black] (0.15,0.5) -- (0.15,5.5) node[midway,left,popdark,align=center]{\small $U_{\text{min}}$};
\end{tikzpicture}
\end{center}
\legende{Comparaison des potentiels théoriques (d'équilibre) et réels de fonctionnement des deux électrodes. La surtension cathodique $\eta_c$ abaisse le potentiel de la cathode, la surtension anodique $\eta_a$ élève celui de l'anode : la tension minimale réelle $U_{\text{min}}$ dépasse la tension thermodynamique $U_{\text{th}}$.}
\end{popfigure}

\begin{attention}[La surtension n'est pas une perte ohmique !]
Erreur très fréquente au Probatoire : confondre la surtension avec la chute de tension due à la résistance de l'électrolyte (l'effet Joule, $r\,I$). Ce sont deux phénomènes de nature différente :
\begin{itemize}
\item La \textbf{chute ohmique} $r\,I$ est une perte d'énergie par échauffement de la solution ; elle est proportionnelle à l'intensité et disparaît quand $I = 0$.
\item La \textbf{surtension} est liée à la \emph{cinétique} des réactions aux électrodes : certaines réactions électrochimiques sont lentes, « difficiles à amorcer » à la surface du métal, et exigent un « coup de pouce » supplémentaire sous forme de potentiel plus énergique. Elle existe dès que le courant commence à circuler, même faiblement.
\end{itemize}
Retiens : la surtension est un phénomène de \cle{cinétique électrochimique}, pas de résistance électrique.
\end{attention}

\courssub{De quoi dépend la surtension ?}

La surtension n'est pas une constante universelle : c'est une grandeur expérimentale qui dépend de plusieurs facteurs.

\begin{propriete}[Facteurs influençant la surtension]
La valeur de la surtension $\eta$ dépend :
\begin{enumerate}
\item de la \textbf{nature du métal de l'électrode} : pour une même réaction, la surtension varie énormément d'un métal à l'autre (platine, mercure, zinc, plomb, cuivre...) ;
\item de la \textbf{nature de l'espèce qui se dépose ou se dégage} : le dégagement d'un gaz (comme \ce{H2} ou \ce{O2}) présente généralement une surtension bien plus forte que le dépôt d'un métal ;
\item de la \textbf{densité de courant} $j = I/S$ (intensité par unité de surface de l'électrode, en \si{A.m^{-2}} ou \si{A.cm^{-2}}) : plus on veut forcer la réaction à aller vite, plus la surtension augmente (loi de Tafel) ;
\item accessoirement, de l'état de surface de l'électrode (polie, rugueuse, oxydée) et de la température.
\end{enumerate}
\end{propriete}

Le cas le plus célèbre et le plus important pour nous est celui du \textbf{dégagement de dihydrogène} (réduction des ions \ce{H+} ou de l'eau) selon le métal de la cathode :

\begin{center}
\begin{tabularx}{\linewidth}{|l|X|}
\hline
\rowcolor{popblueL} \textbf{Métal de la cathode} & \textbf{Surtension de dégagement de \ce{H2} (ordre de grandeur à $j \approx \SI{1}{A.dm^{-2}}$)} \\
\hline
Platine (platiné) & très faible, $\approx \SI{0,05}{V}$ \\
\hline
Cuivre, fer, nickel & moyenne, $\approx \SI{0,4}{V}$ à \SI{0,6}{V} \\
\hline
Zinc & forte, $\approx \SI{0,7}{V}$ à \SI{0,8}{V} \\
\hline
Plomb & forte, $\approx \SI{0,8}{V}$ à \SI{1,0}{V} \\
\hline
Mercure & très forte, $\approx \SI{1,0}{V}$ à \SI{1,2}{V} \\
\hline
\end{tabularx}
\end{center}

Le dégagement de dioxygène à l'anode présente aussi une surtension notable (environ \SI{0,4}{V} à \SI{0,6}{V} sur platine, davantage sur d'autres métaux). En revanche, le \textbf{dépôt d'un métal} (réduction d'un cation métallique comme \ce{Cu^2+} ou \ce{Zn^2+}) se fait presque toujours avec une surtension négligeable : le métal se dépose dès que le potentiel thermodynamique est atteint.

\begin{exemplebox}[La surtension de \ce{H2} sur le mercure]
Pour le couple \ce{H+}/\ce{H2} sur une électrode de \textbf{mercure}, la surtension cathodique peut atteindre environ \SI{1}{V} (voire \SI{1,2}{V} à forte densité de courant). Concrètement, alors que la thermodynamique prévoit le dégagement de dihydrogène dès que le potentiel de l'électrode atteint environ \SI{0}{V} (à pH = 0, couple \ce{H+}/\ce{H2}), il faut en réalité descendre jusqu'à environ \SI{-1}{V} pour observer les premières bulles de \ce{H2} sur le mercure. Pendant tout cet intervalle, la réduction de \ce{H+} est « bloquée » cinétiquement : d'autres réductions, thermodynamiquement moins favorables, peuvent alors avoir lieu à sa place !
\end{exemplebox}

\courssub{Conséquence majeure : la surtension peut inverser les prévisions}

Voici le point le plus subtil — et le plus riche — de cette leçon. La règle de prévision vue à la section précédente repose sur les potentiels \emph{d'équilibre} (thermodynamique). Mais l'électrolyse réelle est gouvernée par les potentiels \emph{réels}, qui incluent les surtensions (cinétique). Il peut donc arriver que \textbf{l'ordre prévu par les potentiels d'équilibre soit inversé} : une espèce thermodynamiquement moins facile à réduire se réduit en pratique, parce que sa concurrente « souffre » d'une forte surtension.

\begin{exoresolu}[Le dépôt de zinc en solution aqueuse : une « impossibilité » thermodynamique qui marche]
\textbf{Énoncé.} On veut électrolyser une solution aqueuse de sulfate de zinc \ce{ZnSO4} (acidifiée, pH voisin de 5) avec une cathode de zinc. On donne : $E^{\circ}(\ce{Zn^2+}/\ce{Zn}) = \SI{-0,76}{V}$ ; pour le couple de l'eau à pH = 5, $E_{\text{éq}}(\ce{H2O}/\ce{H2}) \approx \SI{-0,30}{V}$ ; la surtension de dégagement de \ce{H2} sur le zinc vaut environ $\eta_c \approx \SI{-0,75}{V}$ ; la surtension de dépôt du zinc est négligeable.

Prévoir, par la seule thermodynamique, ce qui devrait se passer à la cathode. Montrer ensuite que la surtension renverse la conclusion.
\tcblower
\textbf{Solution détaillée.}

\textbf{Étape 1 : prévision thermodynamique.} À la cathode, deux réductions sont envisageables :
\begin{itemize}
\item réduction des ions zinc : $\ce{Zn^2+ + 2e- -> Zn}$, de potentiel d'équilibre $E_{\text{éq}} \approx \SI{-0,76}{V}$ (conditions standards) ;
\item réduction de l'eau : $\ce{2H2O + 2e- -> H2 ^ + 2HO-}$, de potentiel d'équilibre $E_{\text{éq}} \approx \SI{-0,30}{V}$ à pH = 5.
\end{itemize}
La règle de prévision dit que l'oxydant le plus fort (potentiel d'équilibre le plus élevé) est réduit : c'est l'eau ($\SI{-0,30}{V} > \SI{-0,76}{V}$). La thermodynamique prévoit donc un \textbf{dégagement de dihydrogène} et \emph{pas} de dépôt de zinc.

\textbf{Étape 2 : prise en compte des surtensions.} Les potentiels réels de fonctionnement sont :
\begin{align*}
E_{\text{réel}}(\ce{H2}) &= E_{\text{éq}}(\ce{H2O}/\ce{H2}) + \eta_c = -0,30 + (-0,75) = \SI{-1,05}{V} \\
E_{\text{réel}}(\ce{Zn}) &= E_{\text{éq}}(\ce{Zn^2+}/\ce{Zn}) + 0 \approx \SI{-0,76}{V}
\end{align*}

\textbf{Étape 3 : nouvelle comparaison.} En réalité, c'est le couple \ce{Zn^2+}/\ce{Zn} qui présente le potentiel réel le plus élevé ($\SI{-0,76}{V} > \SI{-1,05}{V}$) : \textbf{c'est donc le zinc qui se dépose} à la cathode, et non le dihydrogène qui se dégage !

\textbf{Conclusion.} Grâce à la forte surtension de dégagement de \ce{H2} sur le zinc, le dépôt électrolytique de zinc est possible en solution aqueuse, alors que la thermodynamique seule le jugeait impossible. C'est le principe du \textbf{zincage électrolytique} (galvanisation), utilisé pour protéger l'acier de la corrosion.
\end{exoresolu}

\begin{methode}[Prévoir une électrolyse en tenant compte des surtensions]
\begin{enumerate}
\item Faire l'inventaire des espèces présentes au voisinage de chaque électrode (soluté, solvant eau, électrode elle-même si elle est attaquable).
\item Écrire les demi-équations de réduction possibles à la cathode et d'oxydation possibles à l'anode, avec les potentiels d'équilibre (formule de Nernst si nécessaire).
\item Corriger chaque potentiel d'équilibre par la surtension correspondante : $E_{\text{réel}} = E_{\text{éq}} + \eta$ (avec $\eta_c < 0$ à la cathode, $\eta_a > 0$ à l'anode). Les surtensions sont des données expérimentales : elles te seront fournies dans l'énoncé.
\item Comparer les potentiels \emph{réels} : à la cathode, c'est le couple de potentiel réel le plus élevé qui réagit ; à l'anode, celui de potentiel réel le plus bas.
\item Écrire l'équation bilan de l'électrolyse et conclure sur les produits obtenus.
\end{enumerate}
\end{methode}

\courssub{Intérêt industriel de la surtension}

Loin d'être un simple « défaut » qui gaspille de l'énergie, la surtension est une aubaine pour l'industrie : elle rend possibles des électrolyses que la thermodynamique interdisait.

\begin{itemize}
\item \textbf{Électrométallurgie du zinc.} Comme nous venons de le voir, la forte surtension de \ce{H2} sur le zinc permet d'extraire le zinc de ses solutions aqueuses par électrolyse, avec un excellent rendement. Une grande partie du zinc mondial est produite ainsi (procédé hydrométallurgique : lixiviation $\rightarrow$ purification $\rightarrow$ électrolyse).
\item \textbf{Galvanoplastie.} Les dépôts électrolytiques de métaux (chromage des pièces de moto, nickelage, dorure des bijoux fantaisie, cuivrage des circuits imprimés) reposent sur le même principe : le dépôt métallique, à surtension quasi nulle, « gagne » contre le dégagement de \ce{H2}, ralenti par sa surtension.
\item \textbf{Accumulateur au plomb} (la batterie des voitures et motos de nos villes) : la forte surtension de \ce{H2} et de \ce{O2} sur le plomb (électrodes en alliage Pb-Sb ou Pb-Ca) empêche la décomposition parasite de l'eau de l'électrolyte pendant la charge, ce qui rend la batterie rechargeable avec un bon rendement.
\end{itemize}

\begin{savaistu}[La soude très pure grâce au mercure : l'amalgame change la donne]
Pendant près d'un siècle, l'industrie a produit la soude \ce{NaOH} et le dichlore \ce{Cl2} par électrolyse de la saumure (solution concentrée de \ce{NaCl}) avec une \textbf{cathode de mercure} (procédé « à cathode de mercure » ou Castner-Kellner). Le secret du procédé ne repose pas \emph{uniquement} sur la surtension !

Rappelons que $E^{\circ}(\ce{Na+}/\ce{Na}) = \SI{-2,71}{V}$. Même avec une surtension de \ce{H2} de \SI{-1}{V} sur Hg ($E_{\text{réel}}(\ce{H2}) \approx \SI{-1}{V}$ à pH 0), le sodium métallique pur resterait thermodynamiquement inaccessible ($-2,71 < -1$). L'astuce géniale est que le sodium formé se dissout immédiatement dans le mercure pour former un \textbf{amalgame} \ce{Na(Hg)}. L'activité du sodium dans cet amalgame est très faible, ce qui, d'après l'équation de Nernst, \textbf{décale fortement le potentiel du couple \ce{Na+}/\ce{Na(Hg)} vers des valeurs moins négatives} (environ \SI{-1,0}{V} à \SI{-1,2}{V} vs ECS).

Ainsi, sur mercure :
\begin{itemize}
\item $E_{\text{réel}}(\ce{H2}) \approx \SI{-1,0}{V}$ (surtension forte) ;
\item $E_{\text{éq}}(\ce{Na+}/\ce{Na(Hg)}) \approx \SI{-1,0}{V}$ (potentiel décalé par l'amalgame).
\end{itemize}
Les deux potentiels réels deviennent comparables, et la réduction du \ce{Na+} a lieu ! L'amalgame \ce{Na(Hg)} est ensuite décomposé par l'eau dans un compartiment séparé (le « décomposeur ») pour donner une soude d'une pureté remarquable et régénérer le mercure :
\[ \ce{2 Na(Hg) + 2 H2O -> 2 NaOH + H2 ^ + 2 Hg} \]
Ce procédé a été progressivement abandonné depuis les années 1980 à cause de la toxicité du mercure pour l'environnement (catastrophe de Minamata au Japon), au profit des électrolyses à membrane (sans mercure). Mais il reste l'exemple historique le plus spectaculaire du pouvoir combiné de la \textbf{surtension} et de la \textbf{thermodynamique des solutions} (amalgame) pour faire réagir un couple « impossible » !
\end{savaistu}

\begin{attention}[Ce qu'il faut retenir sans se tromper]
\begin{itemize}
\item La surtension \textbf{dépend du couple électrode--réaction} : on ne peut pas parler de « la » surtension d'un métal sans préciser quelle réaction s'y déroule. Le mercure présente une forte surtension pour le dégagement de \ce{H2}, mais pas nécessairement pour d'autres réactions.
\item Une surtension élevée \textbf{ralentit} la réaction concernée ; elle ne la rend pas rigoureusement impossible : à potentiel suffisamment énergique, le dégagement de \ce{H2} finit toujours par se produire, même sur le mercure.
\item Ne conclus jamais une prévision d'électrolyse sans vérifier si l'énoncé fournit des données de surtension : leur présence est un signal fort qu'il faut corriger les potentiels d'équilibre avant de comparer.
\item Attention au vocabulaire : on dit « la surtension cathodique \emph{vaut} \SI{-0,8}{V} » (valeur algébrique) ou « la surtension \emph{est de} \SI{0,8}{V} » (valeur absolue). Sois précis dans tes rédactions.
\end{itemize}
\end{attention}

\begin{aretenir}[L'essentiel sur la surtension]
\begin{itemize}
\item La \cle{surtension} $\eta = E_{\text{réel}} - E_{\text{éq}}$ est l'écart entre le potentiel réel de fonctionnement d'une électrode et le potentiel thermodynamique du couple mis en jeu.
\item Elle est d'origine \textbf{cinétique} (lenteur des réactions à la surface de l'électrode), à ne pas confondre avec la chute ohmique $r\,I$.
\item Elle dépend de la nature du métal de l'électrode, de l'espèce qui se dégage ou se dépose, et de la densité de courant.
\item Le dégagement des gaz \ce{H2} et \ce{O2} présente des surtensions fortes sur certains métaux (\ce{Hg}, \ce{Pb}, \ce{Zn}) ; le dépôt des métaux a une surtension quasi nulle.
\item Conséquence : la surtension peut \textbf{inverser l'ordre prévu par les potentiels d'équilibre} et modifier les produits de l'électrolyse (exemple : dépôt du zinc en solution aqueuse).
\item Intérêt industriel : électrométallurgie du zinc, galvanoplastie, accumulateurs au plomb, ancien procédé de fabrication de la soude à cathode de mercure (amalgame).
\end{itemize}
\end{aretenir}

\begin{exercice}[Surtension et électrolyse du chlorure de sodium]
On électrolyse une solution aqueuse de chlorure de sodium \ce{NaCl} entre une cathode de plomb et une anode inerte. On donne : $E^{\circ}(\ce{Na+}/\ce{Na}) = \SI{-2,71}{V}$ ; $E_{\text{éq}}(\ce{H2O}/\ce{H2}) \approx \SI{-0,42}{V}$ à pH = 7 ; surtension de \ce{H2} sur le plomb : $\eta_c \approx \SI{-0,9}{V}$.
\begin{enumerate}
\item Sans tenir compte de la surtension, quelle réaction est prévue à la cathode ?
\item En tenant compte de la surtension, cette conclusion change-t-elle ? Justifier par un calcul.
\end{enumerate}
\end{exercice}

\begin{corrige}[Exercice 1]
\textbf{1.} À la cathode, deux réductions sont possibles : $\ce{Na+ + e- -> Na}$ ($E^{\circ} = \SI{-2,71}{V}$) et $\ce{2H2O + 2e- -> H2 ^ + 2HO-}$ ($E_{\text{éq}} \approx \SI{-0,42}{V}$). L'oxydant le plus fort est l'eau (potentiel d'équilibre le plus élevé) : la thermodynamique prévoit le dégagement de \ce{H2}.

\textbf{2.} Potentiels réels :
\begin{align*}
E_{\text{réel}}(\ce{H2}) &= E_{\text{éq}}(\ce{H2O}/\ce{H2}) + \eta_c = -0,42 + (-0,9) = \SI{-1,32}{V} \\
E_{\text{réel}}(\ce{Na}) &\approx E^{\circ}(\ce{Na+}/\ce{Na}) = \SI{-2,71}{V} \quad (\text{surtension de dépôt négligeable})
\end{align*}
Même corrigé de la surtension, le potentiel réel du couple de l'eau ($\SI{-1,32}{V}$) reste très supérieur à celui du couple \ce{Na+}/\ce{Na} ($\SI{-2,71}{V}$) : c'est toujours l'eau qui est réduite, avec dégagement de \ce{H2}. La conclusion ne change donc pas ici : la surtension de \SI{0,9}{V} ne suffit pas à combler l'écart thermodynamique de plus de \SI{2}{V}. (C'est précisément pour cela que le procédé industriel à cathode de mercure utilisait... du mercure, dont la surtension est encore plus forte \emph{et} qui stabilise le sodium sous forme d'amalgame, décalant son potentiel d'équilibre.)
\end{corrige}

\coursec{Aspects quantitatifs de l'électrolyse}

Dans la section précédente, nous avons vu que le phénomène de \cle{surtension} permet de prévoir quelles espèces réagissent réellement aux électrodes, même lorsque les potentiels thermodynamiques laissent plusieurs possibilités. Nous savons donc désormais \emph{ce qui} se produit lors d'une électrolyse. Il reste une question fondamentale, celle que se pose tout industriel, tout bijoutier, tout artisan qui chrome une pièce de moto : \emph{combien} de produit obtient-on ? Combien de grammes de cuivre se déposent-ils ? Quel volume de dichlore se dégage-t-il ? Pendant combien de temps faut-il faire passer le courant ?

C'est toute l'ambition de cette section : passer du \textbf{qualitatif} (quelle réaction ?) au \textbf{quantitatif} (quelle quantité de matière ?). La bonne nouvelle, c'est que la réponse tient en une idée simple et lumineuse : \textbf{le courant électrique est un débit d'électrons}, et chaque électron qui circule participe à la réaction. Compter les électrons, c'est compter les atomes transformés. Toute la stœchiométrie de l'électrolyse découle de cette observation.

\courssub{La quantité d'électricité : compter les électrons qui passent}

\courssub{Intuition : le courant comme un fleuve d'électrons}

Imagine le fleuve Wouri qui traverse Douala. Pour connaître la quantité d'eau qui est passée sous le pont en une journée, il suffit de connaître le \emph{débit} du fleuve (combien de mètres cubes par seconde) et de le multiplier par la durée. Le courant électrique fonctionne exactement pareil :

\begin{itemize}
\item l'\cle{intensité} $I$ du courant, mesurée en \textbf{ampères} (A), est le \emph{débit d'électrons} : un ampère correspond à une charge d'un coulomb qui passe chaque seconde ;
\item la \textbf{quantité d'électricité} $Q$ est la \emph{charge totale} qui a traversé le circuit pendant la durée de l'expérience.
\end{itemize}

\begin{definition}[Quantité d'électricité]
La \cle{quantité d'électricité} $Q$ qui traverse un électrolyseur parcouru par un courant d'intensité constante $I$ pendant une durée $t$ est :
\[ Q = I \times t \]
avec $Q$ en \textbf{coulombs} (C), $I$ en \textbf{ampères} (A) et $t$ en \textbf{secondes} (s).
\end{definition}

\begin{attention}[Le piège du temps]
L'erreur la plus fréquente au Probatoire est d'oublier de convertir la durée en \textbf{secondes}. Si l'énoncé dit « 30 minutes », il faut écrire $t = 30 \times 60 = \SI{1800}{s}$. Si l'énoncé dit « 2 heures », alors $t = 2 \times 3600 = \SI{7200}{s}$. Un temps laissé en minutes fausse tout le calcul d'un facteur 60 !
\end{attention}

\begin{exemplebox}[Calcul d'une quantité d'électricité]
Un électrolyseur est traversé par un courant d'intensité $I = \SI{2,0}{A}$ pendant $30$ minutes. La quantité d'électricité mise en jeu est :
\[ Q = I \times t = 2{,}0 \times (30 \times 60) = 2{,}0 \times 1800 = \SI{3600}{C}. \]
\end{exemplebox}

\courssub{Le faraday : le « paquet » d'une mole d'électrons}

Les électrons sont des grains de charge minuscules : la charge d'un seul électron vaut, en valeur absolue, $e = 1{,}602 \times 10^{-19}\ \mathrm{C}$. Compter les électrons un par un serait aussi absurde que compter les grains de sable d'une plage un par un. Les chimistes raisonnent donc en \textbf{moles}, comme toujours.

\begin{definition}[Le faraday]
Le \cle{faraday} (symbole $\mathcal{F}$) est la valeur absolue de la charge électrique portée par \textbf{une mole d'électrons} :
\[ \mathcal{F} = N_A \times e = 6{,}022 \times 10^{23} \times 1{,}602 \times 10^{-19} \approx \SI{96500}{C.mol^{-1}}. \]
On retiendra la valeur approchée : $1\ \mathcal{F} \approx \SI{96500}{C.mol^{-1}}$.
\end{definition}

\begin{savaistu}[Michael Faraday, le fils de forgeron devenu géant de la science]
Michael Faraday (1791--1867), fils d'un forgeron anglais presque autodidacte, établit vers 1833 les lois quantitatives de l'électrolyse qui portent son nom : la masse de substance transformée à une électrode est proportionnelle à la quantité d'électricité qui traverse le circuit. C'est en son honneur que l'unité $\SI{96500}{C.mol^{-1}}$ s'appelle le \emph{faraday}. Sans ses lois, pas d'industrie de l'aluminium, pas de chrome sur nos motos, pas d'eau de Javel industrielle !
\end{savaistu}

La conséquence immédiate est la relation centrale de toute cette section : connaissant la charge totale $Q$, on en déduit la \textbf{quantité de matière d'électrons} qui ont été échangés aux électrodes :

\begin{aretenir}[Quantité de matière d'électrons échangés]
\[ n(e^-) = \frac{Q}{\mathcal{F}} = \frac{I \times t}{\mathcal{F}} \]
avec $n(e^-)$ en moles, $Q$ en coulombs et $\mathcal{F} \approx \SI{96500}{C.mol^{-1}}$.
\end{aretenir}

\begin{exemplebox}[De la charge aux moles d'électrons]
Reprenons $Q = \SI{3600}{C}$. Le nombre de moles d'électrons échangés est :
\[ n(e^-) = \frac{3600}{96500} \approx \SI{3,73 \times 10^{-2}}{mol} \approx 0{,}037\ \mathrm{mol}. \]
\end{exemplebox}

\courssub{Des électrons aux produits : la stœchiométrie des demi-équations}

Les électrons ne se promènent pas seuls : ils sont \emph{consommés} à la cathode (réduction) et \emph{fournis} à l'anode (oxydation), selon des demi-équations électroniques. C'est la demi-équation qui fixe le \textbf{lien de proportionnalité} entre les moles d'électrons et les moles de produit.

Prenons l'exemple du dépôt de cuivre à la cathode lors de l'électrolyse d'une solution de sulfate de cuivre :
\[ \ce{Cu^2+ + 2e- -> Cu} \]
Cette demi-équation nous dit que \textbf{2 moles d'électrons} sont nécessaires pour former \textbf{1 mole de cuivre}. D'où :
\[ n(\ce{Cu}) = \frac{n(e^-)}{2}. \]

De même, pour le dégagement de dichlore à l'anode : $\ce{2Cl- -> Cl2 + 2e-}$, soit $n(\ce{Cl2}) = \dfrac{n(e^-)}{2}$. Pour le dépôt d'argent : $\ce{Ag+ + e- -> Ag}$, soit $n(\ce{Ag}) = n(e^-)$.

\begin{attention}[Vérifier le nombre d'électrons de la demi-équation]
Avant tout calcul, \textbf{écris la demi-équation électronique équilibrée} et relève le nombre $n$ d'électrons qu'elle met en jeu. Confondre $\ce{Cu^2+ + 2e- -> Cu}$ (2 électrons) avec $\ce{Ag+ + e- -> Ag}$ (1 électron) double ou divise par deux le résultat final. C'est le deuxième piège classique après la conversion du temps !
\end{attention}

\courssub{Masse de produit formé et volume de gaz dégagé}

\courssub{Masse déposée ou consommée à une électrode}

Une fois la quantité de matière du produit connue, la masse s'obtient par la relation habituelle $m = n \times M$. En enchaînant toutes les étapes :

\begin{aretenir}[Masse de produit formé à une électrode]
Pour une demi-équation mettant en jeu $n$ électrons par mole de produit :
\[ m = \frac{M \times I \times t}{n \times \mathcal{F}} \]
avec $m$ en grammes, $M$ la masse molaire du produit ($\mathrm{g.mol^{-1}}$), $I$ en ampères, $t$ en secondes, $n$ le nombre d'électrons de la demi-équation et $\mathcal{F} \approx \SI{96500}{C.mol^{-1}}$.
\end{aretenir}

Cette formule traduit les \textbf{lois de Faraday} : à intensité constante, la masse déposée est \emph{proportionnelle au temps} ; à durée constante, elle est \emph{proportionnelle à l'intensité}. La courbe $m = f(t)$ est donc une droite passant par l'origine.

\begin{popfigure}
\begin{center}
\begin{tikzpicture}
\begin{axis}[
    width=0.85\linewidth, height=6.2cm,
    xlabel={Temps $t$ (min)}, ylabel={Masse de cuivre déposée $m$ (g)},
    xmin=0, xmax=65, ymin=0, ymax=2.8,
    axis lines=left,
    xtick={0,15,30,45,60}, ytick={0,0.5,1,1.19,1.5,2,2.5},
    grid=both, grid style={popline!40},
    label style={font=\small}, tick label style={font=\small},
]
\addplot[popcurve, domain=0:62, samples=100] {0.0395*x};
\addplot[mark=*, mark size=2pt, poporange] coordinates {(30,1.19)};
\draw[dashed, popmuted] (axis cs:30,0) -- (axis cs:30,1.19) -- (axis cs:0,1.19);
\node[font=\small, popdark, anchor=south west] at (axis cs:31,1.19) {$I = \SI{2}{A}$, $t = 30$ min : $m \approx \SI{1,19}{g}$};
\end{axis}
\end{tikzpicture}
\end{center}
\legende{Masse de cuivre déposée en fonction du temps lors de l'électrolyse de \ce{CuSO4} à intensité constante $I = \SI{2}{A}$ : droite passant par l'origine (loi de Faraday).}
\end{popfigure}

\courssub{Volume de gaz dégagé}

Lorsque le produit est un gaz (dihydrogène à la cathode, dichlore ou dioxygène à l'anode), on mesure un \textbf{volume} plutôt qu'une masse. On utilise alors le volume molaire $V_m$ dans les conditions de l'expérience :

\begin{aretenir}[Volume de gaz dégagé]
\[ V = n(\text{gaz}) \times V_m \]
où $V_m$ est le volume molaire des gaz dans les conditions de l'expérience (par exemple $V_m = \SI{24,0}{L.mol^{-1}}$ à \SI{20}{\celsius} sous pression normale, ou $V_m = \SI{22,4}{L.mol^{-1}}$ dans les conditions normales de température et de pression). Vérifie toujours la valeur de $V_m$ donnée par l'énoncé !
\end{aretenir}

\begin{exemplebox}[Dégagement de dihydrogène]
Lors de l'électrolyse d'une solution de soude, la cathode est le siège de la réduction de l'eau :
\[ \ce{2H2O + 2e- -> H2 + 2HO-} \]
Avec $n(e^-) = 0{,}0373\ \mathrm{mol}$, on obtient $n(\ce{H2}) = \dfrac{0{,}0373}{2} \approx 0{,}0187\ \mathrm{mol}$, soit, à \SI{20}{\celsius} ($V_m = \SI{24,0}{L.mol^{-1}}$) :
\[ V(\ce{H2}) = 0{,}0187 \times 24{,}0 \approx \SI{0,45}{L}. \]
\end{exemplebox}

\courssub{La démarche générale de résolution}

Toutes les étapes se résument en une chaîne logique unique, qu'il faut savoir dérouler les yeux fermés :

\begin{popfigure}
\begin{center}
\begin{tikzpicture}[
    node distance=0.55cm,
    etape/.style={rectangle, rounded corners=3pt, draw=popblue, fill=popblueL, thick, text width=3.1cm, minimum height=1.15cm, align=center, font=\small},
    formule/.style={rectangle, draw=poporange, fill=poporangeL, thick, text width=3.0cm, minimum height=1.15cm, align=center, font=\small},
    fleche/.style={-{Stealth[length=3mm]}, thick, popdark}
]
\node[etape] (q) {\textbf{Étape 1}\\ Quantité d'électricité $Q$};
\node[formule, right=of q] (ne) {\textbf{Étape 2}\\ Moles d'électrons $n(e^-)$};
\node[etape, right=of ne] (np) {\textbf{Étape 3}\\ Moles de produit $n$};
\node[formule, right=of np] (mv) {\textbf{Étape 4}\\ Masse $m$ ou volume $V$};
\draw[fleche] (q) -- node[above, font=\footnotesize]{$Q = I \times t$} (ne);
\draw[fleche] (ne) -- node[above, font=\footnotesize]{$n(e^-) = \dfrac{Q}{\mathcal{F}}$} (np);
\draw[fleche] (np) -- node[above, font=\footnotesize]{demi-équation} (mv);
\node[below=0.35cm of ne, font=\footnotesize\itshape, popmuted, text width=9cm, align=center] {Étape 3 : $n(\text{produit}) = \dfrac{n(e^-)}{n}$, où $n$ est le nombre d'électrons de la demi-équation. \quad Étape 4 : $m = n \times M$ \ ou \ $V = n \times V_m$.};
\end{tikzpicture}
\end{center}
\legende{Organigramme de la démarche de calcul quantitatif d'une électrolyse : de la charge électrique jusqu'à la masse ou au volume de produit.}
\end{popfigure}

\begin{methode}[Résoudre un exercice quantitatif d'électrolyse]
\begin{enumerate}
\item \textbf{Identifier les réactions aux électrodes} et écrire les demi-équations électroniques équilibrées (en appliquant les règles de prévision et en tenant compte des surtensions).
\item \textbf{Calculer la quantité d'électricité} : $Q = I \times t$, en convertissant impérativement le temps en secondes.
\item \textbf{Calculer la quantité de matière d'électrons} : $n(e^-) = \dfrac{Q}{\mathcal{F}}$.
\item \textbf{Utiliser la demi-équation} pour trouver la quantité de produit : $n(\text{produit}) = \dfrac{n(e^-)}{n}$.
\item \textbf{Conclure} : masse $m = n \times M$ ou volume de gaz $V = n \times V_m$, avec les unités et les chiffres significatifs.
\end{enumerate}
\end{methode}

\begin{exoresolu}[Électrolyse d'une solution de sulfate de cuivre]
On réalise l'électrolyse d'une solution de sulfate de cuivre \ce{CuSO4} acidifiée, entre une cathode de cuivre et une anode de graphite inerte. L'intensité du courant vaut $I = \SI{2,0}{A}$ et l'électrolyse dure $30$ minutes. On donne $M(\ce{Cu}) = \SI{63,5}{g.mol^{-1}}$ et $\mathcal{F} = \SI{96500}{C.mol^{-1}}$.
\begin{enumerate}
\item Écrire la demi-équation à la cathode.
\item Calculer la masse de cuivre déposée sur la cathode.
\end{enumerate}
\tcblower
\textbf{1. Demi-équation cathodique.} À la cathode, les ions cuivre(II), meilleurs oxydants que l'eau, sont réduits en cuivre métallique :
\[ \ce{Cu^2+ + 2e- -> Cu} \]
Cette demi-équation met en jeu $n = 2$ électrons par mole de cuivre formé.

\textbf{2. Masse de cuivre déposée.} On déroule la démarche en quatre étapes.

\emph{Étape 1 — Quantité d'électricité} (attention : convertir $30$ min en secondes) :
\[ Q = I \times t = 2{,}0 \times (30 \times 60) = \SI{3600}{C}. \]

\emph{Étape 2 — Moles d'électrons} :
\[ n(e^-) = \frac{Q}{\mathcal{F}} = \frac{3600}{96500} \approx \SI{3,73 \times 10^{-2}}{mol}. \]

\emph{Étape 3 — Moles de cuivre} (2 électrons pour 1 atome de cuivre) :
\[ n(\ce{Cu}) = \frac{n(e^-)}{2} = \frac{0{,}0373}{2} \approx \SI{1,87 \times 10^{-2}}{mol}. \]

\emph{Étape 4 — Masse de cuivre} :
\[ m(\ce{Cu}) = n(\ce{Cu}) \times M(\ce{Cu}) = 0{,}0187 \times 63{,}5 \approx \SI{1,19}{g}. \]

\textbf{Conclusion :} il se dépose environ $\boxed{\SI{1,19}{g}}$ de cuivre sur la cathode. On vérifie avec la formule directe : $m = \dfrac{M \times I \times t}{n \times \mathcal{F}} = \dfrac{63{,}5 \times 2{,}0 \times 1800}{2 \times 96500} \approx \SI{1,19}{g}$. 
\end{exoresolu}

\courssub{Bilan de matière : le cas de l'anode soluble}

Dans certaines électrolyses, l'anode n'est pas inerte : elle est constituée du \emph{même métal} que celui qu'on dépose à la cathode. C'est le cas de l'électrolyse de \ce{CuSO4} entre deux électrodes de cuivre. On observe alors :

\begin{itemize}
\item à la \textbf{cathode} : réduction $\ce{Cu^2+ + 2e- -> Cu}$ : le métal se \emph{dépose}, la cathode grossit ;
\item à l'\textbf{anode} : oxydation du métal de l'électrode $\ce{Cu -> Cu^2+ + 2e-}$ : l'anode se \emph{dissout}, elle maigrit.
\end{itemize}

Comme les deux demi-équations mettent en jeu le même nombre d'électrons, la conservation de la matière impose un résultat remarquable :

\begin{propriete}[Conservation de la masse à anode soluble]
Lors d'une électrolyse à anode soluble, la \textbf{masse de métal perdue par l'anode est exactement égale à la masse de métal déposée sur la cathode}. La concentration de la solution en ions métalliques reste constante : tout se passe comme si le métal « migrait » de l'anode vers la cathode à travers la solution.
\end{propriete}

C'est précisément ce principe qui est exploité dans le \textbf{raffinage électrolytique} du cuivre : une anode de cuivre impur se dissout, seul le cuivre pur se dépose à la cathode, et les impuretés tombent au fond de la cuve sous forme de « boues anodiques » (qui contiennent parfois de l'or et de l'argent, que l'on récupère !).

\courssub{Applications industrielles et artisanales}

\begin{savaistu}[La galvanoplastie, de l'usine aux rues de Douala et Yaoundé]
La \cle{galvanoplastie} consiste à recouvrir un objet conducteur d'une fine couche de métal précieux ou protecteur (or, argent, chrome, nickel) par électrolyse : l'objet à recouvrir sert de \textbf{cathode}. Les jantes, guidons et pots d'échappement chromés des motos-taxis (\emph{bendskins}) et des bicyclettes très prisés à Douala et Yaoundé doivent leur éclat à cette technique. Les lois de Faraday permettent au plateur de calculer exactement la durée nécessaire pour déposer l'épaisseur de chrome voulue : trop peu de temps, le revêtement est terne et fragile ; trop de temps, on gaspille du métal précieux et de l'électricité !
\end{savaistu}

Les applications quantitatives de l'électrolyse sont nombreuses :

\begin{itemize}
\item \textbf{Galvanoplastie et électrodéposition} : dorure de bijoux, argenture de couverts, chromage de pièces métalliques, cuivrage des circuits imprimés. La formule $m = \dfrac{M I t}{n \mathcal{F}}$ permet de régler le couple (intensité, durée) pour obtenir l'épaisseur de dépôt souhaitée.
\item \textbf{Raffinage des métaux} : purification du cuivre (pureté supérieure à $99{,}9\ \%$), du zinc, du nickel, par électrolyse à anode soluble.
\item \textbf{Coulométrie} : technique d'analyse de laboratoire qui « remonte » le raisonnement : on mesure précisément la quantité d'électricité $Q$ consommée par une électrolyse totale d'une espèce, et on en déduit la quantité de matière (donc la concentration) de cette espèce dans l'échantillon. C'est une méthode de dosage très précise, sans étalonnage.
\item \textbf{Préparation industrielle} : production de dichlore et de soude (électrolyse de la saumure), d'aluminium (électrolyse de l'alumine fondue), où le calcul des quantités d'électricité conditionne directement le coût de production.
\end{itemize}

\begin{exercice}[Dépôt d'argent sur une cuillère]
Pour argenter une cuillère, on la place comme cathode dans une solution contenant des ions \ce{Ag+}, avec une anode d'argent. On fait passer un courant d'intensité $I = \SI{0,50}{A}$ pendant $2$ heures. On donne $M(\ce{Ag}) = \SI{107,9}{g.mol^{-1}}$, $\mathcal{F} = \SI{96500}{C.mol^{-1}}$.
\begin{enumerate}
\item Écrire les demi-équations aux deux électrodes.
\item Calculer la masse d'argent déposée sur la cuillère.
\item Quelle masse l'anode a-t-elle perdue ? Justifier.
\end{enumerate}
\end{exercice}

\begin{corrige}[Exercice : dépôt d'argent]
\textbf{1.} Cathode (réduction) : $\ce{Ag+ + e- -> Ag}$. Anode soluble (oxydation) : $\ce{Ag -> Ag+ + e-}$.

\textbf{2.} Durée : $t = 2 \times 3600 = \SI{7200}{s}$. Charge : $Q = 0{,}50 \times 7200 = \SI{3600}{C}$. Moles d'électrons : $n(e^-) = \dfrac{3600}{96500} \approx \SI{3,73 \times 10^{-2}}{mol}$. La demi-équation met en jeu $n = 1$ électron, donc $n(\ce{Ag}) = n(e^-) \approx 0{,}0373\ \mathrm{mol}$. Masse : $m = 0{,}0373 \times 107{,}9 \approx \SI{4,03}{g}$.

\textbf{3.} L'anode étant soluble et les deux demi-équations mettant en jeu le même nombre d'électrons, la conservation de la matière impose que l'anode a perdu exactement la même masse, soit environ $\SI{4,03}{g}$ d'argent.
\end{corrige}

\begin{aretenir}[L'essentiel des aspects quantitatifs]
\begin{itemize}
\item Quantité d'électricité : $Q = I \times t$ (avec $t$ en \textbf{secondes}).
\item Le faraday : $1\ \mathcal{F} \approx \SI{96500}{C.mol^{-1}}$, charge d'une mole d'électrons.
\item Moles d'électrons : $n(e^-) = \dfrac{Q}{\mathcal{F}}$.
\item Masse de produit : $m = \dfrac{M \times I \times t}{n \times \mathcal{F}}$, où $n$ est le nombre d'électrons de la demi-équation.
\item Volume de gaz : $V = n(\text{gaz}) \times V_m$.
\item À anode soluble : masse perdue à l'anode = masse déposée à la cathode.
\end{itemize}
\end{aretenir}

Ces lois quantitatives, établies par Faraday il y a près de deux siècles, gouvernent aujourd'hui encore des industries entières. Mais l'électricité n'est pas toujours une alliée du métal : lorsque ce sont les couples redox de la nature qui travaillent contre nos constructions, le fer se transforme lentement en rouille. C'est le phénomène de \textbf{corrosion}, et les moyens de s'en protéger, que nous étudions dans la dernière section de cette leçon.

\coursec{Corrosion et protection des métaux}

Dans la section précédente, nous avons vu que l'électrolyse permet de \emph{forcer} une transformation chimique grâce à l'énergie électrique : c'est ainsi qu'on obtient l'aluminium, le dichlore ou qu'on dépose une fine couche de métal précieux sur un bijou. Mais la nature, elle, n'a pas besoin de générateur : certaines transformations d'oxydoréduction se produisent \emph{spontanément}, et elles ne sont pas toujours souhaitables. Qui n'a jamais vu une tôle de toiture rouillée, un portail métallique piqué de taches brunes, ou la carcasse d'une vieille voiture abandonnée se désagréger sous le soleil de Douala ? Ce phénomène, c'est la \cle{corrosion}, et il coûte chaque année des milliards de francs CFA aux États et aux particuliers. Comprendre la corrosion, c'est comprendre une réaction d'oxydoréduction spontanée... et apprendre à la combattre, souvent grâce à l'électrolyse elle-même !

\courssub{Qu'est-ce que la corrosion ?}

\courssub{Définition et exemple de la rouille}

Observe une clôture en fer laissée dehors pendant la saison des pluies : après quelques mois, des taches brun-rouge apparaissent, le métal s'effrite, perd sa solidité. Le fer a \emph{réagi} avec son environnement. Ce phénomène porte un nom bien précis.

\begin{definition}[La corrosion]
La \cle{corrosion} est la dégradation d'un métal par \textbf{oxydation} sous l'action du milieu environnant (air humide, eau, sol, eau de mer). C'est une transformation d'oxydoréduction \textbf{spontanée} au cours de laquelle le métal perd des électrons et passe à l'état d'ion ou d'oxyde :
\[ \ce{M -> M^{n+} + n e-} \]
Le métal retourne ainsi, en quelque sorte, à l'état dans lequel on le trouve dans les minerais : la corrosion est l'ennemi naturel de la métallurgie.
\end{definition}

Le cas le plus familier est celui du \textbf{fer}, qui se couvre de \cle{rouille}. La rouille n'est pas un composé simple : c'est un \textbf{oxyde de fer(III) hydraté}, de formule approchée \ce{Fe2O3, x H2O}, poreux et friable. Sa formation nécessite \textbf{à la fois} l'eau \emph{et} le dioxygène de l'air : un clou placé dans de l'eau bouillie (désaérée) et recouverte d'huile ne rouille pas ; un clou dans l'air sec non plus.

L'oxydation du fer s'écrit :
\[ \ce{Fe -> Fe^{2+} + 2 e-} \quad \text{(oxydation, le fer est le réducteur)} \]
puis les ions \ce{Fe^{2+}} sont oxydés en \ce{Fe^{3+}} par le dioxygène, et forment avec l'eau l'oxyde hydraté. On peut résumer le bilan global par :
\[ \ce{4 Fe + 3 O2 + 2x H2O -> 2 Fe2O3, x H2O} \]
Cette transformation est \textbf{lente}, \textbf{spontanée} et se déroule à température ambiante, sans aucun apport d'énergie extérieur.

\courssub{Interprétation électrochimique : les micro-piles}

Pourquoi parle-t-on de corrosion dans une leçon consacrée à l'électrolyse ? Parce que la corrosion est un phénomène \textbf{électrochimique} : elle fonctionne comme une multitude de \emph{micro-piles} à la surface du métal.

Imagine une goutte d'eau de pluie déposée sur une plaque de fer. L'eau contient toujours des ions dissous (elle est légèrement conductrice). À la surface du métal, deux types de zones apparaissent :

\begin{itemize}
\item Des \textbf{zones anodiques} (souvent au centre de la goutte, là où le dioxygène arrive moins facilement) : le fer s'y oxyde et libère des électrons :
\[ \ce{Fe -> Fe^{2+} + 2 e-} \]
\item Des \textbf{zones cathodiques} (au bord de la goutte, là où le dioxygène de l'air est abondant) : le dioxygène dissous s'y réduit en consommant les électrons qui circulent dans le métal :
\[ \ce{O2 + 2 H2O + 4 e- -> 4 HO-} \]
\end{itemize}

Les ions \ce{Fe^{2+}} et \ce{HO-} migrent dans la goutte, se rencontrent, et le dioxygène achève l'oxydation en donnant la rouille \ce{Fe2O3, x H2O}, qui se dépose généralement \emph{entre} les deux zones. Le métal joue le rôle de fil conducteur, l'eau salée celui d'électrolyte : tous les ingrédients d'une pile sont réunis !

\begin{popfigure}
\begin{center}
\begin{tikzpicture}[scale=1.0, every node/.style={font=\small}]
% Plaque de fer
\fill[popmuted!60] (-5,-0.6) rectangle (5,0);
\draw[popdark, thick] (-5,0) -- (5,0);
\node[below] at (0,-0.65) {\textbf{Plaque de fer}};
% Goutte d'eau
\fill[popblue!25] (-2.6,0) arc[start angle=180, end angle=0, radius=2.6] -- cycle;
\draw[popblue, thick] (-2.6,0) arc[start angle=180, end angle=0, radius=2.6];
\node[popblue] at (0,1.6) {\textbf{Goutte d'eau}};
\node[popblue!80!black] at (0,1.1) {(ions dissous)};
% Zone anodique (centre)
\draw[poporange, very thick] (-1,0) -- (1,0);
\node[poporange, below=2pt] at (0,0) {zone anodique};
\node[poporange, align=center] at (0,0.45) {$\ce{Fe -> Fe^{2+} + 2 e-}$};
% Zones cathodiques (bords)
\draw[popgreen!70!black, very thick] (-2.6,0) -- (-1.7,0);
\draw[popgreen!70!black, very thick] (1.7,0) -- (2.6,0);
\node[popgreen!60!black, below=2pt] at (-2.15,0) {cathode};
\node[popgreen!60!black, below=2pt] at (2.15,0) {cathode};
\node[popgreen!60!black, align=center] at (3.9,1.5) {$\ce{O2 + 2 H2O + 4 e- -> 4 HO-}$};
\draw[-{Stealth}, popgreen!60!black] (2.9,1.35) -- (2.3,0.35);
% Circulation des électrons dans le métal
\draw[-{Stealth}, poppurple, thick] (-0.9,-0.3) .. controls (-1.6,-0.5) .. (-2.1,-0.28);
\draw[-{Stealth}, poppurple, thick] (0.9,-0.3) .. controls (1.6,-0.5) .. (2.1,-0.28);
\node[poppurple] at (-3.4,-0.35) {électrons $e^-$};
% O2 de l'air
\draw[-{Stealth}, popdark] (-3.2,2.6) -- (-2.35,1.5);
\node[popdark] at (-3.4,2.8) {$\ce{O2}$ de l'air};
\draw[-{Stealth}, popdark] (3.2,2.6) -- (2.35,1.5);
\node[popdark] at (3.4,2.8) {$\ce{O2}$ de l'air};
% Rouille
\fill[popgold!80!poporange] (-1.55,0.05) circle (0.09);
\fill[popgold!80!poporange] (1.5,0.08) circle (0.12);
\fill[popgold!80!poporange] (1.35,0.03) circle (0.07);
\fill[popgold!80!poporange] (-1.4,0.1) circle (0.11);
\node[poporange!80!black] at (-3.6,0.75) {rouille \ce{Fe2O3, x H2O}};
\draw[-{Stealth}, poporange!80!black] (-2.6,0.65) -- (-1.5,0.18);
\end{tikzpicture}
\end{center}
\legende{Micro-pile de corrosion sous une goutte d'eau : le fer s'oxyde au centre (anode), le dioxygène se réduit aux bords (cathode), et la rouille se forme entre les deux zones.}
\end{popfigure}

\begin{aretenir}[L'essentiel sur le mécanisme]
La corrosion du fer est une \textbf{oxydoréduction spontanée} fonctionnant en micro-piles : oxydation du fer aux zones anodiques, réduction du dioxygène dissous aux zones cathodiques. Elle exige la présence simultanée d'\textbf{eau} (électrolyte) et de \textbf{dioxygène} (oxydant).
\end{aretenir}

\courssub{Facteurs favorisant la corrosion}

Tous les environnements ne corrodent pas à la même vitesse. L'expérience quotidienne au Cameroun le confirme : une tôle rouille beaucoup plus vite à Kribi ou Limbe, au bord de l'océan, qu'à Bafoussam en zone sèche. Les principaux facteurs sont :

\begin{itemize}
\item \textbf{L'humidité} : sans eau, pas d'électrolyte, donc pas de micro-pile. C'est pourquoi la corrosion est rapide pendant la saison des pluies.
\item \textbf{Les sels dissous} : l'eau de mer, riche en \ce{Na+} et \ce{Cl-}, est un excellent électrolyte ; elle accélère la circulation ionique et donc la corrosion. Les véhicules et constructions du littoral (Douala, Limbe) se dégradent plus vite.
\item \textbf{Le contact avec un métal moins réducteur} : si le fer touche du cuivre en milieu humide, c'est le fer (plus réducteur) qui s'oxyde et qui est sacrifié. La corrosion est alors fortement accélérée au niveau du contact.
\item \textbf{Les rayures et les tensions mécaniques} : une rayure sur une peinture ou une tôle galvanisée crée une zone anodique active.
\end{itemize}

\begin{attention}[La rouille ne protège pas le fer !]
Erreur fréquente : croire que la couche de rouille, une fois formée, empêche la corrosion de continuer. C'est \textbf{faux} ! La rouille est \textbf{poreuse} et \textbf{s'écaille} : elle laisse passer l'eau et le dioxygène, et la corrosion progresse en profondeur jusqu'à la destruction complète de la pièce. À l'inverse, l'aluminium se couvre spontanément d'une couche d'\textbf{alumine} \ce{Al2O3} \emph{continue, adhérente et imperméable} qui bloque la corrosion : on dit que l'aluminium se \emph{passive}. C'est pourquoi les casseroles en aluminium durent des années alors qu'une casserole en fer non protégé rouille en quelques mois.
\end{attention}

\courssub{Les méthodes de protection des métaux}

Puisque la corrosion exige le contact du métal avec l'eau et le dioxygène, ou un couplage électrochimique défavorable, deux grandes stratégies s'offrent à nous : \textbf{isoler} le métal du milieu, ou \textbf{détourner} l'oxydation vers un métal « sacrifié ».

\courssub{Protection par revêtement : la barrière}

L'idée la plus simple : interposer une couche imperméable entre le métal et l'atmosphère. C'est le rôle de la \textbf{peinture} (portails, carrosseries, ponts), du \textbf{graissage} (outils, pièces de machines), du \textbf{vernissage} et de la \textbf{plastification} (grillages plastifiés des enclos). Tant que le revêtement est intact, le métal est hors d'atteinte. Mais attention : dès qu'une rayure perce la barrière, la corrosion reprend au point de rupture, parfois même accélérée localement.

\courssub{Protection par métallisation}

On peut aussi recouvrir le métal d'une fine couche d'un \emph{autre} métal, déposée le plus souvent \textbf{par électrolyse} (c'est l'électrodéposition vue dans les sections précédentes) :

\begin{itemize}
\item la \textbf{galvanisation} : dépôt de \textbf{zinc} sur le fer (tôles zinguées des toitures, fils de fer barbelés, seaux) ;
\item le \textbf{chromage} (robinetterie, pare-chocs) et le \textbf{nickelage} (pièces de monnaie, ustensiles), à la fois décoratifs et protecteurs ;
\item l'\textbf{étamage} : couche d'étain sur les boîtes de conserve.
\end{itemize}

Le cas du zinc est particulièrement astucieux : non seulement il couvre le fer, mais en cas de rayure, le zinc étant \emph{plus réducteur} que le fer, c'est lui qui s'oxyde à la place du fer. La protection continue même là où le fer est mis à nu !

\begin{exemplebox}[Les tôles galvanisées des toitures camerounaises]
Les tôles ondulées qui couvrent la plupart des maisons au Cameroun sont des tôles d'acier \textbf{galvanisées} : elles ont reçu un dépôt de zinc d'environ \SI{20}{\micro m} d'épaisseur. Une tôle d'acier nue rouillerait en 2 à 3 ans sous le climat équatorial ; protégée par le zinc, elle résiste couramment \textbf{20 à 30 ans}. Le zinc joue un double rôle : barrière physique tant qu'il est intact, puis anode sacrificielle aux endroits rayés ou percés (trous de clous !). C'est pourquoi les tôles rouillent d'abord autour des fixations.
\end{exemplebox}

\courssub{Protection cathodique : l'anode sacrificielle}

La méthode la plus élégante consiste à \emph{forcer le fer à devenir une cathode}, c'est-à-dire une zone où il ne peut que recevoir des électrons, jamais en perdre. Pour cela, on le met en contact électrique avec un métal \textbf{plus réducteur} que lui : le \textbf{magnésium} ou le \textbf{zinc}. Ce métal s'oxyde à la place du fer :
\[ \ce{Mg -> Mg^{2+} + 2 e-} \quad \text{au lieu de} \quad \ce{Fe -> Fe^{2+} + 2 e-} \]
Le bloc de magnésium se consume lentement : on l'appelle \cle{anode sacrificielle} (ou anode soluble). Il suffit de le remplacer périodiquement. Cette technique protège les canalisations enterrées, les pipelines, les chauffe-eau et les coques de navires.

Une variante, la \textbf{protection par courant imposé}, relie la structure à protéger au pôle négatif d'un générateur : imposée comme cathode, elle ne peut plus s'oxyder. C'est l'électrolyse mise au service de la conservation !

\begin{popfigure}
\begin{center}
\begin{tikzpicture}[scale=1.0, every node/.style={font=\small}]
% Sol
\fill[popgold!25] (-6,0) rectangle (6,-2.8);
\draw[popdark, thick] (-6,0) -- (6,0);
\node[popdark] at (-5.2,0.35) {surface du sol};
% Canalisation
\fill[popmuted!70] (-5.5,-1.2) rectangle (5.5,-1.9);
\draw[popdark, thick] (-5.5,-1.2) rectangle (5.5,-1.9);
\node at (0,-1.55) {\textbf{Canalisation en acier (protégée : cathode)}};
% Anode de magnésium
\fill[poporange!70] (2.8,-2.1) rectangle (3.4,-2.7);
\draw[popdark, thick] (2.8,-2.1) rectangle (3.4,-2.7);
\node[poporange!80!black, align=center] at (4.6,-2.4) {bloc de \textbf{magnésium}\\ (anode sacrificielle)};
% Fil de connexion
\draw[poppurple, very thick] (3.1,-2.1) -- (3.1,-0.5) -- (1.2,-0.5) -- (1.2,-1.2);
\node[poppurple] at (2.15,-0.3) {connexion électrique};
% Circulation des électrons
\draw[-{Stealth}, poppurple, thick] (2.7,-0.5) -- (1.7,-0.5);
\node[poppurple] at (2.2,-0.75) {$e^-$};
% Réactions
\node[poporange!80!black, align=center] at (-3.2,-2.4) {anode : $\ce{Mg -> Mg^{2+} + 2 e-}$};
\node[popgreen!60!black, align=center] at (-0.5,-2.55) {cathode : $\ce{O2 + 2 H2O + 4 e- -> 4 HO-}$};
\draw[-{Stealth}, poporange!80!black] (-1.5,-2.35) -- (2.75,-2.4);
% Ions Mg2+ qui partent
\draw[-{Stealth}, poporange, dashed] (3.4,-2.3) -- (4.0,-1.9);
\node[poporange] at (4.35,-1.85) {$\ce{Mg^{2+}}$};
\end{tikzpicture}
\end{center}
\legende{Protection cathodique d'une canalisation enterrée : le bloc de magnésium, plus réducteur que le fer, s'oxyde à sa place. La canalisation devient une cathode et ne se corrode plus.}
\end{popfigure}

\begin{methode}[Choisir une méthode de protection]
Face à une situation concrète, raisonne en trois étapes :
\begin{enumerate}
\item \textbf{Identifier le milieu} : air sec, air humide, eau salée, sol ? Plus le milieu est conducteur, plus la protection doit être sérieuse.
\item \textbf{Choisir la stratégie} : revêtement \emph{passif} (peinture, vernis, graisse) si l'objet est accessible et le milieu peu agressif ; protection \emph{active} (galvanisation, anode sacrificielle, courant imposé) si l'objet est enterré, immergé ou difficile d'accès.
\item \textbf{Vérifier la règle d'oxydoréduction} : pour une protection active, le métal ajouté doit être \emph{plus réducteur} que le métal protégé (Zn ou Mg pour le fer). Un métal moins réducteur (Cu, Sn en cas de rayure profonde) aggraverait la corrosion !
\end{enumerate}
\end{methode}

\begin{exoresolu}[Pourquoi le zinc protège-t-il le fer, même rayé ?]
Une tôle de fer galvanisée est rayée : le fer est mis à nu au contact du zinc, en milieu humide. Expliquer, à l'aide des demi-équations, pourquoi le fer ne rouille pas. On donne les couples : \ce{Zn^{2+}}/\ce{Zn} et \ce{Fe^{2+}}/\ce{Fe}, le zinc étant plus réducteur que le fer.
\tcblower
\textbf{Solution.} En milieu humide, le contact fer-zinc crée une micro-pile. Le réducteur le plus fort s'oxyde : c'est le zinc.
\begin{itemize}
\item Oxydation (anode, sur le zinc) : \ce{Zn -> Zn^{2+} + 2 e-}
\item Réduction (cathode, sur le fer mis à nu) : \ce{O2 + 2 H2O + 4 e- -> 4 HO-}
\end{itemize}
Le fer, devenu cathode, ne fait que recevoir des électrons : il ne peut pas être oxydé, donc il ne rouille pas. Le zinc se consomme lentement à sa place — il est « sacrifié ». C'est exactement le principe de l'anode sacrificielle, et c'est ce qui explique la longévité des toitures en tôle galvanisée au Cameroun.
\end{exoresolu}

\begin{savaistu}[Les blocs de zinc du port de Douala]
Les coques des navires, en acier, baignent en permanence dans l'eau de mer — un électrolyte redoutable. Au port de Douala, les coques sont protégées par des \textbf{blocs de zinc} (parfois de magnésium ou d'aluminium) boulonnés directement sur la coque, sous la ligne de flottaison. Plus réducteurs que l'acier, ces blocs se corrodent à sa place : les plongeurs d'entretien les remplacent quand ils sont trop rongés. À l'échelle économique, les enjeux sont immenses : on estime que la corrosion coûte chaque année environ 3 à 4 \% du produit intérieur brut des pays industrialisés — ponts à repeindre, pipelines à réparer, carcasses de véhicules, tôles de toiture à remplacer. Comprendre l'oxydoréduction, c'est aussi économiser des milliards !
\end{savaistu}

\begin{exercice}[À toi de jouer]
Un chauffe-eau contient une cuve en acier et une « résistance » entourée d'une tige de magnésium.
\begin{enumerate}
\item Quel est le rôle de la tige de magnésium ? Comment l'appelle-t-on ?
\item Écrire les demi-équations mises en jeu à la surface de la tige et sur la cuve.
\item Que se passerait-il si l'on remplaçait la tige de magnésium par une tige de cuivre, métal moins réducteur que le fer ?
\end{enumerate}
\end{exercice}

\begin{corrige}[Exercice]
\begin{enumerate}
\item La tige de magnésium est une \textbf{anode sacrificielle} : plus réductrice que le fer, elle s'oxyde à la place de la cuve en acier et la protège de la corrosion.
\item À l'anode (tige de Mg) : \ce{Mg -> Mg^{2+} + 2 e-}. À la cathode (cuve en acier) : \ce{O2 + 2 H2O + 4 e- -> 4 HO-}. La cuve, devenue cathode, ne peut pas s'oxyder.
\item Avec une tige de cuivre, c'est le fer, plus réducteur que le cuivre, qui deviendrait l'anode : \ce{Fe -> Fe^{2+} + 2 e-}. La corrosion de la cuve serait \textbf{accélérée} au lieu d'être empêchée. La « protection » serait pire que l'absence de protection !
\end{enumerate}
\end{corrige}

\begin{aretenir}[Bilan de la section]
\begin{itemize}
\item La \textbf{corrosion} est l'oxydation spontanée d'un métal par son milieu ; la \textbf{rouille} (\ce{Fe2O3, x H2O}) résulte de micro-piles : anode \ce{Fe -> Fe^{2+} + 2 e-}, cathode \ce{O2 + 2 H2O + 4 e- -> 4 HO-}.
\item La rouille, poreuse, ne protège pas ; l'alumine, compacte, passivante, protège l'aluminium.
\item Trois familles de protection : \textbf{revêtement} (peinture, graisse), \textbf{métallisation} (galvanisation, chromage, souvent par électrolyse), \textbf{protection cathodique} (anode sacrificielle en Zn ou Mg, ou courant imposé).
\item Règle d'or : le métal sacrifié doit être \textbf{plus réducteur} que le métal protégé.
\end{itemize}
\end{aretenir}

\newpage

\coursec{Travaux pratiques}

\coursec{TP : L'électrolyse en solution aqueuse}

\courssub{Objectifs du TP}

À l'issue de cette séance de travaux pratiques, tu seras capable de :

\begin{itemize}
\item réaliser, en toute sécurité, l'électrolyse de solutions aqueuses de \ce{HCl}, \ce{NaCl}, \ce{NaOH} et \ce{CuSO4} à l'aide d'un montage simple ;
\item recueillir et identifier les gaz dégagés aux électrodes (\ce{H2}, \ce{O2}, \ce{Cl2}) par des tests caractéristiques ;
\item observer et interpréter le dépôt métallique de cuivre à la cathode lors de l'électrolyse de \ce{CuSO4} ;
\item distinguer, pour chaque électrode, l'\cle{oxydation anodique} de la \cle{réduction cathodique} ;
\item confronter les observations expérimentales aux \cle{règles de prévision} issues des potentiels standards, et discuter le rôle de l'eau et du phénomène de \cle{surtension}.
\end{itemize}

\courssub{Matériel et produits}

\begin{center}
\begin{tabularx}{\linewidth}{|l|X|}
\hline
\rowcolor{popblueL} \textbf{Matériel} & \textbf{Produits} \\
\hline
Tube en U (ou bécher de \SI{250}{mL}) & Solution d'acide chlorhydrique \ce{HCl} à \SI{1}{mol.L^{-1}} \\
\hline
2 électrodes de graphite (mine de crayon ou tige de pile usagée) & Solution de chlorure de sodium \ce{NaCl} à \SI{1}{mol.L^{-1}} (eau salée concentrée) \\
\hline
Générateur de tension continue 6--\SI{12}{V} (ou 4 piles de \SI{1,5}{V} en série) & Solution d'hydroxyde de sodium \ce{NaOH} à \SI{1}{mol.L^{-1}} (lessive de soude diluée) \\
\hline
Fils de connexion avec pinces crocodiles & Solution de sulfate de cuivre \ce{CuSO4} à \SI{0,5}{mol.L^{-1}} (bouillie bordelaise clarifiée possible) \\
\hline
2 petits tubes à essais renversés sur les électrodes & Papier iodure d'amidon (ou papier trempé dans un mélange KI + amidon, ou jus de manioc + iodure) \\
\hline
Bûchettes en bois, allumettes, bec Bunsen (ou bougie) & Eau distillée (ou eau de pluie bouillie et refroidie) \\
\hline
Balance au centième (si disponible), chronomètre & Gants, lunettes de protection, blouse \\
\hline
\end{tabularx}
\end{center}

\begin{attention}[Sécurité : les gestes qui sauvent]
\begin{itemize}
\item \textbf{Pictogramme « corrosif »} : \ce{HCl} et \ce{NaOH} attaquent la peau et les yeux. Porte obligatoirement \textbf{blouse, gants et lunettes}. En cas de contact, rince abondamment à l'eau claire pendant plusieurs minutes et préviens l'enseignant.
\item \textbf{Pictogramme « toxique »} : le dichlore \ce{Cl2} dégagé lors des électrolyses de \ce{HCl} et \ce{NaCl} est un gaz vert-jaune irritant pour les voies respiratoires. Travaille \textbf{près d'une fenêtre ouverte} ou sous la hotte, ne respire jamais directement au-dessus du tube, et ne fais pas durer l'électrolyse inutilement.
\item \textbf{Risque d'explosion} : le mélange dihydrogène--dioxygène est détonant. Le test de l'aboiement se fait avec de \textbf{petites quantités} de gaz, dans un tube tenu à distance du visage, ouverture dirigée vers le haut et loin de tout élève.
\item \textbf{Électricité} : utilise uniquement la \textbf{tension continue 6--12 V} fournie par le générateur ou les piles. Jamais de secteur 220 V ! Vérifie les connexions avant de fermer le circuit.
\item Le sulfate de cuivre est \textbf{polluant} : ne jette aucune solution dans l'évier ; verse les résidus dans le bidon de récupération prévu.
\end{itemize}
\end{attention}

\courssub{Schéma du dispositif}

\begin{popfigure}
\begin{center}
\begin{tikzpicture}[scale=0.95, >=Stealth]
% Générateur
\draw[thick, fill=popgoldL] (-3.2,3.4) rectangle (-1.4,4.4);
\node at (-2.3,3.9) {\small Générateur};
\node at (-2.3,4.7) {\small 6--12 V};
\node at (-3.45,3.9) {\small $+$};
\node at (-1.15,3.9) {\small $-$};
% Tube en U
\draw[thick, popblue] (-0.8,0) -- (-0.8,2.6);
\draw[thick, popblue] (0.8,0) -- (0.8,2.6);
\draw[thick, popblue] (-0.8,0) arc (180:360:0.8);
% Solution
\fill[popblue!15] (-0.72,0.05) -- (-0.72,1.9) -- (0.72,1.9) -- (0.72,0.05) arc (180:360:0.72) -- cycle;
\draw[popblue!15, line width=1.4pt] (-0.72,0) arc (180:360:0.72);
% Électrodes graphite
\draw[thick, fill=popdark] (-0.55,1.0) rectangle (-0.35,3.1);
\draw[thick, fill=popdark] (0.35,1.0) rectangle (0.55,3.1);
% Tubes renversés
\draw[thick, poppurple] (-0.68,1.6) rectangle (-0.22,3.0);
\draw[thick, poppurple] (0.22,1.6) rectangle (0.68,3.0);
\node[popmuted] at (-0.45,2.05) {\tiny gaz};
\node[popmuted] at (0.45,2.05) {\tiny gaz};
% Fils (propres)
\draw[thick, popink] (-0.45,3.1) -- (-0.45,3.9) -- (-3.2,3.9); % Anode vers +
\draw[thick, popink] (0.45,3.1) -- (0.45,3.9) -- (-1.4,3.9);  % Cathode vers -
% Bulles
\foreach \y in {0.6,0.9,1.2} {\fill[popblue!50] (-0.45,\y) circle (0.045);}
\foreach \y in {0.7,1.05,1.35} {\fill[popblue!50] (0.45,\y) circle (0.045);}
% Légendes électrodes
\node[below, align=center] at (-0.45,-0.9) {\small \textbf{Anode} ($+$)\\ \small oxydation};
\node[below, align=center] at (0.45,-0.9) {\small \textbf{Cathode} ($-$)\\ \small réduction};
% Légende solution
\node[popmuted, align=center] at (2.6,1.0) {\small solution\\ \small à électrolyser};
\draw[->, popmuted] (2.0,1.0) -- (0.85,1.0);
\end{tikzpicture}
\end{center}
\legende{Dispositif d'électrolyse en tube en U : les tubes renversés emprisonnent les gaz dégagés à chaque électrode. L'anode est reliée à la borne $+$ du générateur, la cathode à la borne $-$.}
\end{popfigure}

\courssub{Protocole}

\begin{enumerate}
\item \textbf{Préparation du montage.} Remplis le tube en U aux trois quarts avec environ \SI{40}{mL} de la première solution (\ce{HCl} à \SI{1}{mol.L^{-1}}). Introduis les deux électrodes de graphite, puis renverse sur chacune un petit tube à essais préalablement rempli de la même solution (aucune bulle d'air ne doit rester piégée au sommet).
\item \textbf{Mise sous tension.} Relie l'électrode de gauche à la borne $+$ du générateur (c'est l'\cle{anode}) et celle de droite à la borne $-$ (c'est la \cle{cathode}). Règle la tension sur \SI{6}{V} puis ferme le circuit. Démarre le chronomètre.
\item \textbf{Recueil des gaz.} Laisse l'électrolyse se dérouler pendant 5 à 10 minutes, jusqu'à recueillir environ \SI{5}{mL} de gaz dans chaque tube. Note l'aspect des électrodes et compare les volumes de gaz recueillis.
\item \textbf{Identification des gaz.} Coupe le circuit. Retourne le tube de la cathode en le maintenant bouché avec le pouce, approche une bûchette enflammée de son ouverture : une \textbf{détonation} (aboiement) caractérise le dihydrogène \ce{H2}. Pour le tube de l'anode, présente une bûchette \textbf{incandescente} (braise sans flamme) : si elle se ravive, c'est le dioxygène \ce{O2} ; si le gaz est verdâtre, âcre, et bleuit un papier iodure d'amidon humide, c'est le dichlore \ce{Cl2}.
\item \textbf{Autres solutions.} Vide le tube en U, rince-le à l'eau distillée, puis recommence les étapes 1 à 4 successivement avec les solutions de \ce{NaCl}, puis de \ce{NaOH}.
\item \textbf{Électrolyse de \ce{CuSO4}.} Avec la solution bleue de sulfate de cuivre, pèse chaque électrode au centième de gramme si la balance est disponible, note les masses $m_1$ et $m_2$. Électrolyse pendant \SI{15}{min} sous \SI{6}{V} en relevant l'intensité $I$ affichée. Observe la couleur de la cathode et l'évolution de la teinte bleue de la solution. Sèche délicatement les électrodes et pèse-les à nouveau.
\item \textbf{Rangement.} Verse les solutions usagées dans le bidon de récupération, rince le matériel, lave-toi les mains.
\end{enumerate}

\courssub{Observations et résultats attendus}

\begin{center}
\begin{tabularx}{\linewidth}{|l|X|X|X|}
\hline
\rowcolor{popblueL} \textbf{Solution} & \textbf{Cathode ($-$) : réduction} & \textbf{Anode ($+$) : oxydation} & \textbf{Autres observations} \\
\hline
\ce{HCl} \SI{1}{mol.L^{-1}} & Dégagement gazeux incolore ; aboiement à la flamme : \ce{H2} & Gaz verdâtre, odeur âcre, bleuit le papier iodure-amidon : \ce{Cl2} & Volumes de gaz voisins ($V_{\ce{H2}} \approx V_{\ce{Cl2}} \approx \SI{5}{mL}$) \\
\hline
\ce{NaCl} \SI{1}{mol.L^{-1}} & Dégagement de \ce{H2} (aboiement) & Dégagement de \ce{Cl2} (test iodure positif) & La solution devient basique près de la cathode (phénolphtaléine rosit) \\
\hline
\ce{NaOH} \SI{1}{mol.L^{-1}} & Dégagement de \ce{H2} & Gaz incolore ravivant la braise : \ce{O2} & $V_{\ce{H2}} \approx 2\,V_{\ce{O2}}$ (ex. \SI{8}{mL} et \SI{4}{mL}) \\
\hline
\ce{CuSO4} \SI{0,5}{mol.L^{-1}} & Dépôt \textbf{rouge} de cuivre métallique sur le graphite & Dégagement de \ce{O2} (braise ravivée) & Le bleu de la solution pâlit ; gain de masse cathodique $\approx \SI{0,15}{g}$ pour $I = \SI{0,5}{A}$ pendant \SI{15}{min} \\
\hline
\end{tabularx}
\end{center}

\courssub{Exploitation}

\begin{exoresolu}[Interprétation des résultats]
\textbf{Questions.} (1) Écris les demi-équations aux électrodes et l'équation bilan pour l'électrolyse de \ce{HCl}. (2) Pour \ce{NaCl}, le sodium métallique ne se forme pas à la cathode : explique pourquoi à l'aide des potentiels standards. (3) Justifie le rapport de volumes $V_{\ce{H2}} = 2\,V_{\ce{O2}}$ observé avec \ce{NaOH}. (4) Pour \ce{CuSO4}, calcule la masse théorique de cuivre déposé par un courant de \SI{0,50}{A} pendant \SI{15}{min} et compare à la mesure. Données : $E°(\ce{Na+}/\ce{Na}) = \SI{-2,71}{V}$, $E°(\ce{H2O}/\ce{H2}) = \SI{-0,83}{V}$, $M(\ce{Cu}) = \SI{63,5}{g.mol^{-1}}$, $F = \SI{96500}{C.mol^{-1}}$.
\tcblower
\textbf{Solution.} 
(1) \textbf{Cathode (réduction)} : $\ce{2H+ + 2e- -> H2}$. \textbf{Anode (oxydation)} : $\ce{2Cl- -> Cl2 + 2e-}$. \textbf{Bilan} : $\ce{2H+ + 2Cl- -> H2 + Cl2}$.
(2) Le réducteur le plus fort (celui dont le couple a le potentiel standard le plus bas) se forme à la cathode. L'eau ($E° = \SI{-0,83}{V}$) est beaucoup plus facile à réduire que \ce{Na+} ($E° = \SI{-2,71}{V}$). C'est donc l'eau qui est réduite : $\ce{2H2O + 2e- -> H2 + 2HO-}$, ce qui explique aussi le rosissement de la phénolphtaléine près de la cathode.
(3) À la cathode, produire 1 mol de \ce{H2} consomme 2 mol d'électrons ($\ce{2H2O + 2e- -> H2 + 2HO-}$) ; à l'anode, produire 1 mol de \ce{O2} en libère 4 ($\ce{4HO- -> O2 + 2H2O + 4e-}$). Pour la même quantité d'électricité traversée, il se forme donc deux fois plus de moles de \ce{H2} que de \ce{O2}, d'où $V_{\ce{H2}} = 2\,V_{\ce{O2}}$ (loi d'Avogadro-Ampère).
(4) Quantité d'électricité : $Q = I\,t = 0{,}50 \times 15 \times 60 = \SI{450}{C}$. Quantité d'électrons : $n(e^-) = \dfrac{450}{96500} = \SI{4,66e-3}{mol}$. La demi-équation cathodique est $\ce{Cu^2+ + 2e- -> Cu}$, donc $n(\ce{Cu}) = \dfrac{n(e^-)}{2} = \SI{2,33e-3}{mol}$. Masse théorique : $m = n \times M = 2{,}33 \times 10^{-3} \times 63{,}5 \approx \SI{0,148}{g} \approx \SI{0,15}{g}$. La masse mesurée est généralement légèrement inférieure (pertes mécaniques lors du séchage, rendement de courant inférieur à 100 \%, passivation partielle).
\end{exoresolu}

\begin{attention}[Le cas \ce{NaCl} et la surtension]
D'après les seuls potentiels standards, l'oxydation de l'eau en \ce{O2} ($E°_{\ce{O2}/\ce{H2O}} = \SI{+1,23}{V}$) devrait être thermodynamiquement favorisée à l'anode par rapport à celle de \ce{Cl-} ($E°_{\ce{Cl2}/\ce{Cl-}} = \SI{+1,36}{V}$). Pourtant, c'est bien le \textbf{dichlore} qui se dégage : la \cle{surtension} du dioxygène sur le graphite (électrode inerte) rend le dégagement de \ce{O2} cinétiquement très lent (forte surpotentiel anodique), tandis que la concentration élevée en \ce{Cl-} favorise la formation de \ce{Cl2}. C'est exactement le principe industriel de la préparation de l'\textbf{eau de Javel} (usine de Bassa, Douala) par électrolyse de saumure !
\end{attention}

\courssub{Corrosion et protection des métaux}

\begin{aretenir}[Mécanisme de la corrosion (rouille du fer)]
La corrosion est une \cle{oxydation spontanée} d'un métal par l'oxygène de l'air en présence d'eau. Pour le fer, il se forme des \textbf{couples locaux} à la surface du métal :
\begin{itemize}
\item \textbf{Anode (zone d'oxydation)} : $\ce{Fe(s) -> Fe^2+(aq) + 2e-}$ \quad ($E°_{\ce{Fe^2+}/\ce{Fe}} = \SI{-0,44}{V}$)
\item \textbf{Cathode (zone de réduction)} : $\ce{O2(g) + 2H2O(l) + 4e- -> 4HO-(aq)}$ \quad ($E° = \SI{+0,40}{V}$ en milieu basique/neutre)
\end{itemize}
\textbf{Bilan global} : $\ce{2Fe(s) + O2(g) + 2H2O(l) -> 2Fe(OH)2(s)}$.
L'hydroxyde de fer(II) vert \ce{Fe(OH)2} s'oxyde rapidement à l'air en \textbf{rouille} : $\ce{4Fe(OH)2 + O2 + 2H2O -> 4Fe(OH)3}$ (brun), qui se déshydrate en \ce{Fe2O3.nH2O} (poudre friable, non protectrice).
\cle{Conditions nécessaires} : présence simultanée d'\textbf{eau} (électrolyte) et de \textbf{dioxygène}.
\end{aretenir}

\begin{methode}[Principales méthodes de protection]
\begin{enumerate}
\item \textbf{Protection par barrière (passivation)} : Isoler le métal de l'air et de l'eau (peinture, vernis, revêtement plastique, émaillage). Ex. : peinture des ponts (pont sur le Wouri), carrosseries de voitures, fûts à huile de palme.
\item \textbf{Galvanisation (protection cathodique sacrificielle par revêtement)} : Recouvrir le fer d'une couche de \textbf{zinc} (galvanisation à chaud ou électrolytique). Le zinc ($E°_{\ce{Zn^2+}/\ce{Zn}} = \SI{-0,76}{V}$) est \textbf{plus réducteur} que le fer. Si la couche est rayée, le zinc s'oxyde à la place du fer : \ce{Zn -> Zn^2+ + 2e-}. Ex. : toitures en tôle ondulée, pylônes électriques, seaux.
\item \textbf{Protection cathodique par anode sacrificielle} : Fixer un bloc de métal plus réducteur (Mg, Zn, Al) sur la structure à protéger (coque de navire, pipeline SONARA, chauffe-eau). L'anode se consume et doit être remplacée périodiquement.
\item \textbf{Protection cathodique par courant imposé} : Imposer un courant continu qui force la structure à être la cathode (pipelines enterrés longue distance).
\item \textbf{Alliages inoxydables} : Ajouter du Chrome ($\ge 10,5\%$) qui forme une couche passive d'oxyde \ce{Cr2O3} adhérente et imperméable (inox pour ustensiles de cuisine, matériel hospitalier, industrie agroalimentaire).
\end{enumerate}
\end{methode}

\begin{savaistu}[Corrosion et contexte camerounais]
\begin{itemize}
\item \textbf{Climat tropical} : L'humidité élevée et les températures chaudes accélèrent considérablement la corrosion (toitures en tôle rouillées après quelques années seulement).
\item \textbf{Eau de mer} : La conductivité élevée due aux sels dissous (NaCl) aggrave la corrosion des navires, pontons et pipelines côtiers (terminal SONARA à Limbé). La protection cathodique est indispensable.
\item \textbf{Acidité des sols} : Dans les zones volcaniques (Mont Cameroun, Ouest), l'acidité des sols attaque les canalisations enterrées et les fondations métalliques.
\item \textbf{Valorisation} : La récupération des métaux (ferrailles) est une filière économique majeure à Douala (Marché de la Ferraille) et Yaoundé, limitant l'impact environnemental et fournissant la matière première aux aciéries locales (CIMENCAM, ALUCAM pour l'aluminium).
\end{itemize}
\end{savaistu}

\courssub{Conclusion}

Ce TP t'a permis de vérifier expérimentalement les lois de l'électrolyse : à la \cle{cathode} (borne $-$) se produit toujours une \textbf{réduction} (dégagement de \ce{H2} ou dépôt de cuivre), à l'\cle{anode} (borne $+$) une \textbf{oxydation} (dégagement de \ce{Cl2} ou de \ce{O2}). Les règles de prévision fondées sur les potentiels standards, complétées par le rôle de l'eau solvant et par le phénomène de surtension, rendent compte de toutes les observations. Enfin, la relation $m = \dfrac{M\,I\,t}{n\,F}$ (lois de Faraday) permet de prévoir quantitativement les masses déposées : c'est le fondement des applications industrielles (électrolyse de la saumure pour l'eau de Javel et la soude, raffinage du cuivre, production d'aluminium à Edea/ALUCAM). La compréhension des couples rédox en jeu explique aussi la \cle{corrosion} des métaux et guide les stratégies de protection (galvanisation, anodes sacrificielles, inox) essentielles pour nos infrastructures.

\newpage

\coursec{Méthodes et exercices résolus}

\coursec{Leçon 09 : Transformations chimiques forcées : l'électrolyse}

\courssub{1. Généralités sur l'électrolyse}

\begin{definition}[Électrolyse]
Une \cle{électrolyse} est une transformation chimique \textbf{forcée} : elle ne se produit que sous l'action d'un courant électrique continu fourni par un générateur externe. Elle a lieu dans un \cle{électrolyseur} constitué d'une cuve contenant l'\cle{électrolyte} (solution aqueuse conductrice ou sel fondu) et de deux \cle{électrodes} (anode et cathode) reliées au générateur.
\end{definition}

\begin{propriete}[Pôles de l'électrolyseur]
\begin{itemize}
\item L'\cle{anode} est l'électrode reliée au pôle \textbf{positif} ($+$) du générateur. C'est le siège d'une \textbf{oxydation} (perte d'électrons).
\item La \cle{cathode} est l'électrode reliée au pôle \textbf{négatif} ($-$) du générateur. C'est le siège d'une \textbf{réduction} (gain d'électrons).
\end{itemize}
Moyen mnémotechnique : \emph{« AnOx -- CatRed »}. \textbf{Attention :} dans une pile, l'anode est le pôle $-$ ; en électrolyse, l'anode est le pôle $+$.
\end{propriete}

\begin{aretenir}[Électrolyses types au programme (électrodes inertes)]
\begin{center}
\begin{tabularx}{\linewidth}{|l|X|X|X|}
\hline
\rowcolor{popblueL}
\textbf{Électrolyte} & \textbf{Cathode ($-$) : Réduction} & \textbf{Anode ($+$) : Oxydation} & \textbf{Bilan global} \\
\hline
\ce{HCl} aqueux & \ce{2 H+ + 2 e- -> H2} & \ce{2 Cl- -> Cl2 + 2 e-} & \ce{2 H+ + 2 Cl- -> H2 + Cl2} \\
\hline
\ce{NaCl} concentré (saumure) & \ce{2 H2O + 2 e- -> H2 + 2 HO-} & \ce{2 Cl- -> Cl2 + 2 e-} & \ce{2 Cl- + 2 H2O -> Cl2 + H2 + 2 HO-} \\
\hline
\ce{NaOH} aqueux & \ce{2 H2O + 2 e- -> H2 + 2 HO-} & \ce{4 HO- -> O2 + 2 H2O + 4 e-} & \ce{4 HO- -> 2 H2 + O2 + 4 HO-} \textit{(soit $2 H2O -> 2 H2 + O2$)} \\
\hline
\ce{CuSO4} aqueux & \ce{Cu^2+ + 2 e- -> Cu} (dépôt rouge) & \ce{2 H2O -> O2 + 4 H+ + 4 e-} & \ce{2 Cu^2+ + 2 H2O -> 2 Cu + O2 + 4 H+} \\
\hline
\end{tabularx}
\end{center}
\textit{N.B. : Pour \ce{NaCl} dilué, l'anode donne \ce{O2} (surtension du \ce{Cl-} sur graphite). Pour \ce{ZnSO4}, voir surtension (Méthode 3).}
\end{aretenir}

\courssub{2. Réactions aux électrodes : réalisation et interprétation}

\begin{methode}[Conduire l'étude expérimentale d'une électrolyse aqueuse]
Pour étudier une électrolyse (\ce{HCl}, \ce{NaCl}, \ce{NaOH}, \ce{CuSO4} en solution aqueuse), procède dans cet ordre :
\begin{enumerate}
\item \cle{Identifier le circuit} : repérer le générateur, l'électrolyte, les deux électrodes. L'électrode au pôle $+$ est l'anode ; au pôle $-$ la cathode.
\item \cle{Observer} : dépôt métallique (rougeâtre pour Cu), dégagement gazeux (bulles), changement de couleur, variation de pH (papier indicateur).
\item \cle{Identifier les gaz} :
\begin{itemize}
\item \ce{H2} (incolore, inodore) : détonation à la flamme.
\item \ce{O2} (incolore, inodore) : ravive une braise incandescente.
\item \ce{Cl2} (jaune-vert, odeur âcre) : bleuit le papier ioduro-amidoné (réaction \ce{Cl2 + 2 I- -> I2 + 2 Cl-}, \ce{I2} + amidon = bleu) et décolore l'indigo.
\end{itemize}
\item \cle{Attribuer les réactions} : Cathode = Réduction (gain $e^-$) ; Anode = Oxydation (perte $e^-$). \emph{« AnOx, CatRed »}.
\item \cle{Écrire le bilan} : additionner les deux demi-équations après égalisation des électrons échangés.
\end{enumerate}
\end{methode}

\begin{exoresolu}[Électrolyse d'une solution de chlorure de sodium concentrée (saumure)]
On électrolyse une saumure entre deux électrodes inertes en graphite. À l'anode : gaz jaune-vert bleuit le papier ioduro-amidoné. À la cathode : gaz incolore détonant, solution basique.
\begin{enumerate}
\item Identifier les gaz.
\item Écrire les demi-équations et leur nature.
\item Écrire le bilan et justifier le caractère forcé.
\end{enumerate}
\tcblower
\textbf{1. Identification.}
Anode : \cle{dichlore} \ce{Cl2} (bleuit papier ioduro-amidoné). Cathode : \cle{dihydrogène} \ce{H2} (détonation).

\textbf{2. Demi-équations.}
Anode (oxydation) : \[ \ce{2 Cl- -> Cl2 + 2 e-} \]
Cathode (réduction) : l'eau est réduite (ions \ce{Na+} non réductibles en aqueux) : \[ \ce{2 H2O + 2 e- -> H2 + 2 HO-} \]
La formation de \ce{HO-} explique la basicité locale.

\textbf{3. Bilan et caractère forcé.}
\[ \ce{2 Cl- + 2 H2O -> Cl2 + H2 + 2 HO-} \]
Cette transformation est \cle{forcée} : elle nécessite l'énergie électrique du générateur et se fait en sens inverse de la réaction spontanée \ce{Cl2 + H2 -> 2 HCl}. C'est le principe industriel de l'\textbf{eau de Javel} : le \ce{Cl2} produit réagit à froid avec \ce{HO-} pour former l'hypochlorite \ce{ClO-} (principe actif).
\end{exoresolu}

\courssub{3. Prévision des réactions d'électrolyse}

\begin{methode}[Appliquer la règle de prévision (électrodes inertes)]
\begin{enumerate}
\item \cle{Inventorier} toutes les espèces : ions du soluté, eau (\ce{H2O}), électrodes (inertes : Pt, graphite).
\item \cle{À la cathode (réduction)} : l'oxydant le plus fort (potentiel standard $E^\circ$ le \textbf{plus élevé}) est réduit.
Comparer le cation métallique à l'eau : $E^\circ(\ce{H2O}/\ce{H2}) \approx \SI{-0,42}{V}$ (pH $= 7$).
\begin{itemize}
\item Si $E^\circ(\ce{M^{n+}/M}) > \SI{-0,42}{V}$ (ex. \ce{Cu^2+}, \ce{Ag+}) $\rightarrow$ dépôt métallique \ce{M}.
\item Si $E^\circ(\ce{M^{n+}/M}) < \SI{-0,42}{V}$ (ex. \ce{Na+}, \ce{K+}, \ce{Ca^2+}, \ce{Al^3+}, \ce{Zn^2+}$^*$) $\rightarrow$ réduction de l'eau : \ce{2 H2O + 2 e- -> H2 + 2 HO-}.
\end{itemize}
$^*$ Sauf surtension (voir Méthode 3).
\item \cle{À l'anode (oxydation)} : le réducteur le plus fort (potentiel $E^\circ$ le \textbf{plus bas}) est oxydé.
\begin{itemize}
\item Anions oxydables (\ce{Cl-}, \ce{Br-}, \ce{I-}, \ce{S^2-}) $\rightarrow$ halogène ou soufre.
\item Anions \textbf{non oxydables} (\ce{SO4^2-}, \ce{NO3-}, \ce{ClO4-}, \ce{F-}) $\rightarrow$ oxydation de l'eau : \ce{2 H2O -> O2 + 4 H+ + 4 e-}.
\end{itemize}
\item \cle{Écrire} les demi-équations équilibrées (atomes + charges) puis le bilan (électrons simplifiés).
\end{enumerate}
\end{methode}

\begin{exoresolu}[Électrolyse de sulfate de cuivre(II) \ce{CuSO4} \SI{0,50}{mol.L^{-1}}]
Données : $E^\circ(\ce{Cu^2+}/\ce{Cu}) = \SI{+0,34}{V}$ ; $E^\circ(\ce{O2}/\ce{H2O}) = \SI{+1,23}{V}$ ; $E^\circ(\ce{H2O}/\ce{H2}) \approx \SI{-0,42}{V}$ (pH $7$).
\begin{enumerate}
\item Électrodes en graphite (inertes) : dépôt rougeâtre cathode, gaz ravivant braise anode. Demi-équations et bilan.
\item Anode en cuivre (soluble) : anode s'amincit. Expliquer et bilan.
\end{enumerate}
\tcblower
\textbf{1. Électrodes inertes (graphite).}
Espèces : \ce{Cu^2+}, \ce{SO4^2-}, \ce{H2O}.
Cathode : $E^\circ(\ce{Cu^2+}/\ce{Cu}) = \SI{+0,34}{V} > \SI{-0,42}{V}$ $\rightarrow$ \ce{Cu^2+} réduit : \[ \ce{Cu^2+ + 2 e- -> Cu} \] (dépôt rougeâtre).
Anode : \ce{SO4^2-} non oxydable $\rightarrow$ eau oxydée : \[ \ce{2 H2O -> O2 + 4 H+ + 4 e-} \] (\ce{O2} ravive braise).
Bilan ($\times 2$ cathode) : \[ \ce{2 Cu^2+ + 2 H2O -> 2 Cu + O2 + 4 H+} \] (solution s'acidifie).

\textbf{2. Anode en cuivre (anode soluble).}
Le Cu métallique ($E^\circ = \SI{+0,34}{V}$) est un réducteur plus fort que l'eau ($E^\circ(\ce{O2}/\ce{H2O}) = \SI{+1,23}{V}$). L'anode s'oxyde : \[ \ce{Cu -> Cu^2+ + 2 e-} \]
Cathode inchangée : \[ \ce{Cu^2+ + 2 e- -> Cu} \]
Bilan : \[ \ce{Cu_{(anode)} -> Cu_{(cathode)}} \]
Transfert de cuivre : concentration \ce{Cu^2+} constante. Principe de l'\textbf{affinage électrolytique} et du \textbf{cuivrage}.
\end{exoresolu}

\courssub{4. Le phénomène de surtension}

\begin{definition}[Surtension]
La \cle{surtension} ($\eta$) est l'écart entre la tension \textbf{réellement nécessaire} pour observer une réaction à une électrode et la tension thermodynamique prévue par les potentiels standards. Elle dépend de la \textbf{nature de l'électrode} et de l'\textbf{espèce} réagissante. C'est un phénomène \textbf{cinétique} (lenteur du transfert de charge ou de l'évolution gazeuse).
\end{definition}

\begin{methode}[Raisonner avec les surtensions]
\begin{enumerate}
\item \cle{Rappeler la prévision thermodynamique} (potentiels standards).
\item \cle{Constater l'écart expérimental} : la réaction prévue n'a pas lieu, une autre se produit.
\item \cle{Invoquer la surtension} : « La surtension de \ldots sur l'électrode en \ldots décale fortement le potentiel réel de dégagement/dépôt, rendant la réaction thermodynamiquement défavorisée plus facile cinétiquement. »
\item \cle{Citer l'application} : extraction du Zn (cathode Zn), accumulateur Pb (surtension \ce{H2} sur Pb), électrolyse saumure (surtension \ce{O2} sur graphite favorise \ce{Cl2}).
\end{enumerate}
\end{methode}

\begin{exoresolu}[Pourquoi dépose-t-on du zinc lors de l'électrolyse de \ce{ZnSO4} aqueux ?]
Donnée : $E^\circ(\ce{Zn^2+}/\ce{Zn}) = \SI{-0,76}{V}$ ; $E^\circ(\ce{H2O}/\ce{H2}) \approx \SI{-0,42}{V}$ (pH $7$). Cathode en zinc.
\begin{enumerate}
\item Prévision thermodynamique cathode.
\item Explication observation (dépôt Zn).
\item Application industrielle.
\end{enumerate}
\tcblower
\textbf{1. Prévision thermodynamique.}
$E^\circ(\ce{H2O}/\ce{H2}) = \SI{-0,42}{V} > E^\circ(\ce{Zn^2+}/\ce{Zn}) = \SI{-0,76}{V}$.
L'eau est l'oxydant le plus fort. Prévision : \[ \ce{2 H2O + 2 e- -> H2 + 2 HO-} \] (pas de dépôt Zn).

\textbf{2. Explication par la surtension.}
Sur une cathode en \textbf{zinc}, le dégagement de \ce{H2} présente une \textbf{forte surtension cathodique} ($\eta_c \approx \SI{-0,6}{V}$ à \SI{-0,8}{V}). Le potentiel réel de réduction de l'eau devient $\approx \SI{-1,0}{V}$ à $\SI{-1,2}{V}$, bien \textbf{inférieur} à $E^\circ(\ce{Zn^2+}/\ce{Zn})$.
$\rightarrow$ \ce{Zn^2+} devient l'oxydant le plus fort \emph{réellement} accessible : \[ \ce{Zn^2+ + 2 e- -> Zn} \]
La surtension a \textbf{inversé la prévision thermodynamique}.

\textbf{3. Application.}
\textbf{Hydrométallurgie du zinc} : extraction électrolytique du Zn depuis des solutions de \ce{ZnSO4} (cathode en Al ou Zn). \textbf{Accumulateur au plomb} : forte surtension de \ce{H2} sur Pb empêche l'électrolyse de l'eau lors de la charge.
\end{exoresolu}

\courssub{5. Aspects quantitatifs de l'électrolyse (Loi de Faraday)}

\begin{methode}[Calculer masses, volumes, durées, énergie]
\begin{enumerate}
\item \cle{Demi-équation} à l'électrode concernée $\rightarrow$ nombre $n_e$ d'électrons par mole de produit.
\item \cle{Quantité d'électricité} : $Q = I \times \Delta t$ ($I$ en A, $\Delta t$ en \textbf{s}, $Q$ en C).
\item \cle{Quantité d'électrons} : $n(e^-) = \dfrac{Q}{F}$ avec $F = \SI{96500}{C.mol^{-1}}$ (constante de Faraday).
\item \cle{Quantité de matière du produit} : $n(\text{produit}) = \dfrac{n(e^-)}{n_e}$.
\item \cle{Grandeurs finales} : masse $m = n \times M$ ; volume gaz $V = n \times V_m$ (préciser $T, P$ ; souvent $V_m = \SI{24,0}{L.mol^{-1}}$ à \SI{20}{\celsius}, \SI{1}{atm}).
\item \cle{Énergie consommée} : $W = U \times I \times \Delta t = U \times Q$ (en J) ; $\SI{1}{kWh} = \SI{3,6e6}{J}$.
\item \cle{Rendement} $\eta$ : $m_{\text{réel}} = \eta \times m_{\text{théo}}$.
\end{enumerate}
\textbf{Réflexes :} Temps en \textbf{secondes} ! Diviser $n(e^-)$ par $n_e$ ! Vérifier chiffres significatifs.
\end{methode}

\begin{exoresolu}[Dépôt de cuivre et énergie]
Électrolyse \ce{CuSO4} (électrodes graphite), $I = \SI{2,0}{A}$, $\Delta t = \SI{30}{min}$.
Données : $M(\ce{Cu}) = \SI{63,5}{g.mol^{-1}}$ ; $F = \SI{96500}{C.mol^{-1}}$ ; $V_m = \SI{24,0}{L.mol^{-1}}$.
\begin{enumerate}
\item Masse Cu déposée.
\item Volume \ce{O2} dégagé anode.
\item Durée pour \SI{5,0}{g} Cu (même $I$).
\end{enumerate}
\tcblower
\textbf{1. Masse Cu.}
Cathode : \ce{Cu^2+ + 2 e- -> Cu} ($n_e = 2$).
$\Delta t = 30 \times 60 = \SI{1800}{s}$.
$Q = 2,0 \times 1800 = \SI{3,6 \times 10^3}{C}$.
$n(e^-) = \frac{3,6 \times 10^3}{96\,500} = \SI{3,73 \times 10^{-2}}{mol}$.
$n(\ce{Cu}) = \frac{3,73 \times 10^{-2}}{2} = \SI{1,87 \times 10^{-2}}{mol}$.
$m(\ce{Cu}) = 1,87 \times 10^{-2} \times 63,5 = \boxed{\SI{1,19}{g} \approx \SI{1,2}{g}}$.

\textbf{2. Volume \ce{O2}.}
Anode : \ce{2 H2O -> O2 + 4 H+ + 4 e-} ($n_e = 4$ par mol \ce{O2}).
$n(\ce{O2}) = \frac{3,73 \times 10^{-2}}{4} = \SI{9,33 \times 10^{-3}}{mol}$.
$V(\ce{O2}) = 9,33 \times 10^{-3} \times 24,0 = \boxed{\SI{0,224}{L} = \SI{224}{mL}}$.

\textbf{3. Durée pour \SI{5,0}{g} Cu.}
$n(\ce{Cu}) = \frac{5,0}{63,5} = \SI{7,87 \times 10^{-2}}{mol}$.
$n(e^-) = 2 \times 7,87 \times 10^{-2} = \SI{0,157}{mol}$.
$Q = 0,157 \times 96\,500 = \SI{1,52 \times 10^4}{C}$.
$\Delta t = \frac{1,52 \times 10^4}{2,0} = \SI{7,6 \times 10^3}{s} = \boxed{\SI{2}{h}\,\SI{7}{min}}$.
\end{exoresolu}

\begin{exoresolu}[Énergie pour production d'aluminium (Hall-Héroult)]
Électrolyse alumine fondue \ce{Al2O3}. Cathode : \ce{Al^3+ + 3 e- -> Al}. $U = \SI{4,5}{V}$, $I = \SI{1,0 \times 10^5}{A}$.
Données : $M(\ce{Al}) = \SI{27,0}{g.mol^{-1}}$ ; $F = \SI{96500}{C.mol^{-1}}$.
\begin{enumerate}
\item Masse Al en \SI{24}{h}.
\item Énergie par kg Al (en kWh).
\end{enumerate}
\tcblower
\textbf{1. Masse Al / jour.}
$\Delta t = 24 \times 3600 = \SI{8,64 \times 10^4}{s}$.
$Q = 1,0 \times 10^5 \times 8,64 \times 10^4 = \SI{8,64 \times 10^9}{C}$.
$n(e^-) = \frac{8,64 \times 10^9}{96\,500} = \SI{8,95 \times 10^4}{mol}$.
$n(\ce{Al}) = \frac{8,95 \times 10^4}{3} = \SI{2,98 \times 10^4}{mol}$.
$m(\ce{Al}) = 2,98 \times 10^4 \times 27,0 = \SI{8,05 \times 10^5}{g} = \boxed{\SI{805}{kg}}$.

\textbf{2. Énergie / kg.}
$W = U \times Q = 4,5 \times 8,64 \times 10^9 = \SI{3,89 \times 10^{10}}{J}$.
$W = \frac{3,89 \times 10^{10}}{3,6 \times 10^6} = \SI{1,08 \times 10^4}{kWh}$.
$\frac{W}{m} = \frac{1,08 \times 10^4}{805} = \boxed{\SI{13,4}{kWh.kg^{-1}}}$.
\textit{Contexte : ALUCAM à Edéa utilise l'hydroélectricité des barrages camerounais pour cette production énergivore.}
\end{exoresolu}

\courssub{6. Corrosion et protection des métaux}

\begin{definition}[Corrosion du fer]
La \cle{corrosion} est l'oxydation \textbf{spontanée} d'un métal par le dioxygène de l'air en présence d'eau. Pour le fer, elle forme la \cle{rouille} (oxyde de fer(III) hydraté \ce{Fe2O3,n H2O} ou hydroxyde \ce{Fe(OH)3}). L'eau salée (bord de mer : Douala, Kribi, Limbé) augmente la conductivité et \textbf{accélère} fortement le phénomène.
\end{definition}

\begin{methode}[Analyser la corrosion et choisir une protection]
\begin{enumerate}
\item \cle{Mécanisme local} (goutte d'eau sur Fe) :
Anode (zone peu oxygénée) : \ce{Fe -> Fe^2+ + 2 e-} (oxydation).
Cathode (zone aérée) : \ce{O2 + 2 H2O + 4 e- -> 4 HO-} (réduction).
\ce{Fe^2+} migre, s'oxyde en \ce{Fe^3+} par \ce{O2}, précipite en \ce{Fe(OH)3} $\rightarrow$ rouille.
\item \cle{Contact bimétallique} : le métal \textbf{plus réducteur} ($E^\circ$ plus bas) s'oxyde (anode) et protège l'autre (cathode).
\item \cle{Protections} :
\begin{itemize}
\item \textbf{Passive} : barrière étanche (peinture, vernis, graisse, plastique, émaillage).
\item \textbf{Revêtement métallique} : \textbf{Galvanisation} (Zn, $E^\circ = \SI{-0,76}{V}$) protège activement même si rayé ; \textbf{Étagage} (Sn, $E^\circ = \SI{-0,14}{V}$) ne protège que si intact (sinon corrosion accélérée du Fe).
\item \textbf{Cathodique} : Anode sacrificielle (Zn, Mg) sur coques, pipelines ; ou courant imposé.
\end{itemize}
\end{enumerate}
\end{methode}

\begin{exoresolu}[Clôture ferreuse en zone côtière camerounaise]
Rouille aux soudures et rayures. Données : $E^\circ(\ce{Fe^2+}/\ce{Fe}) = \SI{-0,44}{V}$ ; $E^\circ(\ce{Zn^2+}/\ce{Zn}) = \SI{-0,76}{V}$ ; $E^\circ(\ce{Cu^2+}/\ce{Cu}) = \SI{+0,34}{V}$.
\begin{enumerate}
\item Demi-équations corrosion + nature.
\item Effet contact Cu vs Zn.
\item Deux méthodes protection + principe.
\end{enumerate}
\tcblower
\textbf{1. Mécanisme corrosion.}
Anode (Fe) : \ce{Fe -> Fe^2+ + 2 e-} \textbf{(Oxydation)}.
Cathode (O2) : \ce{O2 + 2 H2O + 4 e- -> 4 HO-} \textbf{(Réduction)}.
\ce{Fe^2+} + \ce{O2} + \ce{H2O} $\rightarrow$ \ce{Fe^3+} $\rightarrow$ \ce{Fe(OH)3} (rouille). Sel marin $\uparrow$ conductivité $\uparrow$ vitesse.

\textbf{2. Contact bimétallique.}
\textbf{Cuivre} : $E^\circ(\ce{Fe}) = \SI{-0,44}{V} < E^\circ(\ce{Cu}) = \SI{+0,34}{V}$. Fe plus réducteur $\rightarrow$ Fe s'oxyde \textbf{plus vite} (corrosion accélérée).
\textbf{Zinc} : $E^\circ(\ce{Zn}) = \SI{-0,76}{V} < E^\circ(\ce{Fe})$. Zn plus réducteur $\rightarrow$ Zn s'oxyde \textbf{à la place} du Fe. Fe devient cathode, protégé (anode sacrificielle).

\textbf{3. Protections proposées.}
\begin{itemize}
\item \textbf{Peinture antirouille / revêtement plastique} (passive) : barrière \ce{O2}/\ce{H2O}. Entretien rayures indispensable.
\item \textbf{Galvanisation à chaud} (revêtement Zn) \textbf{OU} anodes sacrificielles Zn/Mg fixées sur la clôture (cathodique) : le Zn se consume, le Fe reste intact.
\end{itemize}
\end{exoresolu}

\begin{attention}[Les 5 erreurs fatales au Probatoire]
\begin{enumerate}
\item \textbf{Confondre anode/cathode.} Électrolyse : Anode = pôle $+$ (Oxydation), Cathode = pôle $-$ (Réduction). \emph{AnOx -- CatRed}. (Inverse de la pile !)
\item \textbf{Oublier l'eau réactif.} Aqueux : \ce{Na+}, \ce{K+}, \ce{Ca^2+}, \ce{Al^3+} \textbf{jamais} réduits $\rightarrow$ \ce{H2}. \ce{SO4^2-}, \ce{NO3-} \textbf{jamais} oxydés $\rightarrow$ \ce{O2}.
\item \textbf{Bilan mal équilibré en électrons.} Multiplier les demi-équations pour égaliser les $e^-$ \textbf{avant} d'additionner. Les $e^-$ disparaissent du bilan.
\item \textbf{Erreurs calculs Faraday.} (1) Temps en minutes au lieu de \textbf{secondes}. (2) Oublier de diviser $n(e^-)$ par $n_e$ (nb $e^-$ par mol produit). (3) Confondre $F$ (\SI{96500}{C.mol^{-1}}) et $Q$ (C).
\item \textbf{Ignorer surtension / anode soluble.} Prévision thermodynamique faussée par surtension (ex. Zn, Pb). Anode métal attaquable (Cu, Zn, Ag) $\rightarrow$ \textbf{elle s'oxyde} (se dissout), pas l'eau !
\end{enumerate}
\end{attention}

\newpage

\coursec{Exercices}

\courssub{Je vérifie mes connaissances}

\begin{exercice}[QCM : les bons réflexes]
Pour chaque affirmation, choisir la (ou les) bonne(s) réponse(s).
\begin{enumerate}
\item Une électrolyse est une transformation :
\begin{enumerate}
\item spontanée \quad b) forcée \quad c) toujours totale \quad d) qui ne nécessite aucune énergie.
\end{enumerate}
\item À la cathode se produit toujours :
\begin{enumerate}
\item une oxydation \quad b) une réduction \quad c) un dégagement de dioxygène \quad d) une dissolution du métal.
\end{enumerate}
\item Lors de l'électrolyse d'une solution aqueuse de chlorure de sodium concentrée, on recueille à l'anode :
\begin{enumerate}
\item \ce{H2} \quad b) \ce{O2} \quad c) \ce{Cl2} \quad d) du sodium métallique.
\end{enumerate}
\item La quantité d'électricité ayant traversé un électrolyseur parcouru par un courant $I$ pendant une durée $\Delta t$ vaut :
\begin{enumerate}
\item $Q = \dfrac{I}{\Delta t}$ \quad b) $Q = I \times \Delta t$ \quad c) $Q = \dfrac{\Delta t}{I}$ \quad d) $Q = I^2 \times \Delta t$.
\end{enumerate}
\item Pour protéger une coque de bateau en fer, on fixe sur elle un bloc de :
\begin{enumerate}
\item cuivre \quad b) plomb \quad c) magnésium \quad d) argent.
\end{enumerate}
\end{enumerate}
\end{exercice}

\begin{exercice}[Vrai ou faux, en justifiant]
Répondre par \textbf{vrai} ou \textbf{faux}, puis justifier chaque réponse par une phrase ou une demi-équation électronique.
\begin{enumerate}
\item Dans un électrolyseur, l'anode est reliée à la borne négative du générateur.
\item L'électrolyse de l'eau acidulée produit du dihydrogène à la cathode et du dioxygène à l'anode.
\item La surtension peut empêcher une espèce thermodynamiquement favorisée de réagir à une électrode.
\item Lors du raffinage du cuivre, l'anode de cuivre impur voit sa masse augmenter.
\item Une électrode de graphite est dite \og attaquable \fg.
\end{enumerate}
\end{exercice}

\begin{exercice}[Texte à trous : le vocabulaire de l'électrolyse]
Compléter le texte suivant avec les mots ou expressions : \emph{forcée ; anode ; réduction ; oxydation ; cathode ; générateur ; ions ; électrons ; énergie électrique}.

Une électrolyse est une transformation chimique \ldots\ldots\ldots\ldots\ldots qui ne peut avoir lieu que grâce à l'apport d'\ldots\ldots\ldots\ldots\ldots fournie par un \ldots\ldots\ldots\ldots\ldots L'électrode reliée à la borne positive du générateur s'appelle l'\ldots\ldots\ldots\ldots\ldots ; il s'y produit une \ldots\ldots\ldots\ldots\ldots (perte d'\ldots\ldots\ldots\ldots\ldots). L'électrode reliée à la borne négative s'appelle la \ldots\ldots\ldots\ldots\ldots ; il s'y produit une \ldots\ldots\ldots\ldots\ldots. Dans la solution, le courant est assuré par le déplacement des \ldots\ldots\ldots\ldots\ldots
\end{exercice}

\begin{exercice}[Définitions éclair]
Définir en une ou deux phrases chacun des termes suivants :
\begin{enumerate}
\item électrolyse ;
\item électrode inerte ;
\item surtension ;
\item anode sacrificielle.
\end{enumerate}
\end{exercice}

\courssub{Je m'entraîne}

\begin{exercice}[Électrolyse d'une solution d'acide chlorhydrique]
On réalise l'électrolyse d'une solution aqueuse d'acide chlorhydrique (\ce{H3O+} + \ce{Cl-}) entre deux électrodes de graphite.
\begin{enumerate}
\item Faire le schéma légendé du montage (générateur, électrodes, sens du courant et des électrons).
\item Écrire les demi-équations des réactions se produisant à chaque électrode, en précisant leur nature (oxydation ou réduction).
\item Écrire l'équation bilan de l'électrolyse.
\item Proposer un test expérimental permettant d'identifier chacun des gaz recueillis.
\end{enumerate}
\end{exercice}

\begin{exercice}[Électrolyse d'une solution de sulfate de cuivre]
On électrolyse une solution de sulfate de cuivre (\ce{Cu^2+} + \ce{SO4^2-}) entre une cathode de platine et une anode de graphite. On observe un dépôt rougeâtre sur la cathode et un dégagement gazeux à l'anode.
\begin{enumerate}
\item Identifier les espèces susceptibles de réagir à chaque électrode.
\item À l'aide des potentiels standards $E^\circ(\ce{Cu^2+}/\ce{Cu}) = \SI{0,34}{V}$ et $E^\circ(\ce{H3O+}/\ce{H2}) = \SI{0,00}{V}$, justifier la nature du dépôt observé à la cathode.
\item Écrire les demi-équations aux électrodes et l'équation bilan de l'électrolyse.
\item Que devient la couleur bleue de la solution au cours de l'électrolyse ? Justifier.
\end{enumerate}
\end{exercice}

\begin{exercice}[Dépôt de cuivre : calculs quantitatifs]
On électrolyse une solution de sulfate de cuivre avec un courant d'intensité constante $I = \SI{5,0}{A}$ pendant $\Delta t = \SI{1,0}{h}$. On recueille du dioxygène à l'anode.
\begin{enumerate}
\item Calculer la quantité d'électricité $Q$ ayant traversé le circuit.
\item En déduire la quantité de matière d'électrons échangés, puis la masse de cuivre déposée à la cathode.
\item Calculer le volume de dioxygène dégagé à l'anode.
\end{enumerate}
\textbf{Données :} $M(\ce{Cu}) = \SI{63,5}{g.mol^{-1}}$ ; $1\,\mathcal{F} = \SI{96500}{C.mol^{-1}}$ ; volume molaire des gaz dans les conditions de l'expérience : $V_m = \SI{24}{L.mol^{-1}}$.
\end{exercice}

\begin{exercice}[Raffinage du cuivre]
Dans une raffinerie, on électrolyse une solution de sulfate de cuivre acidifiée entre une anode de cuivre impur et une cathode de cuivre pur. L'anode perd une masse $m = \SI{3,2}{g}$ de cuivre, le courant ayant une intensité $I = \SI{2,0}{A}$.
\begin{enumerate}
\item Écrire la demi-équation se produisant à l'anode. Comment appelle-t-on ce type d'électrolyse ?
\item Calculer la quantité d'électricité ayant traversé le circuit.
\item En déduire la durée de l'électrolyse, exprimée en heures et minutes.
\end{enumerate}
\textbf{Données :} $M(\ce{Cu}) = \SI{63,5}{g.mol^{-1}}$ ; $1\,\mathcal{F} = \SI{96500}{C.mol^{-1}}$.
\end{exercice}

\courssub{Je me prépare au Probatoire}

\begin{exercice}[Fabrication de l'eau de Javel par électrolyse de la saumure]
Dans une petite unité de production installée à Douala, on fabrique de l'eau de Javel en électrolysant une solution concentrée de chlorure de sodium (saumure) entre des électrodes inertes en graphite. On observe un dégagement gazeux à chaque électrode.

\textbf{Données :} potentiels standards : $E^\circ(\ce{Na+}/\ce{Na}) = \SI{-2,71}{V}$ ; $E^\circ(\ce{H2O}/\ce{H2}) = \SI{-0,83}{V}$ (à pH = 7) ; $E^\circ(\ce{Cl2}/\ce{Cl-}) = \SI{1,36}{V}$ ; $E^\circ(\ce{O2}/\ce{H2O}) = \SI{0,82}{V}$ (à pH = 7). $M(\ce{Cl}) = \SI{35,5}{g.mol^{-1}}$ ; $1\,\mathcal{F} = \SI{96500}{C.mol^{-1}}$ ; $V_m = \SI{24}{L.mol^{-1}}$.

\begin{enumerate}
\item Faire le schéma légendé de l'électrolyseur en précisant la polarité des électrodes et le sens de circulation des électrons.
\item À la cathode, deux espèces peuvent être réduites : \ce{Na+} et \ce{H2O}. À l'aide des potentiels standards, montrer que c'est l'eau qui est réduite et non les ions sodium. Écrire la demi-équation correspondante.
\item À l'anode, la surtension du dioxygène sur graphite est élevée, contrairement à celle du dichlore. Expliquer pourquoi c'est \ce{Cl2} qui se dégage et non \ce{O2}. Écrire la demi-équation anodique.
\item Écrire l'équation bilan de l'électrolyse.
\item Le dichlore formé réagit ensuite avec les ions hydroxyde présents en solution selon l'équation :
\[ \ce{Cl2 + 2 HO- -> Cl- + ClO- + H2O} \]
Nommer l'ion responsable des propriétés désinfectantes de l'eau de Javel.
\item L'électrolyse fonctionne sous un courant $I = \SI{10,0}{A}$ pendant $\Delta t = \SI{2,0}{h}$. Calculer le volume de dichlore recueilli à l'anode, puis la quantité de matière d'ions hypochlorite \ce{ClO-} obtenue.
\end{enumerate}
\end{exercice}

\begin{exercice}[Étude de la courbe caractéristique d'un électrolyseur]
Au laboratoire d'un lycée de Yaoundé, un groupe d'élèves électrolyse une solution aqueuse d'acide sulfurique dilué entre deux électrodes de platine. Ils font varier la tension $U$ aux bornes de l'électrolyseur et relèvent l'intensité $I$ du courant :

\begin{center}
\begin{tabular}{|c|c|c|c|c|c|c|c|}
\hline
$U$ (V) & 0,5 & 1,0 & 1,5 & 2,0 & 2,5 & 3,0 & 3,5 \\
\hline
$I$ (mA) & 0 & 0 & 5 & 40 & 90 & 140 & 190 \\
\hline
\end{tabular}
\end{center}

\textbf{Données :} $E^\circ(\ce{O2}/\ce{H2O}) = \SI{1,23}{V}$ ; $E^\circ(\ce{H+}/\ce{H2}) = \SI{0,00}{V}$ ; $1\,\mathcal{F} = \SI{96500}{C.mol^{-1}}$ ; $V_m = \SI{24}{L.mol^{-1}}$.

\begin{enumerate}
\item Tracer la courbe $I = f(U)$ sur papier millimétré (échelles : \SI{1}{cm} pour \SI{0,5}{V} et \SI{1}{cm} pour \SI{20}{mA}).
\item Déterminer graphiquement la tension de seuil $U_s$ de l'électrolyseur.
\item Calculer la tension théorique minimale $E^\circ(\ce{O2}/\ce{H2O}) - E^\circ(\ce{H+}/\ce{H2})$ de l'électrolyse. Comparer avec $U_s$ et interpréter l'écart observé en invoquant le phénomène de surtension.
\item Écrire les demi-équations aux électrodes et l'équation bilan de l'électrolyse.
\item On fixe $U = \SI{3,5}{V}$ pendant $\Delta t = \SI{30}{min}$. En utilisant le tableau, calculer les volumes de dihydrogène et de dioxygène recueillis.
\end{enumerate}
\end{exercice}

\begin{exercice}[Protection d'une canalisation enterrée et galvanisation au zinc]
\textbf{Partie A.} Une canalisation en fer enterrée dans le sol humide de la région de Bafoussam se corrode lentement. Pour la protéger, un technicien la relie par un fil conducteur à un bloc de magnésium enfoui à proximité.

\textbf{Données :} $E^\circ(\ce{Fe^2+}/\ce{Fe}) = \SI{-0,44}{V}$ ; $E^\circ(\ce{Mg^2+}/\ce{Mg}) = \SI{-2,37}{V}$ ; $E^\circ(\ce{Zn^2+}/\ce{Zn}) = \SI{-0,76}{V}$ ; $E^\circ(\ce{H2O}/\ce{H2}) = \SI{-0,83}{V}$ (à pH = 7). $M(\ce{Mg}) = \SI{24,3}{g.mol^{-1}}$ ; $M(\ce{Zn}) = \SI{65,4}{g.mol^{-1}}$ ; $M(\ce{Fe}) = \SI{55,8}{g.mol^{-1}}$ ; $1\,\mathcal{F} = \SI{96500}{C.mol^{-1}}$.

\begin{enumerate}
\item À l'aide des potentiels standards, expliquer pourquoi c'est le magnésium qui s'oxyde et non le fer. Comment nomme-t-on ce type de protection ?
\item Écrire les demi-équations mises en jeu au niveau du bloc de magnésium et à la surface de la canalisation.
\item Le bloc de magnésium a une masse initiale de \SI{2,0}{kg}. Sachant que le courant de protection moyen vaut $I = \SI{50}{mA}$, estimer la durée de vie du bloc en années (on prendra 1 an $= \SI{3,15e7}{s}$).
\end{enumerate}

\textbf{Partie B.} Dans un atelier de Douala, on veut recouvrir des pièces d'acier d'une couche de zinc par électrolyse d'une solution aqueuse de sulfate de zinc.
\begin{enumerate}[resume]
\item Le bain de galvanisation est acide. La réduction des ions \ce{H+} ($E^\circ = \SI{0,00}{V}$) est thermodynamiquement plus favorable que celle des ions \ce{Zn^2+} ($E^\circ = \SI{-0,76}{V}$). Expliquer, en invoquant la surtension du dihydrogène sur le zinc, pourquoi le zinc peut néanmoins être déposé en solution aqueuse.
\item On veut déposer une masse $m = \SI{6,54}{g}$ de zinc avec un courant $I = \SI{4,0}{A}$. Calculer la durée nécessaire de l'électrolyse.
\end{enumerate}
\end{exercice}



\newpage

\coursec{Corrigés des exercices}

\courssub{Je vérifie mes connaissances}

\begin{corrige}[Exercice 1 --- QCM : les bons réflexes]
\begin{enumerate}
\item \textbf{b)} Une électrolyse est une transformation \cle{forcée} : elle ne se produit que sous l'action d'un courant électrique imposé par un générateur extérieur. Sans apport d'énergie électrique, elle s'arrête.
\item \textbf{b)} Par définition, la \cle{cathode} est l'électrode où se produit une \cle{réduction} (gain d'électrons). Moyen mnémotechnique : \emph{cathode -- réduction} (mêmes voyelles), \emph{anode -- oxydation}.
\item \textbf{c)} En solution aqueuse concentrée en \ce{Cl-}, l'oxydation des ions chlorure $\ce{2Cl- -> Cl2 + 2e-}$ est favorisée par rapport à l'oxydation de l'eau, grâce à la \cle{surtension} élevée du dioxygène sur le graphite.
\item \textbf{b)} Définition de la quantité d'électricité : $Q = I \times \Delta t$, en coulombs (C) si $I$ est en ampères (A) et $\Delta t$ en secondes (s).
\item \textbf{c)} Le magnésium a un potentiel standard plus négatif que le fer : $E^\circ(\ce{Mg^2+}/\ce{Mg}) = \SI{-2,37}{V} < E^\circ(\ce{Fe^2+}/\ce{Fe}) = \SI{-0,44}{V}$. C'est une \cle{anode sacrificielle} qui s'oxyde à la place du fer. Le cuivre, le plomb et l'argent, moins réducteurs que le fer, accéléreraient au contraire sa corrosion.
\end{enumerate}
\end{corrige}

\begin{corrige}[Exercice 2 --- Vrai ou faux, en justifiant]
\begin{enumerate}
\item \textbf{Faux.} Dans un électrolyseur, l'\cle{anode} est reliée à la \cle{borne positive} du générateur. C'est le générateur qui impose les polarités ; l'anode est le siège de l'oxydation.
\item \textbf{Vrai.} Cathode (réduction) : $\ce{2H3O+ + 2e- -> H2 + 2H2O}$ ; anode (oxydation) : $\ce{2H2O -> O2 + 4H+ + 4e-}$. Bilan : $\ce{2H2O -> 2H2 + O2}$.
\item \textbf{Vrai.} La \cle{surtension} traduit une lenteur cinétique à l'interface électrode/solution. Exemple : sur graphite, l'oxydation de l'eau en \ce{O2} subit une forte surtension, ce qui permet l'oxydation des ions \ce{Cl-} en \ce{Cl2}, pourtant thermodynamiquement moins favorisée ($E^\circ(\ce{Cl2}/\ce{Cl-}) = \SI{1,36}{V} > E^\circ(\ce{O2}/\ce{H2O}) = \SI{1,23}{V}$).
\item \textbf{Faux.} Lors du raffinage, l'anode de cuivre impur est \emph{attaquable} : elle s'oxyde selon $\ce{Cu -> Cu^2+ + 2e-}$, donc sa masse \textbf{diminue}. C'est la cathode de cuivre pur qui grossit par dépôt : $\ce{Cu^2+ + 2e- -> Cu}$.
\item \textbf{Faux.} Le graphite est une électrode \cle{inerte} : il ne participe pas à la réaction d'oxydoréduction, il assure seulement le transfert des électrons. Une électrode \emph{attaquable} (cuivre, zinc, fer...) participe à la réaction et se consomme.
\end{enumerate}
\end{corrige}

\begin{corrige}[Exercice 3 --- Texte à trous : le vocabulaire de l'électrolyse]
Une électrolyse est une transformation chimique \textbf{forcée} qui ne peut avoir lieu que grâce à l'apport d'\textbf{énergie électrique} fournie par un \textbf{générateur}. L'électrode reliée à la borne positive du générateur s'appelle l'\textbf{anode} ; il s'y produit une \textbf{oxydation} (perte d'\textbf{électrons}). L'électrode reliée à la borne négative s'appelle la \textbf{cathode} ; il s'y produit une \textbf{réduction}. Dans la solution, le courant est assuré par le déplacement des \textbf{ions}.
\end{corrige}

\begin{attention}[Point de vigilance]
Dans les fils et les électrodes, le courant est dû au déplacement des \textbf{électrons} ; dans l'électrolyte, il est dû au déplacement des \textbf{ions} (cations vers la cathode, anions vers l'anode). Ne pas confondre les deux porteurs de charge.
\end{attention}

\begin{corrige}[Exercice 4 --- Définitions éclair]
\begin{enumerate}
\item \textbf{Électrolyse :} transformation chimique forcée (non spontanée) provoquée par le passage d'un courant électrique dans un électrolyte (solution ionique ou composé fondu), se traduisant par une oxydation à l'anode et une réduction à la cathode.
\item \textbf{Électrode inerte :} électrode (platine, graphite) qui ne participe pas à la réaction chimique ; elle sert uniquement de siège au transfert des électrons entre le circuit et les espèces en solution. Sa masse ne varie pas.
\item \textbf{Surtension :} écart entre la tension réellement nécessaire pour qu'une réaction d'électrode se produise à vitesse appréciable et la tension d'équilibre prévue par la thermodynamique. Elle traduit une barrière cinétique à l'interface électrode/solution et dépend de la nature de l'électrode et de l'espèce.
\item \textbf{Anode sacrificielle :} métal de potentiel standard plus négatif que celui de la structure à protéger (fer, acier), relié électriquement à celle-ci. Il s'oxyde préférentiellement --- il se \og sacrifie \fg{} --- et la structure devient la cathode d'une pile : c'est la \cle{protection cathodique}.
\end{enumerate}
\end{corrige}

\courssub{Je m'entraîne}

\begin{corrige}[Exercice 5 --- Électrolyse d'une solution d'acide chlorhydrique]
\begin{enumerate}
\item \textbf{Schéma légendé :}
\begin{popfigure}
\begin{tikzpicture}[scale=0.85, >=Stealth]
% Becher
\draw[thick] (0,0) -- (0,3.2) (3,0) -- (3,3.2);
\draw[thick] (-0.2,0) -- (3.2,0);
\fill[popblueL!40] (0.05,0.05) rectangle (2.95,2.4);
% Electrodes
\draw[thick, fill=popdark!60] (0.5,1.2) rectangle (0.9,4.2);
\draw[thick, fill=popdark!60] (2.1,1.2) rectangle (2.5,4.2);
\node[above, font=\small] at (0.7,4.2) {Anode ($+$)};
\node[above, font=\small] at (2.3,4.2) {Cathode ($-$)};
% Generateur
\draw[thick] (5.2,4.2) circle (0.55);
\node at (5.2,4.2) {G};
\node[font=\small] at (5.2,5.0) {$+$};
\node[font=\small] at (5.2,3.4) {$-$};
\draw[thick] (0.9,4.2) -- (0.9,4.9) -- (4.65,4.9) -- (4.65,4.2);
\draw[thick] (2.5,4.2) -- (2.5,3.5) -- (5.75,3.5) -- (5.75,4.2);
% Sens electrons
\draw[popink, ->, thick] (1.1,5.15) -- node[above, font=\small] {$e^-$} (4.4,5.15);
\draw[popink, ->, thick] (5.9,3.3) -- node[right, font=\small] {$e^-$} (5.9,2.6);
% Bulles
\foreach \y in {1.4,1.8,2.2} {\draw[poppink, fill=poppink!40] (0.7,\y) circle (0.07);}
\foreach \y in {1.4,1.8,2.2} {\draw[popgreen, fill=popgreen!40] (2.3,\y) circle (0.07);}
\node[poppink, font=\small, left] at (0.45,1.8) {\ce{Cl2}};
\node[popgreen!70!black, font=\small, right] at (2.55,1.8) {\ce{H2}};
\node[font=\small] at (1.5,0.5) {solution \ce{H3O+} + \ce{Cl-}};
\end{tikzpicture}
\legende{Électrolyse de l'acide chlorhydrique entre électrodes de graphite. Les électrons quittent l'anode et arrivent à la cathode par le circuit extérieur.}
\end{popfigure}

\item \textbf{Demi-équations :}
\begin{itemize}
\item \textbf{Cathode (borne $-$), réduction :} les ions oxonium sont réduits :
\[ \ce{2H3O+ + 2e- -> H2 + 2H2O} \]
\item \textbf{Anode (borne $+$), oxydation :} les ions chlorure sont oxydés (l'oxydation de l'eau est défavorisée par la surtension de \ce{O2} sur graphite) :
\[ \ce{2Cl- -> Cl2 + 2e-} \]
\end{itemize}

\item \textbf{Équation bilan :} on additionne membre à membre en simplifiant les électrons :
\[ \ce{2H3O+ + 2Cl- -> H2 + Cl2 + 2H2O} \]

\item \textbf{Tests d'identification :}
\begin{itemize}
\item \textbf{Dihydrogène \ce{H2} (cathode) :} on approche une buchette enflammée de l'ouverture du tube : on entend une petite détonation (aboiement caractéristique).
\item \textbf{Dichlore \ce{Cl2} (anode) :} gaz verdâtre d'odeur suffocante ; il décolore un papier indicateur humide (pouvoir blanchissant) et bleuit un papier imbibé d'iodure de potassium amidonné ($\ce{Cl2 + 2I- -> I2 + 2Cl-}$, le diiode colore l'amidon en bleu-violet).
\end{itemize}
\end{enumerate}
\end{corrige}

\begin{corrige}[Exercice 6 --- Électrolyse d'une solution de sulfate de cuivre]
\begin{enumerate}
\item \textbf{Espèces susceptibles de réagir :}
\begin{itemize}
\item \textbf{À la cathode (réduction) :} les ions \ce{Cu^2+} et l'eau \ce{H2O} (et les traces de \ce{H3O+}).
\item \textbf{À l'anode (oxydation) :} l'eau \ce{H2O} ; les ions \ce{SO4^2-} ne s'oxydent pas en solution aqueuse ; le graphite est inerte.
\end{itemize}

\item \textbf{Nature du dépôt :} à la cathode, c'est l'oxydant du couple de potentiel standard \textbf{le plus élevé} qui est réduit en priorité :
\[ E^\circ(\ce{Cu^2+}/\ce{Cu}) = \SI{0,34}{V} > E^\circ(\ce{H3O+}/\ce{H2}) = \SI{0,00}{V} \]
Les ions \ce{Cu^2+} sont donc réduits à la place des ions \ce{H3O+} : le dépôt rougeâtre est du \cle{cuivre métallique}.

\item \textbf{Demi-équations et bilan :}
\begin{itemize}
\item Cathode : $\ce{Cu^2+ + 2e- -> Cu}$ \quad (réduction, dépôt rougeâtre)
\item Anode : $\ce{2H2O -> O2 + 4H+ + 4e-}$ \quad (oxydation de l'eau, dégagement de dioxygène)
\end{itemize}
On multiplie la demi-équation cathodique par 2 pour égaliser les électrons échangés :
\[ \ce{2Cu^2+ + 2H2O -> 2Cu + O2 + 4H+} \]

\item \textbf{Évolution de la couleur bleue :} la couleur bleue, due aux ions \ce{Cu^2+} hydratés, \textbf{s'atténue progressivement puis disparaît}. En effet, les ions \ce{Cu^2+} sont consommés à la cathode : leur concentration diminue au cours de l'électrolyse. De plus, la solution s'acidifie (formation d'ions \ce{H+} à l'anode).
\end{enumerate}
\end{corrige}

\begin{corrige}[Exercice 7 --- Dépôt de cuivre : calculs quantitatifs]
\begin{enumerate}
\item \textbf{Quantité d'électricité :} on convertit la durée en secondes : $\Delta t = \SI{1,0}{h} = \SI{3600}{s}$.
\[ Q = I \times \Delta t = 5,0 \times 3600 = \SI{1,8e4}{C} \]

\item \textbf{Quantité d'électrons puis masse de cuivre :}
\[ n(e^-) = \frac{Q}{\mathcal{F}} = \frac{1,8 \times 10^4}{96\,500} = \SI{0,187}{mol} \]
À la cathode : $\ce{Cu^2+ + 2e- -> Cu}$, donc $n(\ce{Cu}) = \dfrac{n(e^-)}{2}$ :
\[ n(\ce{Cu}) = \frac{0,187}{2} = \SI{9,3e-2}{mol} \]
\[ m(\ce{Cu}) = n(\ce{Cu}) \times M(\ce{Cu}) = 9,3 \times 10^{-2} \times 63,5 \approx \SI{5,9}{g} \]

\item \textbf{Volume de dioxygène :} à l'anode : $\ce{2H2O -> O2 + 4H+ + 4e-}$, donc $n(\ce{O2}) = \dfrac{n(e^-)}{4}$ :
\[ n(\ce{O2}) = \frac{0,187}{4} = \SI{4,7e-2}{mol} \]
\[ V(\ce{O2}) = n(\ce{O2}) \times V_m = 4,7 \times 10^{-2} \times 24 \approx \SI{1,1}{L} \]
\end{enumerate}
\end{corrige}

\begin{attention}[Points de vigilance]
\begin{itemize}
\item Convertir systématiquement la durée en \textbf{secondes} avant de calculer $Q = I\,\Delta t$.
\item Bien tenir compte de la \textbf{stoechiométrie électronique} : 2 mol d'électrons pour 1 mol de \ce{Cu}, 4 mol d'électrons pour 1 mol de \ce{O2}. C'est l'erreur la plus fréquente au Probatoire.
\item Vérification rapide : $V(\ce{H2})$ serait le double de $V(\ce{O2})$ pour une même quantité d'électricité.
\end{itemize}
\end{attention}

\begin{corrige}[Exercice 8 --- Raffinage du cuivre]
\begin{enumerate}
\item \textbf{Demi-équation anodique :} l'anode de cuivre impur se dissout :
\[ \ce{Cu -> Cu^2+ + 2e-} \]
C'est une électrolyse à \cle{anode soluble} (ou \cle{anode attaquable}) : le métal de l'anode participe à la réaction.

\item \textbf{Quantité d'électricité :}
\[ n(\ce{Cu}) = \frac{m}{M(\ce{Cu})} = \frac{3,2}{63,5} = \SI{5,04e-2}{mol} \]
D'après la demi-équation, $n(e^-) = 2\,n(\ce{Cu}) = 2 \times 5,04 \times 10^{-2} = \SI{0,101}{mol}$.
\[ Q = n(e^-) \times \mathcal{F} = 0,101 \times 96\,500 \approx \SI{9,7e3}{C} \]

\item \textbf{Durée de l'électrolyse :}
\[ \Delta t = \frac{Q}{I} = \frac{\num{9,7e3}}{2,0} = \SI{4,86e3}{s} \]
Conversion : $\dfrac{4860}{3600} = 1,35\,\text{h}$, et $0,35 \times 60 = 21$ min.
\[ \boxed{\Delta t \approx \text{1 h 21 min}} \]
\end{enumerate}
\end{corrige}

\begin{attention}[Point de vigilance]
Au raffinage, la masse perdue par l'anode n'est pas exactement égale à la masse gagnée par la cathode : les impuretés métalliques (argent, or) ne s'oxydent pas et tombent au fond de la cuve sous forme de \emph{boues anodiques}, dont la récupération est économiquement intéressante.
\end{attention}

\courssub{Je me prépare au Probatoire}

\begin{corrige}[Exercice 9 --- Fabrication de l'eau de Javel par électrolyse de la saumure]
\begin{enumerate}
\item \textbf{Schéma légendé :}
\begin{popfigure}
\begin{tikzpicture}[scale=0.85, >=Stealth]
\draw[thick] (0,0) -- (0,3.2) (3,0) -- (3,3.2);
\draw[thick] (-0.2,0) -- (3.2,0);
\fill[popblueL!40] (0.05,0.05) rectangle (2.95,2.4);
\draw[thick, fill=popdark!60] (0.5,1.2) rectangle (0.9,4.2);
\draw[thick, fill=popdark!60] (2.1,1.2) rectangle (2.5,4.2);
\node[above, font=\small] at (0.7,4.2) {Anode ($+$)};
\node[above, font=\small] at (2.3,4.2) {Cathode ($-$)};
\draw[thick] (5.2,4.2) circle (0.55);
\node at (5.2,4.2) {G};
\node[font=\small] at (5.2,5.0) {$+$};
\node[font=\small] at (5.2,3.4) {$-$};
\draw[thick] (0.9,4.2) -- (0.9,4.9) -- (4.65,4.9) -- (4.65,4.2);
\draw[thick] (2.5,4.2) -- (2.5,3.5) -- (5.75,3.5) -- (5.75,4.2);
\draw[popink, ->, thick] (1.1,5.15) -- node[above, font=\small] {$e^-$} (4.4,5.15);
\foreach \y in {1.4,1.8,2.2} {\draw[poppink, fill=poppink!40] (0.7,\y) circle (0.07);}
\foreach \y in {1.4,1.8,2.2} {\draw[popgreen, fill=popgreen!40] (2.3,\y) circle (0.07);}
\node[poppink, font=\small, left] at (0.45,1.8) {\ce{Cl2}};
\node[popgreen!70!black, font=\small, right] at (2.55,1.8) {\ce{H2}};
\node[font=\small] at (1.5,0.5) {saumure \ce{Na+} + \ce{Cl-}};
\end{tikzpicture}
\legende{Électrolyse de la saumure entre électrodes inertes de graphite : \ce{Cl2} à l'anode, \ce{H2} à la cathode.}
\end{popfigure}

\item \textbf{Cathode :} on compare les potentiels des deux couples en présence (milieu neutre/basique) :
\[ E(\ce{H2O}/\ce{H2}) \approx \SI{-0,83}{V} \gg E^\circ(\ce{Na+}/\ce{Na}) = \SI{-2,71}{V} \]
L'oxydant le plus facile à réduire est celui du couple de potentiel le plus élevé : c'est \textbf{l'eau}, et non les ions sodium. Le sodium métallique ne peut d'ailleurs pas exister au contact de l'eau. Demi-équation :
\[ \ce{2H2O + 2e- -> H2 + 2HO-} \]

\item \textbf{Anode :} thermodynamiquement, l'eau devrait s'oxyder plus facilement que les ions chlorure car $E(\ce{O2}/\ce{H2O}) \approx \SI{0,82}{V}$ (à pH neutre) $< E^\circ(\ce{Cl2}/\ce{Cl-}) = \SI{1,36}{V}$. \textbf{Mais} la surtension du dioxygène sur le graphite est élevée, alors que celle du dichlore est faible : la tension réellement nécessaire pour dégager \ce{O2} devient supérieure à celle requise pour dégager \ce{Cl2}. De plus, la solution est concentrée en \ce{Cl-}. C'est donc le dichlore qui se forme :
\[ \ce{2Cl- -> Cl2 + 2e-} \]

\item \textbf{Équation bilan :} en additionnant les deux demi-équations (2 électrons échangés de part et d'autre) :
\[ \ce{2Cl- + 2H2O -> Cl2 + H2 + 2HO-} \]

\item L'ion responsable des propriétés désinfectantes et blanchissantes de l'eau de Javel est l'ion \cle{hypochlorite} \ce{ClO-}.

\item \textbf{Calculs quantitatifs :} $\Delta t = \SI{2,0}{h} = \SI{7200}{s}$.
\[ Q = I \times \Delta t = 10,0 \times 7200 = \SI{7,2e4}{C} \]
\[ n(e^-) = \frac{Q}{\mathcal{F}} = \frac{\num{7,2e4}}{96\,500} = \SI{0,746}{mol} \]
D'après $\ce{2Cl- -> Cl2 + 2e-}$ : $n(\ce{Cl2}) = \dfrac{n(e^-)}{2} = \SI{0,373}{mol}$.
\[ V(\ce{Cl2}) = n(\ce{Cl2}) \times V_m = 0,373 \times 24 \approx \SI{9,0}{L} \]
D'après $\ce{Cl2 + 2HO- -> Cl- + ClO- + H2O}$, 1 mol de \ce{Cl2} produit 1 mol de \ce{ClO-} :
\[ n(\ce{ClO-}) = n(\ce{Cl2}) \approx \SI{0,37}{mol} \]
\end{enumerate}
\end{corrige}

\begin{attention}[Barème indicatif et points de vigilance (type Probatoire, 10 pts)]
Schéma légendé correct (1,5 pt) ; comparaison des potentiels et conclusion à la cathode (1,5 pt) ; rôle de la surtension à l'anode (1,5 pt) ; demi-équations et bilan équilibrés (2 pts) ; nom de l'ion \ce{ClO-} (0,5 pt) ; calculs de $Q$, $n(e^-)$, $V(\ce{Cl2})$ et $n(\ce{ClO-})$ avec unités (3 pts).
\begin{itemize}
\item Ne jamais écrire $\ce{Na+ + e- -> Na}$ en solution aqueuse : c'est toujours l'eau qui est réduite.
\item Convertir les heures en secondes ; conserver les unités à chaque étape.
\item Le bilan ne doit pas faire apparaître d'électrons.
\end{itemize}
\end{attention}

\begin{corrige}[Exercice 10 --- Étude de la courbe caractéristique d'un électrolyseur]
\begin{enumerate}
\item \textbf{Courbe $I = f(U)$ :}
\begin{popfigure}
\begin{tikzpicture}
\begin{axis}[
width=0.9\linewidth, height=7.5cm,
xlabel={$U$ (V)}, ylabel={$I$ (mA)},
xmin=0, xmax=4, ymin=0, ymax=220,
xtick={0,0.5,1,1.5,2,2.5,3,3.5,4},
ytick={0,20,40,60,80,100,120,140,160,180,200,220},
grid=major, grid style={popmuted!40},
axis lines=left,
tick label style={font=\small},
label style={font=\small},
]
\addplot[popcurve, very thick, mark=*, mark size=1.8] coordinates {
(0.5,0) (1.0,0) (1.5,5) (2.0,40) (2.5,90) (3.0,140) (3.5,190)
};
\addplot[poporange, dashed, thick, domain=1.2:4] {100*x - 150};
\draw[popink, dashed, thick] (axis cs:1.5,0) -- (axis cs:1.5,220);
\node[popink, font=\small, anchor=west] at (axis cs:1.55,200) {$U_s \approx \SI{1,5}{V}$};
\end{axis}
\end{tikzpicture}
\legende{Courbe caractéristique $I = f(U)$. Le prolongement de la partie linéaire coupe l'axe des tensions en $U_s$.}
\end{popfigure}

\item \textbf{Tension de seuil :} en prolongeant la partie linéaire de la courbe jusqu'à $I = 0$, on lit :
\[ U_s \approx \SI{1,5}{V} \]

\item \textbf{Tension théorique minimale :}
\[ \Delta E^\circ = E^\circ(\ce{O2}/\ce{H2O}) - E^\circ(\ce{H+}/\ce{H2}) = 1,23 - 0,00 = \SI{1,23}{V} \]
On constate que $U_s \approx \SI{1,5}{V} > \SI{1,23}{V}$. L'écart, environ \SI{0,3}{V}, s'explique par le phénomène de \cle{surtension} : la réaction de dégagement des gaz (surtout \ce{O2} sur le platine) est cinétiquement lente, il faut appliquer une tension supérieure à la tension thermodynamique pour que l'électrolyse démarre à vitesse appréciable. S'ajoute éventuellement la chute ohmique dans l'électrolyte.

\item \textbf{Demi-équations et bilan} (électrodes de platine inertes, solution d'acide sulfurique dilué) :
\begin{itemize}
\item Cathode (réduction) : $\ce{2H+ + 2e- -> H2}$
\item Anode (oxydation de l'eau, car \ce{SO4^2-} ne s'oxyde pas) : $\ce{2H2O -> O2 + 4H+ + 4e-}$
\end{itemize}
Bilan (on multiplie la première par 2) :
\[ \ce{2H2O -> 2H2 + O2} \]
C'est en réalité l'\textbf{électrolyse de l'eau} ; l'acide sulfurique ne sert qu'à rendre la solution conductrice.

\item \textbf{Volumes de gaz à $U = \SI{3,5}{V}$ :} le tableau donne $I = \SI{190}{mA} = \SI{0,190}{A}$, et $\Delta t = \SI{30}{min} = \SI{1800}{s}$.
\[ Q = I \times \Delta t = 0,190 \times 1800 = \SI{342}{C} \]
\[ n(e^-) = \frac{Q}{\mathcal{F}} = \frac{342}{96\,500} = \SI{3,54e-3}{mol} \]
\begin{itemize}
\item Dihydrogène : $n(\ce{H2}) = \dfrac{n(e^-)}{2} = \SI{1,77e-3}{mol}$, d'où
\[ V(\ce{H2}) = 1,77 \times 10^{-3} \times 24 = \SI{4,3e-2}{L} \approx \SI{43}{mL} \]
\item Dioxygène : $n(\ce{O2}) = \dfrac{n(e^-)}{4} = \SI{8,9e-4}{mol}$, d'où
\[ V(\ce{O2}) = 8,9 \times 10^{-4} \times 24 = \SI{2,1e-2}{L} \approx \SI{21}{mL} \]
\end{itemize}
On vérifie le rapport stoechiométrique : $V(\ce{H2}) = 2\,V(\ce{O2})$, caractéristique de l'électrolyse de l'eau.
\end{enumerate}
\end{corrige}

\begin{attention}[Barème indicatif et points de vigilance (type Probatoire, 10 pts)]
Courbe soignée avec échelles et axes légendés (2 pts) ; détermination graphique de $U_s$ (1 pt) ; calcul de $\Delta E^\circ$ et interprétation par la surtension (2 pts) ; demi-équations et bilan (2 pts) ; calculs de $Q$, $n(e^-)$ et des deux volumes avec unités (3 pts).
\begin{itemize}
\item Attention à la conversion $\SI{190}{mA} = \SI{0,190}{A}$ : oublier ce facteur 1000 fausse tout le calcul.
\item Le volume de \ce{H2} doit être le \textbf{double} de celui de \ce{O2} : si ce n'est pas le cas, reprendre la stoechiométrie électronique.
\end{itemize}
\end{attention}

\begin{corrige}[Exercice 11 --- Protection d'une canalisation enterrée et galvanisation au zinc]
\textbf{Partie A.}
\begin{enumerate}
\item \textbf{Explication :} $E^\circ(\ce{Mg^2+}/\ce{Mg}) = \SI{-2,37}{V} < E^\circ(\ce{Fe^2+}/\ce{Fe}) = \SI{-0,44}{V}$ : le magnésium est un réducteur \textbf{plus fort} que le fer, il s'oxyde donc préférentiellement. Le couple fer/magnésium relié par le fil constitue une pile dans laquelle le magnésium est l'anode (il se consomme) et le fer la cathode (il ne peut plus être oxydé, car il reçoit les électrons). C'est la \cle{protection cathodique par anode sacrificielle}.

\item \textbf{Demi-équations :}
\begin{itemize}
\item Sur le bloc de magnésium (oxydation) : $\ce{Mg -> Mg^2+ + 2e-}$
\item À la surface de la canalisation (réduction, sol humide aéré) : $\ce{O2 + 2H2O + 4e- -> 4HO-}$ ; en milieu peu aéré : $\ce{2H2O + 2e- -> H2 + 2HO-}$.
\end{itemize}

\item \textbf{Durée de vie du bloc :} $m = \SI{2,0}{kg} = \SI{2,0e3}{g}$.
\[ n(\ce{Mg}) = \frac{m}{M(\ce{Mg})} = \frac{\num{2,0e3}}{24,3} = \SI{82,3}{mol} \]
D'après $\ce{Mg -> Mg^2+ + 2e-}$ : $n(e^-) = 2 \times 82,3 = \SI{164,6}{mol}$.
\[ Q = n(e^-) \times \mathcal{F} = 164,6 \times 96\,500 = \SI{1,59e7}{C} \]
Avec $I = \SI{50}{mA} = \SI{0,050}{A}$ :
\[ \Delta t = \frac{Q}{I} = \frac{\num{1,59e7}}{0,050} = \SI{3,18e8}{s} \]
En années : $\Delta t = \dfrac{\num{3,18e8}}{\num{3,15e7}} \approx \boxed{\SI{10}{ans}}$.
\end{enumerate}

\textbf{Partie B.}
\begin{enumerate}
\setcounter{enumi}{3}
\item \textbf{Rôle de la surtension :} thermodynamiquement, les ions \ce{H+} devraient être réduits avant les ions \ce{Zn^2+} car $E^\circ(\ce{H+}/\ce{H2}) = \SI{0,00}{V} > E^\circ(\ce{Zn^2+}/\ce{Zn}) = \SI{-0,76}{V}$. Cependant, la réduction des ions \ce{H+} en \ce{H2} sur une surface de zinc présente une \cle{surtension} élevée : le potentiel réel auquel \ce{H2} se dégage devient plus négatif que celui du dépôt du zinc. La réduction $\ce{Zn^2+ + 2e- -> Zn}$, quasi dépourvue de surtension, devient alors la réaction observée : le zinc peut être déposé en solution aqueuse.

\item \textbf{Durée de l'électrolyse :}
\[ n(\ce{Zn}) = \frac{m}{M(\ce{Zn})} = \frac{6,54}{65,4} = \SI{0,100}{mol} \]
D'après $\ce{Zn^2+ + 2e- -> Zn}$ : $n(e^-) = 2 \times 0,100 = \SI{0,200}{mol}$.
\[ Q = n(e^-) \times \mathcal{F} = 0,200 \times 96\,500 = \SI{1,93e4}{C} \]
\[ \Delta t = \frac{Q}{I} = \frac{\num{1,93e4}}{4,0} = \SI{4,83e3}{s} \]
Conversion : $4825\,\text{s} = 3600\,\text{s} + 1225\,\text{s} = 1\,\text{h} + 20\,\text{min} + 25\,\text{s}$, soit environ \boxed{1\,\text{h}\,20\,\text{min}}.
\end{enumerate}
\end{corrige}

\begin{attention}[Barème indicatif et points de vigilance (type Probatoire, 12 pts)]
Partie A : comparaison des potentiels et nom de la protection (2 pts) ; demi-équations (2 pts) ; calcul de la durée de vie avec conversions (3 pts). Partie B : explication par la surtension (2 pts) ; calcul de la durée (3 pts).
\begin{itemize}
\item Convertir les masses en grammes, les courants en ampères, et utiliser la cohérence des unités ($\SI{50}{mA} = \SI{0,050}{A}$).
\item Pour l'anode sacrificielle, le métal protecteur doit avoir un potentiel standard \textbf{inférieur} à celui du fer : justifier toujours ce choix par une inégalité de potentiels.
\item Ne pas confondre protection cathodique (anode sacrificielle, sans générateur) et galvanisation (électrolyse, avec générateur).
\end{itemize}
\end{attention}

\newpage

\coursec{Fiche bilan}

\begin{popfigure}
\begin{tikzpicture}[scale=0.85, transform shape,
  centre/.style={circle, draw=popdark, fill=popdark, text=white, font=\bfseries\small, minimum size=2.8cm, align=center},
  branche/.style={rounded corners=4pt, draw=#1, fill=#1!15, text=popdark, font=\footnotesize, align=center, inner sep=4pt},
  lien/.style={-{Stealth[length=2.5mm]}, thick, #1}]
\node[centre] (C) {Électrolyse\\transformations\\forcées};
\node[branche=popblue, above left=0.6cm and 2.3cm of C] (B1) {\textbf{Généralités}\\Réaction imposée\\par le générateur\\$\Delta_rG>0$};
\node[branche=poporange, above=0.7cm of C] (B2) {\textbf{Réactions aux électrodes}\\Anode ($+$) : oxydation\\Cathode ($-$) : réduction};
\node[branche=popgreen, above right=0.6cm and 2.3cm of C] (B3) {\textbf{Prévision}\\Cathode : oxydant le\\plus fort ($E^\circ$ max)\\Anode : réducteur le\\plus fort ($E^\circ$ min)};
\node[branche=poppurple, below left=0.6cm and 2.3cm of C] (B4) {\textbf{Surtension $\eta$}\\Écart potentiel réel/théo\\$\Rightarrow$ \ce{H2} au lieu de \ce{Na},\\ \ce{O2} au lieu de \ce{Cl2} (parfois)};
\node[branche=popgold, below=0.7cm of C] (B5) {\textbf{Aspects quantitatifs}\\$Q=I\,t=n_e\,F$\\$F=\SI{96500}{C.mol^{-1}}$\\$m=\dfrac{M\,I\,t}{z\,F}$};
\node[branche=poppink, below right=0.6cm and 2.3cm of C] (B6) {\textbf{Corrosion \& protection}\\Rouille : \ce{Fe -> Fe^2+}\\Protection : revêtement,\\anode sacrificielle (\ce{Zn}, \ce{Mg})};
\draw[lien=popblue] (C) -- (B1);
\draw[lien=poporange] (C) -- (B2);
\draw[lien=popgreen] (C) -- (B3);
\draw[lien=poppurple] (C) -- (B4);
\draw[lien=popgold] (C) -- (B5);
\draw[lien=poppink] (C) -- (B6);
\end{tikzpicture}
\legende{Carte mentale de la leçon 09 : l'électrolyse, du principe aux applications industrielles.}
\end{popfigure}

\begin{definition}[Bilan de la leçon -- Électrolyse]
\textbf{Définitions clés}
\begin{itemize}
  \item \cle{Électrolyse} : transformation chimique \textbf{forcée} (non spontanée, $\Delta_rG > 0$) provoquée par le passage d'un courant continu imposé par un générateur externe. Elle se produit dans le sens \emph{inverse} de la réaction spontanée de la pile correspondante.
  \item \cle{Anode} : électrode reliée au pôle $+$ du générateur ; siège d'une \textbf{oxydation} (perte d'électrons). \cle{Cathode} : électrode reliée au pôle $-$ ; siège d'une \textbf{réduction} (gain d'électrons).
  \item \cle{Surtension ($\eta$)} : potentiel supplémentaire (en valeur absolue) à appliquer par rapport au potentiel d'équilibre thermodynamique ($E^\circ$) pour qu'une réaction cinétiquement lente se produise effectivement à une électrode donnée (ex. dégagement de \ce{H2} sur cathode Hg ou Pt au lieu du dépôt de \ce{Na}).
  \item \cle{Corrosion} : oxydation spontanée d'un métal par le dioxygène de l'air en milieu aqueux (ex. rouille : \ce{Fe2O3.nH2O} sur le fer).
\end{itemize}

\textbf{Demi-équations de référence (à connaître par cœur)}
\begin{center}
\begin{tabularx}{\linewidth}{|l|X|X|}
\hline
\rowcolor{popblueL} \textbf{Milieu / Espèce} & \textbf{Cathode (Réduction)} & \textbf{Anode (Oxydation)} \\
\hline
Eau (milieu neutre/basique) & $\ce{2H2O + 2e- -> H2 + 2HO-}$ & $\ce{4HO- -> O2 + 2H2O + 4e-}$ \\
\hline
Eau (milieu acide) & $\ce{2H+ + 2e- -> H2}$ & $\ce{2H2O -> O2 + 4H+ + 4e-}$ \\
\hline
Ion cuivre & $\ce{Cu^2+ + 2e- -> Cu}$ & $\ce{Cu -> Cu^2+ + 2e-}$ \\
\hline
Ion chlorure & -- & $\ce{2Cl- -> Cl2 + 2e-}$ \\
\hline
Ion sodium & $\ce{Na+ + e- -> Na}$ (théorique) & -- \\
\hline
\end{tabularx}
\end{center}

\textbf{Bilans des électrolyses types au programme (électrodes inertes Pt, sauf indication)}
\begin{center}
\begin{tabularx}{\linewidth}{|l|X|X|X|}
\hline
\rowcolor{popblueL} \textbf{Électrolyte} & \textbf{Cathode ($-$)} & \textbf{Anode ($+$)} & \textbf{Bilan global} \\
\hline
\ce{HCl} aqueux & $\ce{2H+ + 2e- -> H2}$ & $\ce{2Cl- -> Cl2 + 2e-}$ & $\ce{2HCl -> H2 + Cl2}$ \\
\hline
\ce{NaCl} aqueux (saumure) & $\ce{2H2O + 2e- -> H2 + 2HO-}$ & $\ce{2Cl- -> Cl2 + 2e-}$ & $\ce{2NaCl + 2H2O -> Cl2 + H2 + 2NaOH}$ \\
\hline
\ce{NaOH} aqueux & $\ce{2H2O + 2e- -> H2 + 2HO-}$ & $\ce{4HO- -> O2 + 2H2O + 4e-}$ & $\ce{2H2O -> 2H2 + O2}$ \\
\hline
\ce{CuSO4} aqueux & $\ce{Cu^2+ + 2e- -> Cu}$ (dépôt) & $\ce{2H2O -> O2 + 4H+ + 4e-}$ & $\ce{2CuSO4 + 2H2O -> 2Cu + O2 + 2H2SO4}$ \\
\hline
\ce{CuSO4} aqueux (électrodes \ce{Cu}) & $\ce{Cu^2+ + 2e- -> Cu}$ & $\ce{Cu -> Cu^2+ + 2e-}$ (dissolution) & Transfert de \ce{Cu} anode $\to$ cathode \\
\hline
\end{tabularx}
\end{center}

\textbf{Formules et grandeurs quantitatives}
\begin{center}
\begin{tabularx}{\linewidth}{|l|X|l|}
\hline
\rowcolor{popblueL} \textbf{Relation} & \textbf{Signification} & \textbf{Unités} \\
\hline
$Q = I\,t$ & Quantité d'électricité traversant l'électrolyseur & $Q$ en \si{C}, $I$ en \si{A}, $t$ en \si{s} \\
\hline
$Q = n_e\,F$ & Lien avec la quantité d'électrons échangés & $F = \SI{96500}{C.mol^{-1}}$ (constante de Faraday) \\
\hline
$n = \dfrac{I\,t}{z\,F}$ & Quantité de matière de produit formé ($z$ = nb d'électrons par molécule) & $n$ en \si{mol} \\
\hline
$m = \dfrac{M\,I\,t}{z\,F}$ & Masse de métal déposé (ou dissous) & $m$ en \si{g}, $M$ en \si{g.mol^{-1}} \\
\hline
$U_{\text{min}} = E_{\text{anode}} - E_{\text{cathode}} + \eta_a + |\eta_c|$ & Tension minimale à imposer (potentiels de réduction $E$) & \si{V} \\
\hline
\end{tabularx}
\end{center}

\textbf{Méthodes express}
\begin{itemize}
  \item \textbf{Prévoir les produits d'une électrolyse aqueuse} :
  \begin{enumerate}
    \item Lister \textbf{toutes} les espèces présentes : ions du soluté, \ce{H2O}, électrodes (inertes ou actives).
    \item \textbf{Cathode} : l'espèce qui se réduit est l'\textbf{oxydant le plus fort} (potentiel de réduction $E^\circ$ le plus \textbf{élevé} / moins négatif), \emph{en tenant compte des surtensions} (ex. $\eta_{\ce{H2}}$ élevée sur Hg $\Rightarrow$ dépôt \ce{Na} possible en théorie, mais sur Pt $\eta_{\ce{H2}}$ faible $\Rightarrow$ \ce{H2} dégagé).
    \item \textbf{Anode} : l'espèce qui s'oxyde est le \textbf{réducteur le plus fort} (potentiel de réduction $E^\circ$ le plus \textbf{bas} / plus négatif), \emph{en tenant compte des surtensions} (ex. $\eta_{\ce{O2}}$ élevée sur Pt $\Rightarrow$ \ce{Cl-} oxydé avant \ce{H2O} si $[\ce{Cl-}]$ élevé).
    \item Écrire les deux demi-équations équilibrées (atomes + charges), les combiner pour obtenir le bilan global.
  \end{enumerate}
  \item \textbf{Calcul quantitatif (masse, volume, débit)} :
  \begin{enumerate}
    \item Convertir le temps $t$ en \textbf{secondes}.
    \item Calculer $Q = I \times t$ (Coulombs).
    \item Calculer $n_e = Q / F$ (mol d'électrons).
    \item Utiliser un \textbf{tableau d'avancement} reliant $n_e$ aux quantités de produits via les coefficients stœchiométriques des demi-équations.
    \item En déduire la masse $m = n \times M$ ou le volume $V = n \times V_m$ (conditions T, P données).
  \end{enumerate}
  \item \textbf{Protéger un métal de la corrosion} :
  \begin{enumerate}
    \item \textbf{Revêtement barrière} : peinture, vernis, graisse, galvanisation à chaud (zinc), étamage (étain) -- isole le métal de l'air et de l'eau.
    \item \textbf{Protection cathodique (anode sacrificielle)} : fixer un bloc de métal \textbf{plus réducteur} (potentiel $E^\circ$ plus bas) que le fer : \ce{Zn} ($E^\circ = -\SI{0,76}{V}$) ou \ce{Mg} ($E^\circ = -\SI{2,37}{V}$). Ce bloc s'oxyde à la place du fer : $\ce{Zn -> Zn^2+ + 2e-}$ ; $\ce{Fe}$ reste intact (cathode).
  \end{enumerate}
\end{itemize}
\end{definition}

\begin{aretenir}[Checklist avant le Probatoire]
\begin{itemize}
  \item[$\square$] Je sais définir une électrolyse et expliquer pourquoi $\Delta_rG > 0$ (réaction forcée).
  \item[$\square$] Je ne confonds plus : \textbf{Anode} = $+$ = \textbf{Oxydation} ; \textbf{Cathode} = $-$ = \textbf{Réduction} (mnémo : \textbf{An}Ox / \textbf{Ca}Ré).
  \item[$\square$] J'écris correctement les 4 demi-équations de l'eau (acide et neutre/basique) \textbf{équilibrées en atomes et en charges}.
  \item[$\square$] Je sais prévoir les produits des 4 électrolyses types au programme : \ce{HCl}, \ce{NaCl} (saumure $\to$ eau de Javel), \ce{NaOH}, \ce{CuSO4} (électrodes Pt et \ce{Cu}).
  \item[$\square$] J'explique le rôle de la \textbf{surtension} : pourquoi \ce{H2} se dégage à la cathode (Pt) au lieu de \ce{Na} lors de l'électrolyse de \ce{NaCl} aqueux, et pourquoi \ce{Cl2} se dégage à l'anode au lieu de \ce{O2} (si $[\ce{Cl-}]$ suffisant).
  \item[$\square$] Je connais la valeur de la constante de Faraday : $F = \SI{96500}{C.mol^{-1}}$.
  \item[$\square$] J'applique rigoureusement $Q = I\,t = n_e\,F$ avec $t$ en \textbf{secondes}.
  \item[$\square$] Je sais calculer la masse de métal déposée : $m = \dfrac{M\,I\,t}{z\,F}$ (identifier $z$ depuis la demi-équation).
  \item[$\square$] Je cite deux applications industrielles majeures : \textbf{synthèse de l'eau de Javel} (électrolyse saumure) et \textbf{obtention de l'aluminium} (électrolyse alumine fondue cryolithe -- savoir culturel).
  \item[$\square$] J'explique la formation de la \textbf{rouille} (pile de corrosion : anode \ce{Fe/Fe^2+}, cathode \ce{O2/H2O}) et je donne \textbf{deux} méthodes de protection (revêtement + anode sacrificielle).
  \item[$\square$] Je vérifie \textbf{systématiquement} l'équilibre des atomes et des charges de chaque demi-équation rédigée.
\end{itemize}
\end{aretenir}

\findecours

\end{document}
